\documentclass[11pt]{article} % required
\usepackage{geometry}
\geometry{letterpaper, portrait, margin=1in, headheight=110pt}
\usepackage{fancyhdr}
\pagestyle{fancy}

\usepackage{amssymb}
\usepackage{amsmath}
\usepackage{amsthm}

\usepackage{color}
%\usepackage[T1]{fontenc}
\usepackage{droid}
\renewcommand*\familydefault{\sfdefault} %% Only if the base font of the document is to be sans serif

\usepackage{csquotes}
\usepackage{varwidth}
\usepackage{capt-of}
\usepackage{booktabs}
\usepackage[pdftex]{graphicx}
% declare the path(s) where your graphic files are
\graphicspath{{img/pdf/}{img/png/}{img/pdf_tex/}}
% The hyperref package gives us a pdf with properly built
% internal navigation ('pdf bookmarks' for the table of contents,
% internal cross-reference links, web links for URLs, etc.)
\usepackage[]{hyperref}
\usepackage{xcolor}
\hypersetup{
	colorlinks=true,% hyperlinks will be coloured
	linkcolor=black,% hyperlink text will be green
	linkbordercolor=red,% hyperlink border will be red
    citecolor=blue,
    urlcolor=blue,
	urlbordercolor={blue}
}
\makeatletter
\Hy@AtBeginDocument{%
  \def\@pdfborder{0 0 1}% Overrides border definition set with colorlinks=true
  \def\@pdfborderstyle{/S/U/W 1}% Overrides border style set with colorlinks=true
                                % Hyperlink border style will be underline of width 1pt
}
\makeatother

\usepackage{algorithm}
\usepackage{algpseudocode}
\usepackage{textgreek}
\usepackage{mathtools} % for \ceil and \floor

\usepackage{listings}

\lstdefinelanguage{uJS} {
	language = C,
	morekeywords = {var, const, function, true, false, try, throw, catch},
}

\lstdefinelanguage{C99} {
	language = C,
	morekeywords = {int8_t, uint8_t, int16_t, uint16_t, int32_t, uint32_t},
}

%\usepackage{inconsolata}

\providecommand{\tightlist}{\setlength{\itemsep}{0pt}\setlength{\parskip}{0pt}}

\lstset{
	basicstyle=\footnotesize\ttfamily, 
	keywordstyle=\color{blue}\bfseries\ttfamily,
	stringstyle=\color{red}\ttfamily,
	commentstyle=\color{magenta}\ttfamily,
	identifierstyle=\color{black}\ttfamily,          
%		 morecomment=[l][\color{magenta}]{\#},
	numberstyle=\tiny,          
	numbersep=5pt,             
	tabsize=4,                
	extendedchars=true,      
	breaklines=true,        
	showspaces=false,      
	showtabs=false,       
	xleftmargin=17pt,
	framexleftmargin=17pt,
	framexrightmargin=5pt,
	framexbottommargin=4pt,
	showstringspaces=false
}

\lstloadlanguages{
	C, C99, uJS 
}

% Definition of \maketitle
\makeatletter         
\def\@maketitle{
	\raggedright
	\vfill
	\begin{center}
	{\Huge \bfseries \sffamily \@title }\\[4ex] 
	{\Large \sffamily \@author}\\ 
	{\sffamily \textit{\email}}\\[8ex]
	{\sffamily \@date}\\[8ex]
	{\sffamily \bfseries Revision \therevision}
	\vfill
	\includegraphics[height=60pt]{img/pdf/thinkos_embedded.pdf}
	\vfill
	\end{center}
	\noindent\rule{\linewidth}{1pt}\\
	\quad
}
\makeatother

\begingroup
\catcode`_=13
\gdef\var{%
	\begingroup
	\catcode`_=13
	\def_{\_}%
	\dovar
}
\gdef\dovar#1{%
	\mathtt{#1}%
\endgroup
}
\endgroup

\fancyhead[R]{\thepage}
\rhead{\includegraphics[height=20pt]{img/pdf/thinkos_logo.pdf}}
\lhead{\thetitle}
\rfoot{Page \thepage}
\cfoot{}

\usepackage{lettrine}

\newcommand{\dg}{$^\circ$ }

\newcommand{\ThinkOS}{\textbf{ThinkOS} }

\usepackage[nounderscore]{syntax}

\usepackage{vhistory}

%\lettrine[findent=2pt]{\textbf{T}}{his}



% Define Title Author Date and Revision
\makeatletter         
\title{ThinkOS - User Guide}
\author{Bob Mittmann}
\date{January 2017}
\def\email{bobmittmann@gmail.com} % revision number variable
\def\therevision{0.1} % revision number variable
\let\thetitle\@title % title variable
\makeatother

\begin{document}

\maketitle
\thispagestyle{empty}
\newpage
\pagenumbering{Roman}
\tableofcontents
\newpage
\listoffigures
\newpage
\listoftables
\newpage

% Start of the revision history table
\begin{versionhistory}
  \vhEntry{0.1}{21.01.17}{BM}{created}
\end{versionhistory}

\newpage
\pagenumbering{arabic}

\section{Introduction}

\lettrine{T}{his} user guide provides an overview of the \ThinkOS real-time operating system. The target audience for this guide are embedded software developers knowledgeable in C language. Notice that no mention of C++ is made in this text.\\

\ThinkOS was designed specifically for the \textbf{Cortex-M} family of microprocessors from \textbf{ARM}. These are all 32bits CPUs with mostly 16bits instruction set, which produces very compact code, hence suitable for deep embedded systems with stringent memory and other peripheral resources constraints.

At this moment \ThinkOS supports \textbf{M3}, \textbf{M4} and \textbf{M7} cores.

The popular \textbf{M0} family lacks some features needed for some operations in the kernel. Porting to these platforms is feasible but needs someone with time to do that. Any volunteer?

ThinkOS is a Real Time Operating System designed for the ARM Cortex-M core. It takes advantage of some unique features of this processor to enable for very low task switch latency.


\subsection{Why Another OS?}

The real answer for this question perhaps is simply: ``Just because I can''. But as most people won't be satisfied with this shallow reply here are some justifications:

\subsection{Where do I can get ThinkOS from?}

\begin{flushleft}
Repository: [\href{https://github.com/bobmittmann/thinkos}{https://github.com/bobmittmann/thinkos}]\\
Git HTTPS: [\href{https://github.com/bobmittmann/thinkos.git}{https://github.com/bobmittmann/thinkos.git}]
\end{flushleft}

\subsection{Performance}

The main bottleneck for performance is usually the scheduler. The ThinkOS running in a 120MHz Cortex-M3 (STM32F270) have a measured latency time of 0.5 microseconds.

There are some tests and calculations that need to be done to determine the latency of other elements of the system...

\subsubsection{Coding Style}

\ThinkOS source uses the \href{https://www.kernel.org/doc/html/v4.10/process/coding-style.html}{Linux kernel coding style} and I encourage all developers to follow it. This is just little more than a matter of preference, please read the style guide to know the rational behind the rules before despising it ...

----------------------------------------------------------
== Introduction ==
ThinkOS is a Real Time Operating System designed for the ARM Cortex-M core. It takes advantage of some unique features of this processor to enable for very low task switch latency.
== Modes of usage ==
There are basically two ways of using the ThinkOS:
1. As a library linked against your application. (Need a term to decribe this like embedded mode...). 
In this type of usage there is usually no bootloader, the application boots and configure itself. In this scenario the firmware is the application. The advatnages of this mode are:
	- Small footprint...
	- Kernel calls and services are tailored for the application.
2. As part of the bootloader separated from the application. This mode will be called "standalone" mode, or "bootloader" mode.
In this case the application will be a separated entity (file, binary) and run in a separated memory space. The bootloader will be residing on the platform at all times and will be responsible to initialize the platform and to load the application. It's possible then to upload/upgrade/change the application without changing the bootlader and consequently the operating system. Advantages:
	- No need to configure the OS (need to elaborate...) to match the application. At least initially.
	- In field firmware upgrades are safer as the bootloader is the responsible for the upgrade.
	- The Debug Monitor features can be used for development.

One feature that distinguishes the \ThinkOS from most of other Operating Systems is the way the control structures for the several types of objects are held by kernel.

The usual approach for holding this information is to allocate a structure specific for a certain object and pass a pointer to the kernel. The kernel then stores this in a list or array. 
\ThinkOS in other hand holds the core data structure in a  Structure of Arrays (SOA) representation. The cost of this representation is the readability of the code, (it may look a little obscure at the beginning). 
The idea here is simplify the race condition avoidance in the kernel and speed-up the scheduler. Threads


\ThinkOS is a preemptive time-sharing kernel. Each thread has its own stack and a timer. This timer is used to block the thread's execution for a certain period of time as in \lstinline{thinkos_sleep()}.

\subsubsection{System Calls}

\subsection{Scheduler}
ThinkOS scheduler was designed to reduce the context switch time to a minimum. The most important piece on the Kernel is the Scheduler....

\subsection{Time Sharing}
All threads with priority higher than 0 go into a round-robin time sharing scheduling policy, when this feature is enabled. Their execution time will be inversely proportional to their priority.

\subsection{Interrupt Model}
There are two types of interrupt handlers with ThinkOS:

\paragraph{Implicit ISR}{Thread interrupt handlers that executes as normal threads. This is a simple call to \lstinline{thinkos_irq_wait(IRQ)} inside a regular task or function. The thread will block until this interrupt is raised. This mechanism provides a convenient and powerful tool to design drivers. Devices that benefit most of this type of ISR are simplex I/O channels like I2C or SPI, hardware accelerators like cryptographic ...}

\paragraph{Native ISR}{Normal interrupt handlers that shares a common stack. These handlers are subject to to some constraints regarding API calls. It is recommended to use this type of interrupt handler when low latency is a must. But bear in mind that if the handler is too complex the actual latency due to loading registers and housekeeping may be equivalent to a context switch in which case it will be equivalent of a threaded handler in terms of performance. Another case will be a high priority interrupt you want to make sure it will won't be preempted...}

\subsection{Configuration}



\section {Quick Start}

This guide will help you to install the necessary tools to compile and run a ThinkOS demo application on a development board.

\subsection{Windows Tools}

These are the basic steps to follow in order to install the tools: 

\begin{itemize}
\item Install GNU Toolchain for ARM Processros for Windows hosts.
\item Install GNU Make for Windows.
\item Install GNU CoreUtils for Windows.
\item Install the Java JRE from Oracle.
\item Install the Eclipse for C++ Developers IDE.
\item Install the GNU ARM Eclipse Plug-ins.
\item Install the C/C++ GDB Hardware Debugging Plug-in.
\item Install the STM32 ST-LINK Utility from STMicroelectronics.
\end{itemize}

Detailed instructions for each one of the steps can be found in the next sections of this document. 

\subsection{Demo Application}
There are some demo applications using \ThinkOS in the source tree. This is a guide on how to compile and run ThinkOS programs in the STM32F3Discovery board. 


\begin{itemize}
\item Get the \ThinkOS Source Code from GitHub.
\item Import the Eclipse projects for the target board into your workspace.
\item Compile the \ThinkOS BootLoader project using Eclipse IDE.
\item Load the \ThinkOS BootLoader on the board using the STM32 ST-LINK Utility.
\item Install the Virtual COM Driver on the PC to access the ThinkOS Debug-Monitor.
\item Compile the SDK Libraries (sdklib) for the target board using Eclipse.
\item Compile the demo application project using Eclipse.
\item Run the application with GDB Remote Debugging.
\end{itemize}



\subsection{ThinkOS Basics}

\lstset {language=C99}

In this tutorial the most basic \ThinkOS functions are described with some code examples. This is a summary of what it will be covered in this section:

\begin{itemize}
\item How to create and run threads.
\item Create, lock and unlock mutexes.
\item Create, signal and wait on semaphores.
\item Wait for time intervals (sleep).
\item Create a new thread
\end{itemize}

\subsection{Creating a Thread}

The easiest way to create and run a thread is to call \lstinline{thinkos_thread_create()}. Here is an example: 

\begin{lstlisting}[caption={Simple Thread Create Example},label={lst:example1},language=C99,frame=single]
static int my_task(void * arg)
{
    for (;;) {
        /* do something ... */
         thinkos_sleep(50);
    }
    return 0;
}

static uint32_t my_stack[256]; /* 1KB of stack is enough for most porposes */

void my_module_init(void)
{
    thinkos_thread_create(my_task, NULL, my_stack, sizeof(my_stack));
}
\end{lstlisting}

In this example we define a function to be the entry point of the thread, in this case \lstinline{my_task()}. We also define an array of 32bit words, amounting to 1KB, to be used as stack for the thread (\lstinline{my_stack}). We then call \lstinline{thinkos_thread_create()} which will allocate and run the thread. The return of this function is the thread id or a negative error code in case something goes wrong.

\subsubsection{Using semaphores}

ThinkOS semaphores allow processes and threads to synchronize their actions.

A semaphore is an integer whose value is never allowed to fall below zero. Two operations can be performed on semaphores: increment the semaphore value by one (\lstinline{thinkos_sem_post()}); and decrement the semaphore value by one (\lstinline{thinkos_sem_wait()}). If the value of a semaphore is currently zero, then a \lstinline{thinkos_sem_wait()} operation will block until the value becomes greater than zero. A semaphore should be allocated with \lstinline{thinkos_sem_alloc()} which will return a descriptor to be used in subsequent calls to  \lstinline{thinkos_sem_post()}, and \lstinline{thinkos_sem_wait()}.

In the example below, two treads are created: a producer and a consumer. The producer signals the semaphore periodically twice a second, whereas the consumer waits for the signal and prints a message when receives it.

\begin{lstlisting}[caption={Using semaphores example},label={lst:example2},language=C99,frame=single]
/* Global semaphore descriptor */

int my_sem;

/* Producer thread */
static int producer_task(void * arg)
{
    for (;;) {
        /* Wait 500 milliseconds... */
        thinkos_sleep(500);
        /* Signal the semaphore. */
        thinkos_sem_post(my_sem);
    }
    return 0;
}

/* Consumer thread */
static int consumer_task(void * arg)
{
    for (;;) {
        /* Wait on the semaphore... */
        thinkos_sem_wait(my_sem);
        /* Do something... */
        printf("Signal received\n");
    }
    return 0;
}


static uint32_t consumer_stack[256];
static uint32_t producer_stack[256];

void my_test_init(void)
{
    /* Allocate a semaphore with 0 initial count. */
    my_sem = thinkos_sem_alloc(0);
    /* Create a thread to signal the semaphore */
    thinkos_thread_create(producer_task, NULL, producer_stack, sizeof(producer_stack));
    /* Create a thread to wait on the semaphore */
    thinkos_thread_create(consumer_task, NULL, consumer_stack, sizeof(consumer_stack));
}
\end{lstlisting}

\subsubsection{Using MUTEXes}

A mutex is a MUTual EXclusion device, and is useful for protecting shared data structures from concurrent modifications, and implementing critical sections and monitors.

A mutex has two possible states: unlocked (not owned by any thread), and locked (owned by one thread). A mutex can never be owned by two different threads simultaneously. A thread attempting to lock a mutex that is already locked by another thread is suspended until the owning thread unlocks the mutex first.

A call to \lstinline{thinkos_mutex_lock()} locks the given mutex. If the mutex is currently unlocked, it becomes locked and owned by the calling thread, and \lstinline{thinkos_mutex_lock()} returns immediately. If the mutex is already locked by another thread, \lstinline{thinkos_mutex_lock()} suspends the calling thread until the mutex is unlocked by calling \lstinline{thinkos_mutex_unlock()}. A mutex should be allocated with \lstinline{thinkos_mutex_alloc()} which will return a descriptor to be used in subsequent mutex calls.

In the following example the serial port is a resource that needs to be protected to ensure the packets are transmitted in full.


\begin{lstlisting}[caption={Using MutExes examples},label={lst:example3},language=C99,frame=single]
/* Link Layer private data */
static struct {
	uint8_t addr; /* this panel address */
	
	
    struct serial_dev * serial; /* serial device driver */
    int mutex; /* protect the link layer send call */
} rs485;

void netlnk_send(unsigned int daddr, void * data, unsigned int len)
{
    struct lnkhdr hdr;

    /* lock the mutex, blocking the access to the serial port... */
    thinkos_mutex_lock(rs485.mutex);
    /* prepare header */
    hdr.sync = PKT_SYNC;
    hdr.daddr = daddr;
    hdr.saddr = rs485.addr;
    hdr.datalen = len;
    /* send header */
    serial_send(rs485.serial, &hdr, sizeof(struct lnkhdr));
    /* send payload */
    serial_send(rs485.serial, data, len);
    /* unlock mutex, other thread can use the serial port now. */
    thinkos_mutex_unlock(rs485.mutex);
}

void netlnk_init(unsigned int addr)
{
    /* Open the serial port */
    rs485.serial = stm32f_uart1_serial_init(9600, SERIAL_8N1);
    /* Set the local address */
    rs485.addr = addr;
    /* Allocate a mutex */
    rs485.mutex = thinkos_mutex_alloc();
}
\end{lstlisting}

\subsubsection{Time Wait}

Use the \lstinline{thinkos_sleep()} function to wait for a period of time. The calling thread suspends its execution until at least the time specified as parameter has elapsed. The time period is specified in milliseconds.
The example that follows a LED will blink repeatedly 4 times per second:

\begin{lstlisting}[caption={Time Wait Example},label={lst:example4},language=C99,frame=single]
int blink_task(void * arg)
{
    for (;;) {
        led_on();           /* Turn LED on */
        thinkos_sleep(100); /* Wait for 100 milliseconds */
        led_off();          /* Turn LED off */
        thinkos_sleep(150); /* Wait for 150 milliseconds */
    }

    return 0;
}
\end{lstlisting}

\subsubsection{Compiling the Code}

In order to compile the code in a Windows machine you must first install these tools:
\begin{itemize}
\item GNU Toolchain for ARM Processors for Windows hosts.
\item GNU Make for Windows.
\end{itemize}

\subsubsection{Compiling Using the Command Shell}

You can use the windows Command Prompt or a MinGW/MSYS terminal to compile the code.
Open a windows shell: Start->All Programs->Accessories->Command Prompt
Change to the directory  in the ThinkOS source tree.
Type: make

\section{Modules}
 Overview of the different modules.
 - Kernel
 - Debug Monitor
 - GDB server
 - Console

\subsection{Configuring}
 Description on how to configure ThinkOS.
 - bootloader mode
 - embedded mode

\subsection{Compiling }
 How and what is need in order to compile the ThinkOS.
 - Download the source code
 - Tools
 - Configure / or chose and example
 - Compile
 - Load and run

\section{API}

The ThinkOS API can be separated according to the type of object the operation is performed ...
Core - initialization configuration and debug
\subsection{Threads}
functions related to threads
\subsection{Mutexes}
\subsection{Conditional Variables}
\subsection{Semaphores}
\subsection{Events}
\subsection{Interrupts}
\subsection{Time}
\subsection{Kernel Objects}

 - system calls
 - library functions

The ThinkOS API can be separated according to the type of object the operation is performed ...

Core - initialization configuration and debug Threads - functions related to threads Mutexes Conditional Variables Semaphores Events Interrupts Time Kernel Objects

One feature that distinguish the ThinkOS from most of other Operating Systems is the way the control structures for the several type of objects are held by kernel. The usual approach for holding this information is to allocate a structure specific for a certain object and pass a pointer to the kernel. The kernel then stores this in a list or array. ThinkOS in other hand holds the information on a structure of arrays, each array refers to one field of an object. This may look very obscure but the idea here is to protect the memory and speed-up the scheduler.

Threads

The ThinkOS is a preemptive time-sharing kernel. Each thread have it's own stack and a timer. This timer is used to block the thread's execution for a certain period of time as in \lstinline{thinkos_sleep()}.

Scheduler

ThinkOS scheduler was designed to reduce the context switch time to a minimum. The most important piece on the Kernel is the Scheduler....

Time Sharing

All threads with priority higher than 0 go into a round-robin time sharing scheduling policy, when this feature is enabled. Their execution time will be inversely proportional to their priority.

Interrupt Model

There are two types of interrupt handlers with ThinkOS:

Implicit ISR - Thread interrupt handlers that executes as normal threads. This is a simple call to \lstinline{thinkos_irq_wait(IRQ)} inside a regular task or function. The thread will block until this interrupt is raised. This mechanism provides a convenient and powerful tool to design drivers. Devices that benefit most of this type of ISR are simplex I/O channels like I2C or SPI, hardware accelerators like cryptographic ...

Native ISR - Normal interrupt handlers that shares a common stack. These handlers are subject to to some constraints regarding API calls. It is recommended to use this type of interrupt handler when low latency is a must. But bare in mind that if the handler is too complex the actual latency due to loading registers and housekeeping may be equivalent to a context switch in which case it will be equivalent of a threaded handler in terms of performance. Another case will be a high priority interrupt you want to make sure it will won't be preempted...

Configuration

The file \lstinline{include/thinkos_sys.h} contains the definition of the structures used by the OS. In this file there is a series of macros used to set default values for the configuration options.

It is not recommended to change the values directly in \lstinline{thinkos_sys.h}, create a \lstinline{config.h} file instead, and define the macro \lstinline{CONFIG_H} at the compiler's: \lstinline{-DCONFIG_H}.

Note: if you are using the \ThinkOS as part of the YARD-ICE there is a \lstinline{config.h} file located at the source root.

\section{How the kernel works}
 - Data structures.
 - Design decisions.

\subsection{Exception Priorities}

\section {Demo Using Windows}

%{\includegraphics{images/image1.png}}{~ }

{ThinkOS demo on Windows}

\begin{itemize}
\item
  \hypertarget{h.gjdgxs}{}

  {For embedded developers on a budget.}
\item
  {All open source and free of}{~}{charge}{~}{tools.}
\end{itemize}

\section{\texorpdfstring{{Quick Guide}}{Quick Guide}}\label{quick-guide}

{This guide will help you to install the necessary tools to compile and
run a }{ThinkOS}{~demo application on a development board.}

\subsection{\texorpdfstring{{Windows
Tools}}{Windows Tools}}\label{windows-tools}

{These are the basic steps to follow in order to install the tools: }

\begin{enumerate}
\tightlist
\item
  {Install }{\protect\hyperlink{h.30j0zll}{GNU Toolchain for ARM
  Processros}}{\protect\hyperlink{h.30j0zll}{~}}{for Windows hosts.}
\item
  {Install }{\protect\hyperlink{h.1fob9te}{GNU
  Make}}{\protect\hyperlink{h.1fob9te}{~}}{for Windows.}
\item
  {Install }{\protect\hyperlink{h.3znysh7}{GNU
  CoreUtils}}{\protect\hyperlink{h.3znysh7}{~}}{for Windows.}
\item
  {Install the }{\protect\hyperlink{h.2et92p0}{Java
  JRE}}{\protect\hyperlink{h.2et92p0}{~}}{from Oracle.}
\item
  {Install the }{\protect\hyperlink{h.tyjcwt}{Eclipse for C++
  Developers}}{\protect\hyperlink{h.tyjcwt}{~}}{IDE.}
\item
  {Install the }{\protect\hyperlink{h.3dy6vkm}{GNU ARM Eclipse
  Plug-ins}}{\protect\hyperlink{h.3dy6vkm}{.}}
\item
  {Install the }{\protect\hyperlink{h.1t3h5sf}{C/C++ GDB Hardware
  Debugging Plug-in}}{\protect\hyperlink{h.1t3h5sf}{.}}
\item
  {Install the }{\protect\hyperlink{h.4d34og8}{STM32 ST-LINK
  Utility}}{\protect\hyperlink{h.4d34og8}{~}}{from STMicroelectronics.}
\end{enumerate}

{Detailed instructions for each one of the steps can be found in the
next sections of this document. }

\subsection{\texorpdfstring{{Demo
Application}}{Demo Application}}\label{demo-application}

{There are some demo applications using ThinkOS in the source tree. This
is a guide on how to compile and run }{ThinkOS}{~programs in the
}{\protect\hyperlink{h.35nkun2}{STM32F3Discovery}}{\protect\hyperlink{h.35nkun2}{~}}{board.
}

\begin{enumerate}
\tightlist
\item
  {Get the }{\protect\hyperlink{h.2s8eyo1}{ThinkOS Source
  Code}}{\protect\hyperlink{h.2s8eyo1}{~}}{from GitHub.}
\item
  {\protect\hyperlink{h.17dp8vu}{Import the Eclipse
  projects}}{\protect\hyperlink{h.17dp8vu}{~}}{for the target board into
  your workspace.}
\item
  {\protect\hyperlink{h.3rdcrjn}{Compile the ThikOS
  BootLoader}}{\protect\hyperlink{h.3rdcrjn}{~}}{project using Eclipse
  IDE.}
\item
  {Load the }{\protect\hyperlink{h.1ksv4uv}{ThinkOS
  BootLoader}}{\protect\hyperlink{h.1ksv4uv}{~}}{on the board using the
  }{STM32 ST-LINK Utility}{.}
\item
  {Install the }{\protect\hyperlink{h.44sinio}{Virtual COM
  Driver}}{\protect\hyperlink{h.44sinio}{~}}{on the PC to access the
  }{ThinkOS Debug-Monitor}{.}
\item
  {\protect\hyperlink{h.26in1rg}{Compile the SDK
  Libraries}}{\protect\hyperlink{h.26in1rg}{~}}{(sdklib) for the target
  board using Eclipse.}
\item
  {\protect\hyperlink{h.lnxbz9}{Compile the demo
  application}}{\protect\hyperlink{h.lnxbz9}{~}}{project using Eclipse.}
\item
  {Run the application with }{\protect\hyperlink{h.2jxsxqh}{GDB Remote
  Debugging}}{\protect\hyperlink{h.2jxsxqh}{.}}
\end{enumerate}

\hypertarget{h.30j0zll}{\section{\texorpdfstring{{GNU Toolchain for ARM
Processors}}{GNU Toolchain for ARM Processors}}\label{h.30j0zll}}

{The GNU toolchain is a broad collection of programming tools produced
by the GNU Project. These tools form a toolchain (a suite of tools used
in a serial manner) used for developing software applications and
operating systems.}

\begin{itemize}
\tightlist
\item
  {Download and install the }{GNU Toolchain for ARM Processros}{~from
  the site:
  }{\href{https://www.google.com/url?q=https://launchpad.net/gcc-arm-embedded\&sa=D\&ust=1511275046369000\&usg=AFQjCNGxWGT_YonUxM4BFF9geho0LfeKww}{https://launchpad.net/gcc-arm-embedded}}{.
  At the time of writing this document the latest version was
  4.9-2015-q2, which can be downloaded directly from:
  }{\href{https://www.google.com/url?q=https://launchpad.net/gcc-arm-embedded/4.9/4.9-2015-q2-update/\%2Bdownload/gcc-arm-none-eabi-4_9-2015q2-20150609-win32.exe\&sa=D\&ust=1511275046369000\&usg=AFQjCNHgw7yb_djkKsPe-rUggiSx6ZpV-Q}{https://launchpad.net/gcc-arm-embedded/4.9/4.9-2015-q2-update/+download/gcc-arm-none-eabi-4\_9-2015q2-20150609-win32.exe}}
\item
  {Run the downloaded file
  (gcc-arm-none-eabi-4\_9-2015q2-20150609-win32.exe) to start the
  installation.}
\end{itemize}

{\includegraphics{images/image3.png}}

\begin{itemize}
\tightlist
\item
  {In the next steps use the default values suggested by the wizard
  installer.}
\item
  {In the last screen, after the files were copied, select the ``Add
  path to environment variable'' option. }
\end{itemize}

{\includegraphics{images/image5.png}}

\hypertarget{h.1fob9te}{\section{\texorpdfstring{{GNU Make for
Windows}}{GNU Make for Windows}}\label{h.1fob9te}}

{Make is a tool which controls the generation of executables and other
non-source files of a program from the program's source files. Make gets
its knowledge of how to build your program from a file called the
makefile, which lists each of the non-source files and how to compute it
from other files. When you write a program, you should write a makefile
for it, so that it is possible to use Make to build and install the
program.}

\begin{itemize}
\tightlist
\item
  {Download the GNU Make package:
  http://sourceforge.net/projects/gnuwin32/files/make/3.81/make-3.81.exe/download}
\item
  {Run the downloaded program ``make-3.81.exe'' and follow the
  installation steps. }
\end{itemize}

{\includegraphics{images/image4.png}}

\begin{itemize}
\tightlist
\item
  {Install on the same folder as the GCC Toolchain ~so you won't need to
  modify your path. During the installation in the ``Select Destination
  Location:'' page select: ``C:\textbackslash{}Program Files
  (x86)\textbackslash{}GNU Tools ARM Embedded\textbackslash{}4.9
  2015q2'' in the text box.}
\end{itemize}

{\includegraphics{images/image7.png}}

\begin{itemize}
\tightlist
\item
  {Also there is no point of creating a ``Start Menu'' folder. Disable
  this option by checking the ``Don't create a Start Menu folder'' in
  the page: ``Select Start Menu Folder''. }
\end{itemize}

\hypertarget{h.3znysh7}{\section{\texorpdfstring{{GNU CoreUtils for
Windows}}{GNU CoreUtils for Windows}}\label{h.3znysh7}}

{The GNU Core Utilities are the basic file, shell and text manipulation
utilities of the GNU operating system. These are the core utilities
which are expected to exist on every operating system.}

\begin{itemize}
\tightlist
\item
  {Download the GnuCore Utils package from:
  }{\href{https://www.google.com/url?q=http://sourceforge.net/projects/gnuwin32/files/coreutils/5.3.0/coreutils-5.3.0.exe/download\&sa=D\&ust=1511275046372000\&usg=AFQjCNEiYiOPmJnrJ9jTIPo2YGpR3Yz9wg}{http://sourceforge.net/projects/gnuwin32/files/coreutils/5.3.0/coreutils-5.3.0.exe/download}}
\item
  {Run the downloaded program ``coreutils-5.3.0.exe'' and follow the
  installation steps.}
\item
  {Please refer to the }{GNU Make for Windows}{~installation procedure.}
\end{itemize}

\hypertarget{h.2et92p0}{\section{\texorpdfstring{{Java Runtime
Environment (JRE)}}{Java Runtime Environment (JRE)}}\label{h.2et92p0}}

{Eclipse depends on a Java so a Java Runtime Environment must be
installed in your system.}

\subsection{\texorpdfstring{{Do I have Java in my
system?}}{Do I have Java in my system?}}\label{do-i-have-java-in-my-system}

{The easy way to check whether you have the JRE installed in your system
is to type call the java interpreter from a windows console.}

\begin{itemize}
\tightlist
\item
  {Run the command shell: ``Start''--\textgreater{}``All
  Programs''--\textgreater{}``Accessories''--\textgreater{}``Command
  Prompt''}
\item
  {Type: ``java -version''}
\item
  {The output should be something similar to this:}
\end{itemize}

{java version ``1.8.0\_40''}

{Java(TM) SE Runtime Environment (build 1.8.0\_40-b26)}

{Java HotSpot(TM) 64-Bit Server VM (build 25.40-b25, mixed mode)}

\subsection{\texorpdfstring{{Installing the
JRE}}{Installing the JRE}}\label{installing-the-jre}

\begin{itemize}
\tightlist
\item
  {Download the JRE from:
  }{\href{https://www.google.com/url?q=http://www.oracle.com/technetwork/java/javase/downloads/index.html\&sa=D\&ust=1511275046373000\&usg=AFQjCNFFJIGcAOFgmFHQHLk6g_3Ml19LyA}{http://www.oracle.com/technetwork/java/javase/downloads/index.html}}{.}
\item
  {Select the correct version for your platform. Most probably the Java
  version you want is the 64 bits version for ~Windows hosts: "}{Windows
  x64}{" --\textgreater{} "}{jre-8u60-windows-x64.exe}{"}
\end{itemize}

\hypertarget{h.tyjcwt}{\section{\texorpdfstring{{Eclipse IDE for C/C++
Developers}}{Eclipse IDE for C/C++ Developers}}\label{h.tyjcwt}}

\begin{itemize}
\tightlist
\item
  {Download the }{Eclipse IDE for C/C++ Developers}{~from
  }{\href{https://www.google.com/url?q=https://eclipse.org/downloads/\&sa=D\&ust=1511275046374000\&usg=AFQjCNGku66E543IPyqaPfqFKQo3pd2L3A}{https://eclipse.org/downloads/}}{.
  At the time of this writing the latest version is }{Mars}{~wich can be
  downloaded directly from:
  }{\href{https://www.google.com/url?q=http://www.eclipse.org/downloads/download.php?file\%3D/technology/epp/downloads/release/mars/R/eclipse-cpp-mars-R-win32-x86_64.zip\&sa=D\&ust=1511275046375000\&usg=AFQjCNGpm_xGDGG0CV-OUsVfXgLuVmq5zw}{http://www.eclipse.org/downloads/download.php?file=/technology/epp/downloads/release/mars/R/eclipse-cpp-mars-R-win32-x86\_64.zip}}
\item
  {Unzip the downloaded file into a directory of your choice. ~For
  example }{C:/eclipse}
\end{itemize}

\hypertarget{h.3dy6vkm}{\section{\texorpdfstring{{GNU ARM Eclipse
Plug-ins}}{GNU ARM Eclipse Plug-ins}}\label{h.3dy6vkm}}

{To install these plug-ins use the Eclipse standard install/update
mechanism}

\begin{itemize}
\tightlist
\item
  {run the Eclipse IDE by selecting the }{eclipse.exe}{~from the
  installation directory. Ex: ~}{C:/eclipse/eclipse.exe}
\item
  {select the suggested default workspace (just click in the
  }{OK}{~button). We will select another workspace to compile the demo
  application later.}
\item
  {in the Eclipse menu select }{Help--\textgreater{}Install New
  Software\ldots{}}
\item
  {in the }{Install}{~window, click the }{Add\ldots{}}{~button}
\item
  {in the }{Add Repository}{~dialog box fill in the fields:}
\end{itemize}

\begin{itemize}
\tightlist
\item
  {Name: }{GNU ARM Eclipse Plug-ins}
\item
  {Location: }{http://gnuarmeclipse.sourceforge.net/updates}
\end{itemize}

{\includegraphics{images/image6.png}}

\begin{itemize}
\tightlist
\item
  {click the OK button}
\item
  {the main window should list a group named CDT GNU Cross Development
  Tools; expand it}
\item
  {select the ``GNU ARM C/C++ Cross Compile'' plug-in}
\end{itemize}

{\includegraphics{images/image9.png}}

\begin{itemize}
\tightlist
\item
  {click the Next button and follow the usual installation procedure}
\end{itemize}

\hypertarget{h.1t3h5sf}{\section{\texorpdfstring{{C/C++ GDB Hardware
Debugging
Plug-in}}{C/C++ GDB Hardware Debugging Plug-in}}\label{h.1t3h5sf}}

\begin{itemize}
\tightlist
\item
  {In the Eclipse menu select Help --\textgreater{}Install New
  Software\ldots{}}
\item
  {in the Work with drop down select Mars -
  http://download.eclipse.org/releases/mars (if your Eclipse version is
  not the Mars, you should find a similar entry with the corresponding
  version).}
\item
  {the main window should list a group named Mobile and Device
  Development; expand it}
\item
  {select the C/C++ GDB Hardware Debugging plug-in}
\end{itemize}

{\includegraphics{images/image8.png}}

\begin{itemize}
\tightlist
\item
  {click the Next button and follow the usual installation procedure}
\end{itemize}

\hypertarget{h.4d34og8}{\section{\texorpdfstring{{STM32 ST-LINK
Utility}}{STM32 ST-LINK Utility}}\label{h.4d34og8}}

{The STM32 ST-LINK Utility software facilitates fast in-system
programming of the STM32 microcontroller families in development
environments via the tools, ST-LINK and ST-LINK/V2. It will be used to
load the }{ThinkOS Bootloader}{~on the target board.}

\begin{itemize}
\tightlist
\item
  {Download }{STSW-LINK004}{~(STM32 ST-LINK utility), the downloaded
  file is called }{stsw-link004.zip}{. The link:
  }{\href{https://www.google.com/url?q=http://www.st.com/web/catalog/tools/FM147/SC1887/PF258168\&sa=D\&ust=1511275046378000\&usg=AFQjCNG92HuOtBhSNJxbtqHYuxedmOi_CA}{http://www.st.com/web/catalog/tools/FM147/SC1887/PF258168}}{~was
  pointing to the file at the time of this writing. Alternatively you
  can search for: }{STM32 ST-LINK utility}{~in the STMicroelectronics
  website:
  }{\href{https://www.google.com/url?q=http://www.st.com\&sa=D\&ust=1511275046378000\&usg=AFQjCNFcwsIyrxx1y2iQSZFPc5tU-A4JIQ}{http://www.st.com}}
\item
  {Decompress the }{stsw-link004.zip}{~file in a temporary folder and
  execute the installer: }{STM32 ST-LINK Utility\_v3.7.0.exe}
\end{itemize}

\hypertarget{h.2s8eyo1}{\section{\texorpdfstring{{ThinkOS Source
Code}}{ThinkOS Source Code}}\label{h.2s8eyo1}}

\begin{itemize}
\tightlist
\item
  {Get the }{ThinkOS}{~source code from GitHub }{YARD-ICE}{~project:
  }{\href{https://www.google.com/url?q=https://github.com/bobmittmann/yard-ice/archive/0.24.zip\&sa=D\&ust=1511275046379000\&usg=AFQjCNH4WK30jeFdsIOws19RGCXF5Ekb5Q}{https://github.com/bobmittmann/yard-ice/archive/0.24.zip}}
\item
  {Extract the source code into a local directory, ex.:
  }{C:\textbackslash{}devel}
\end{itemize}

\subsection{\texorpdfstring{{Cloning the Git
repository}}{Cloning the Git repository}}\label{cloning-the-git-repository}

{Alternatively you can clone the Git repository from:
}{\href{https://www.google.com/url?q=https://github.com/bobmittmann/yard-ice\&sa=D\&ust=1511275046380000\&usg=AFQjCNFusJgGgAORHVsXadKiiwWkzF5HBQ}{https://github.com/bobmittmann/yard-ice}}

\hypertarget{h.17dp8vu}{\section{\texorpdfstring{{Importing Eclipse
Projects}}{Importing Eclipse Projects}}\label{h.17dp8vu}}

{To compile the }{ThinkOS BootLoader}{~and the demo application you must
import the Eclipse projects existing in the ThinkOS source code tree
into your Eclipse workspace.}

\begin{itemize}
\tightlist
\item
  {Select }{board/stm32f3discovery }{as your Eclipse workspace:
  `}{C:\textbackslash{}devel\textbackslash{}yard-ice\textbackslash{}src\textbackslash{}board\textbackslash{}stm32f3discovery}{'}
\item
  {In the Eclipse menu select }{File--\textgreater{}Import\ldots{}}
\item
  {in the }{Import}{~window, expand the }{General}{~group}
\item
  {Select }{Existing Projects into Workspace}{~and click }{Next
  \textgreater{}}
\end{itemize}

{\includegraphics{images/image11.png}}

\begin{itemize}
\tightlist
\item
  {check `}{Select root directory}{'}
\item
  {select the current workspace directory: Ex.:
  `}{C:\textbackslash{}devel\textbackslash{}yard-ice\textbackslash{}src\textbackslash{}board\textbackslash{}stm32f3discovery}{'}
\end{itemize}

{\includegraphics{images/image10.png}}

\begin{itemize}
\tightlist
\item
  {click the ``Finish'' button }
\end{itemize}

\hypertarget{h.3rdcrjn}{\section{\texorpdfstring{{Compiling the ThinkOS
Bootloader}}{Compiling the ThinkOS Bootloader}}\label{h.3rdcrjn}}

\begin{itemize}
\tightlist
\item
  {In the }{Project Explorer}{~tab select the }{boot}{~project. If the
  Project Explorer tab/window is not visible select }{Window}{→}{Show
  View}{→}{Project Explorer}{~from the menu.}
\item
  {Press the build button (Hammer) or select
  }{Project--\textgreater{}Build Project}
\end{itemize}

{\includegraphics{images/image14.png}}

\subsection{\texorpdfstring{{Using the Command
Shell}}{Using the Command Shell}}\label{using-the-command-shell}

{It is not necessary to have the Eclipse environment to compile the
projects in the source tree, you could use the windows Command Prompt or
a MinGW/MSYS terminal directly.}

\begin{itemize}
\tightlist
\item
  {Open a windows shell: Start}{→}{All
  Programs}{→}{Accessories}{→}{Command Prompt}
\item
  {Change to the directory
  src\textbackslash{}board\textbackslash{}stm32f3discovery\textbackslash{}boot
  in the yard-ice source tree.}
\item
  {Type: }{make}
\end{itemize}

{\includegraphics{images/image12.png}}

{}

\hypertarget{h.26in1rg}{\section{\texorpdfstring{{Compiling the board
SDK libraries}}{Compiling the board SDK libraries}}\label{h.26in1rg}}

\begin{itemize}
\tightlist
\item
  {In the }{Project Explorer}{~tab select the }{sdklibs}{~project.}
\item
  {Press the build button (Hammer) or select
  }{Project--\textgreater{}Build Project}
\end{itemize}

\hypertarget{h.lnxbz9}{\section{\texorpdfstring{{Compiling the demo
application}}{Compiling the demo application}}\label{h.lnxbz9}}

\begin{itemize}
\tightlist
\item
  {In the }{Project Explorer}{~tab select the }{led\_test}{~project}
\item
  {Press the build button (Hammer) or select
  }{Project--\textgreater{}Build Project}
\end{itemize}

\hypertarget{h.35nkun2}{\section{\texorpdfstring{{STM32F3Discovery
board}}{STM32F3Discovery board}}\label{h.35nkun2}}

{The }{STM32F3Discovery}{~is a low cost development board for the STM32
F3 series Cortex-M4 mixed-signals. It includes everything required for
beginners and experienced users to get started quickly. Based on the
}{STM32F303VCT6}{, it includes an ST-LINK/V2 embedded debug tool,
accelerometer, gyroscope and e-compass ST MEMS, USB connection, LEDs and
pushbuttons.}

{\includegraphics{images/image13.jpg}}

{The Stm32f3Discovery board has 2 USB connectors:}

\begin{itemize}
\tightlist
\item
  {One at the center of the board marked as }{USB ST-LINK}{~which is
  connected to an auxiliary chip who in turn connects to the JTAG/SWD
  interface of STM303 MPU. We will use this interface to program the
  Boot-Loader once.}
\item
  {The second one, located at the side of the board, marked as }{USB
  USER}{, is attached directly to the STM32F303 MPU. The ThinkOS
  Debug-Monitor will make use of this port as a console for user
  interface and for remote debugging through GDB.}
\end{itemize}

\hypertarget{h.1ksv4uv}{\section{\texorpdfstring{{ThinkOS on
STM32F3Discovery}}{ThinkOS on STM32F3Discovery}}\label{h.1ksv4uv}}

{The ThinkOS can be configured to run as a Boot-Loader and
Debug-Monitor, this way the OS will be resident in the board adding the
support for Application (Firmware) updates as well as debug support.}

{To install or upgrade the ThinkOS in your Stm32f3Discovery board follow
these instructions:}

\begin{itemize}
\tightlist
\item
  {Compile the }{boot}{~project }
\item
  {Install the STM32 ST-LINK Utility}
\item
  {Connect a USB cable to the USB ST-LINK connector in the board and
  into your PC, this will power the board.}
\item
  {Run the STM32 ST-LINK Utility}
\item
  {Connect the program to the board by clicking the connect button or by
  selecting }{Target}{→}{Connect}{~from the menu.}
\item
  {Open the }{thinkos.bin}{~file compiled at step 1, located at
  src\textbackslash{}board\textbackslash{}stm32f3discovery\textbackslash{}boot\textbackslash{}release.
  Use the }{File}{→}{Open file\ldots{}}{~menu.}
\item
  {Program the board by selecting }{Target}{→}{Program\ldots{}}{~and
  pressing the }{Start}{~button on the }{Download {[} thinkos.bin
  {]}}{~dialog box.}
\end{itemize}

{\includegraphics{images/image15.png}}

\begin{itemize}
\tightlist
\item
  {After the programming completed the LEDs on the circular pattern in
  the board will light up in a spinning pattern.}
\end{itemize}

\hypertarget{h.44sinio}{\subsubsection{\texorpdfstring{{USB Virtual COM
Port}}{USB Virtual COM Port}}\label{h.44sinio}}

{You can now remove the mini USB from the central connector (}{USB
ST-LINK}{) and insert it into the adjacent connector at the side of the
board (}{USB USER}{). At this point a virtual serial port driver will be
installed in your system. If the driver fails to install automatically
you need to install it manually. }

\begin{itemize}
\tightlist
\item
  {Download the STSW-STM32102STM32 Virtual COM Port Driver from the
  STMicroelectonics website:
  }{\href{https://www.google.com/url?q=http://www.st.com/web/en/catalog/tools/PF257938\&sa=D\&ust=1511275046387000\&usg=AFQjCNHybx7g4II_reWXNPTvvvAeYlCa4Q}{http://www.st.com/web/en/catalog/tools/PF257938}}{.
  The file should be `stsw-stm32102.zip'}
\item
  {Decompress the file into a temporary folder and run the installation
  program: }{VCP\_V1.4.0\_Setup.exe}{.}
\end{itemize}

{If the driver is installed the Devices and Printers list will show it:
Start}{→Devices and Printers. The device will be displayed in the
}{Unspecified}{~section and is called `}{STMicroelectronics Virtual COM
Port}{'. If you select the device in the list, more details will be
displayed at the bottom. The Model should read: `}{ThinkOS Debug
Monitor}{'.}

{\includegraphics{images/image16.png}}

\hypertarget{h.2jxsxqh}{\section{\texorpdfstring{{ThinkOS
Eclipse/GDB}}{ThinkOS Eclipse/GDB}}\label{h.2jxsxqh}}

{The }{ThinkOS Debug-Monitor}{~is very high priority application
embedded in the }{ThinkOS BootLoader }{which provides a link between the
host GDB debugger and the ThinkOS Kernel, allowing any front end capable
of using GDB for debugging to load and debug Application (Firmware) on
the target board.}

{The }{ThinkOS Debug-Monitor}{~emulates a serial port over the USB to
communicate with the GDB (debugger) running in the PC, you need to make
sure }{ST USB VCOM Port driver}{~is correctly installed in your system
in order to debug }{ThinkOS}{~programs on the Stm32f3Discovery board.}

{This section describes how to setup a remote GDB session using Eclipse
CDT and a }{Stm32f3Discovery}{~board running }{ThinkOS}{.}

\begin{itemize}
\tightlist
\item
  {Make sure the }{Stm32f3Discovery}{~board is attached to the }{USB
  USER}{~interface.}
\item
  {Check the }{Virtual COM Port}{~number. Go to }{Start}{→}{Devices and
  Printers}{~and locate the `}{STMicroelectronics Virtual COM Port}{'.
  The COM port number will be displayed just under the device.}
\item
  {Open the Eclipse project named }{led\_test}{~in the stm32f3discovery
  workspace.}
\item
  {Build the project.}
\item
  {Open the Debug Configurations window selecting }{Run}{→}{Debug
  Configurations\ldots{}}{~from the menu.}
\item
  {Double click the item }{GDB Hardware Debugging}{~or press the `New'
  button.}
\item
  {Select the Debugger tab and fill in the information accordingly:}
\end{itemize}

\begin{itemize}
\tightlist
\item
  {GDB Command: `arm-none-eabi-gdb -b 38400' (please notice the space
  between the --b and the 38400)}
\item
  {JTAG Device: select Generic Serial}
\item
  {GDB Connection String:
  `\textbackslash{}\textbackslash{}.\textbackslash{}COM14'. ~Use here
  the COM port reported in the }{Devices and Printers}{~list.}
\end{itemize}

{\includegraphics{images/image17.png}}

\begin{itemize}
\tightlist
\item
  {Click the `Apply' button}
\item
  {Click the `Debug' button}
\item
  {The Eclipse will change the perspective to Debug and the Debug
  toolbar should appear.}
\item
  {Press the Resume (F8) button (yellow green arrow) to run the
  program.}
\item
  {The LEDs should flash in sequence.}
\item
  {Stop the program by clicking in the pause button. This will list the
  threads of the system and the source code window will indicate the
  current execution line.}
\item
  {At this point you can perform step-by-step execution, insert
  breakpoints, inspect variables, expression registers etc.}
\item
  {When you finish debugging you can close the debug session by pressing
  the Disconnect button (red zigzag). Before starting a new section
  always disconnect the previous one or the GDB will complain it can't
  open the serial/USB port.}
\end{itemize}

{\includegraphics{images/image18.png}}

\section{\texorpdfstring{{Summary of Packages and
Files}}{Summary of Packages and Files}}\label{summary-of-packages-and-files}

\protect\hypertarget{t.36dcbdc507fe00169c271aaebe6543bb601998bb}{}{}\protect\hypertarget{t.0}{}{}

\begin{longtable}[]{@{}ll@{}}
\toprule
\begin{minipage}[t]{0.47\columnwidth}\raggedright\strut
{File}\strut
\end{minipage} & \begin{minipage}[t]{0.47\columnwidth}\raggedright\strut
{Content}\strut
\end{minipage}\tabularnewline
\begin{minipage}[t]{0.47\columnwidth}\raggedright\strut
{stsw-link009.zip }\strut
\end{minipage} & \begin{minipage}[t]{0.47\columnwidth}\raggedright\strut
{ST-LINK USB driver ~ ~ ~ ~ ~ ~ ~ }\strut
\end{minipage}\tabularnewline
\begin{minipage}[t]{0.47\columnwidth}\raggedright\strut
{stsw-link004.zip}\strut
\end{minipage} & \begin{minipage}[t]{0.47\columnwidth}\raggedright\strut
{STM32 ST-LINK utility}\strut
\end{minipage}\tabularnewline
\begin{minipage}[t]{0.47\columnwidth}\raggedright\strut
{stlinkupgrade.zip }\strut
\end{minipage} & \begin{minipage}[t]{0.47\columnwidth}\raggedright\strut
{ST-LINK firmware upgrade}\strut
\end{minipage}\tabularnewline
\begin{minipage}[t]{0.47\columnwidth}\raggedright\strut
{eclipse-cpp-mars-R-win32-x86\_64.zip}\strut
\end{minipage} & \begin{minipage}[t]{0.47\columnwidth}\raggedright\strut
{Eclipse IDE for C/C++ Developers}\strut
\end{minipage}\tabularnewline
\begin{minipage}[t]{0.47\columnwidth}\raggedright\strut
{jre-8u60-windows-x64.exe}\strut
\end{minipage} & \begin{minipage}[t]{0.47\columnwidth}\raggedright\strut
{Java Runtime Environement (JRE)}\strut
\end{minipage}\tabularnewline
\begin{minipage}[t]{0.47\columnwidth}\raggedright\strut
{gcc-arm-none-eabi-4\_9-2015q2-20150609-win32.exe}\strut
\end{minipage} & \begin{minipage}[t]{0.47\columnwidth}\raggedright\strut
{GNU Toolchain for ARM Processros}\strut
\end{minipage}\tabularnewline
\begin{minipage}[t]{0.47\columnwidth}\raggedright\strut
{make-3.81.exe}\strut
\end{minipage} & \begin{minipage}[t]{0.47\columnwidth}\raggedright\strut
{GNU Make for Windows}\strut
\end{minipage}\tabularnewline
\begin{minipage}[t]{0.47\columnwidth}\raggedright\strut
{coreutils-5.3.0.exe}\strut
\end{minipage} & \begin{minipage}[t]{0.47\columnwidth}\raggedright\strut
{GNU CoreUtils for Windows}\strut
\end{minipage}\tabularnewline
\begin{minipage}[t]{0.47\columnwidth}\raggedright\strut
{yard-ice-0.23.zip}\strut
\end{minipage} & \begin{minipage}[t]{0.47\columnwidth}\raggedright\strut
{ThinkOS Source Code}\strut
\end{minipage}\tabularnewline
\begin{minipage}[t]{0.47\columnwidth}\raggedright\strut
{stsw-stm32102.zip}\strut
\end{minipage} & \begin{minipage}[t]{0.47\columnwidth}\raggedright\strut
{STSW-STM32102STM32 Virtual COM Port Driver}\strut
\end{minipage}\tabularnewline
\bottomrule
\end{longtable}

\section{\texorpdfstring{{References}}{References}}\label{references}

\subsubsection{\texorpdfstring{{Gnu Tools for
Window}}{Gnu Tools for Window}}\label{gnu-tools-for-window}

{More information on the GNU Make and CoreUtils package can be found at:
}

\begin{itemize}
\tightlist
\item
  {\href{https://www.google.com/url?q=http://gnuwin32.sourceforge.net/packages/make.htm\&sa=D\&ust=1511275046397000\&usg=AFQjCNHQhDY13EJKv3bQUa54ST0nRARnLg}{http://gnuwin32.sourceforge.net/packages/make.htm}}
\item
  {\href{https://www.google.com/url?q=http://gnuwin32.sourceforge.net/packages/coreutils.htm\&sa=D\&ust=1511275046397000\&usg=AFQjCNHzCBMuNsNxheAyqUiqalsOFlTK9g}{http://gnuwin32.sourceforge.net/packages/coreutils.htm}}
\end{itemize}

\subsubsection{\texorpdfstring{{GNU ARM Eclipse
Plug-Ins}}{GNU ARM Eclipse Plug-Ins}}\label{gnu-arm-eclipse-plug-ins}

{Quick info on GNU ARM Eclipse Plug-Ins:}

\begin{itemize}
\tightlist
\item
  {Name: `}{GNU ARM Eclipse Plug-ins}{'}
\item
  {Location:
  }{\href{https://www.google.com/url?q=http://gnuarmeclipse.sourceforge.net/updates\&sa=D\&ust=1511275046399000\&usg=AFQjCNEdmX4d-t5pswZkr32N4e9t7J324Q}{http://gnuarmeclipse.sourceforge.net/updates}}
\item
  {Website:
  }{\href{https://www.google.com/url?q=http://gnuarmeclipse.livius.net/blog\&sa=D\&ust=1511275046399000\&usg=AFQjCNHwOOZ7pOFg21RgjhMN2SMiP27Adw}{http://gnuarmeclipse.livius.net/blog}}
\end{itemize}

\subsubsection{\texorpdfstring{{ThinkOS}}{ThinkOS}}\label{h.ajl6qmj6j511}

{The ThinkOS source is part of the YARD-ICE project.}

{\includegraphics{images/image2.png}}

\begin{itemize}
\tightlist
\item
  {YARD-ICE Project Location:
  }{\href{https://www.google.com/url?q=https://github.com/bobmittmann/yard-ice\&sa=D\&ust=1511275046400000\&usg=AFQjCNGufoXP2wfG8botuTmr2t7ckLYnvQ}{https://github.com/bobmittmann/yard-ice}}
\item
  {Latest Release:
  }{\href{https://www.google.com/url?q=https://github.com/bobmittmann/yard-ice/archive/0.24.zip\&sa=D\&ust=1511275046401000\&usg=AFQjCNF-nzx-oGprzrbxGIY-uiZXVNkATg}{https://github.com/bobmittmann/yard-ice/archive/0.24.zip}}
\item
  {Git Repository (HTTPS)}{:
  }{https://github.com/bobmittmann/yard-ice.git}
\item
  {Git Repository (SSH): }{git@github.com:bobmittmann/yard-ice.git}
\end{itemize}

{}



\section {API}


\subsection{\texorpdfstring{{Threads}}{Threads}}\label{threads}

{int
}{\protect\hyperlink{h.1fob9te}{thinkos\_thread\_create}}{\protect\hyperlink{h.1fob9te}{~}}{(int(*task\_ptr)(void
*), void *task\_arg, void *stack\_ptr, unsigned int opt)}

{create a new thread }

{int
}{\protect\hyperlink{h.3znysh7}{thinkos\_thread\_create\_inf}}{\protect\hyperlink{h.3znysh7}{~}}{(int(*task\_ptr)(void
*), void *task\_arg, const struct
}{\protect\hyperlink{h.2iq8gzs}{thinkos\_thread\_inf}}{\protect\hyperlink{h.2iq8gzs}{~}}{*inf)}

{create a new thread with extend information }

{int const
}{\protect\hyperlink{h.2et92p0}{thinkos\_thread\_self}}{\protect\hyperlink{h.2et92p0}{~}}{(void)}

{obtaind a handler for the calling thread }

{int
}{\protect\hyperlink{h.tyjcwt}{thinkos\_cancel}}{\protect\hyperlink{h.tyjcwt}{~}}{(unsigned
int thread\_id, int code)}

{request a thread to terminate }

{int
}{\protect\hyperlink{h.3dy6vkm}{thinkos\_exit}}{\protect\hyperlink{h.3dy6vkm}{~}}{(int
code)}

{cause the thread to terminate }

{int
}{\protect\hyperlink{h.1t3h5sf}{thinkos\_join}}{\protect\hyperlink{h.1t3h5sf}{~}}{(unsigned
int thread\_id)}

{join with a terminated thread }

{int
}{\protect\hyperlink{h.4d34og8}{thinkos\_pause}}{\protect\hyperlink{h.4d34og8}{~}}{(unsigned
int thread\_id)}

{pause the thread execution }

{int
}{\protect\hyperlink{h.2s8eyo1}{thinkos\_resume}}{\protect\hyperlink{h.2s8eyo1}{~}}{(unsigned
int thread\_id)}

{resume a thread operation }

{void
}{\protect\hyperlink{h.17dp8vu}{thinkos\_yield}}{\protect\hyperlink{h.17dp8vu}{~}}{(void)}

{causes the calling thread to relinquish the CPU. The thread is moved to
the end of the queue for its static priority and a new thread gets to
run. }

\subsection{\texorpdfstring{{Time related
calls}}{Time related calls}}\label{h.fj4n6rp8qwm7}

\subsubsection{\texorpdfstring{{Functions}}{Functions}}\label{h.iqcme3p75udy}

{int
}{\protect\hyperlink{h.26in1rg}{thinkos\_sleep}}{\protect\hyperlink{h.26in1rg}{~}}{(unsigned
int ms)}

{pause for a specified amount of time }

{uint32\_t
}{\protect\hyperlink{h.lnxbz9}{thinkos\_clock}}{\protect\hyperlink{h.lnxbz9}{~}}{(void)}

{retrieve the ThikOS clock time }

{int
}{\protect\hyperlink{h.35nkun2}{thinkos\_alarm}}{\protect\hyperlink{h.35nkun2}{~}}{(uint32\_t
clock)}

{pause until a specified alarm time }

{}

{}

\subsubsection{\texorpdfstring{{Detailed
Description}}{Detailed Description}}\label{detailed-description}

{}

{}

\subsubsection{\texorpdfstring{{Function
Documentation}}{Function Documentation}}\label{function-documentation}

\paragraph{\texorpdfstring{{int thinkos\_thread\_create (int(*)(void *)
~task\_ptr, void * ~task\_arg, void * ~stack\_ptr, unsigned int
~opt)}}{int thinkos\_thread\_create (int(*)(void *) ~task\_ptr, void * ~task\_arg, void * ~stack\_ptr, unsigned int ~opt)}}\label{int-thinkos_thread_create-intvoid-task_ptr-void-task_arg-void-stack_ptr-unsigned-int-opt}

{}

{\#include
\textless{}}{\protect\hyperlink{h.pkwqa1}{thinkos.h}}{\protect\hyperlink{h.pkwqa1}{\textgreater{}}}

{create a new thread }

{}

\subparagraph{\texorpdfstring{{Parameters:}}{Parameters:}}\label{parameters}

\protect\hypertarget{t.f5730ae642f88aa6ce2571176ce1f624cbde4ea6}{}{}\protect\hypertarget{t.0}{}{}

\begin{longtable}[]{@{}ll@{}}
\toprule
\begin{minipage}[t]{0.47\columnwidth}\raggedright\strut
{task\_ptr}{~}\strut
\end{minipage} & \begin{minipage}[t]{0.47\columnwidth}\raggedright\strut
{}\strut
\end{minipage}\tabularnewline
\begin{minipage}[t]{0.47\columnwidth}\raggedright\strut
{task\_arg}{~}\strut
\end{minipage} & \begin{minipage}[t]{0.47\columnwidth}\raggedright\strut
{}\strut
\end{minipage}\tabularnewline
\begin{minipage}[t]{0.47\columnwidth}\raggedright\strut
{stack\_ptr}{~}\strut
\end{minipage} & \begin{minipage}[t]{0.47\columnwidth}\raggedright\strut
{}\strut
\end{minipage}\tabularnewline
\begin{minipage}[t]{0.47\columnwidth}\raggedright\strut
{opt}{~}\strut
\end{minipage} & \begin{minipage}[t]{0.47\columnwidth}\raggedright\strut
{}\strut
\end{minipage}\tabularnewline
\bottomrule
\end{longtable}

\subparagraph{\texorpdfstring{{Returns:}}{Returns:}}\label{returns}

{\protect\hyperlink{h.3s49zyc}{THINKOS\_ENOSYS}}{\protect\hyperlink{h.3s49zyc}{~}}{if
call is not implemented,
}{\protect\hyperlink{h.2fk6b3p}{THINKOS\_OK}}{\protect\hyperlink{h.2fk6b3p}{~}}{otherwise.
}

\paragraph{\texorpdfstring{{int thinkos\_thread\_create\_inf
(int(*)(void *) ~task\_ptr, void * ~task\_arg, const struct
}{\protect\hyperlink{h.2iq8gzs}{thinkos\_thread\_inf}}{\protect\hyperlink{h.2iq8gzs}{~}}{*
~inf)}}{int thinkos\_thread\_create\_inf (int(*)(void *) ~task\_ptr, void * ~task\_arg, const struct thinkos\_thread\_inf~* ~inf)}}\label{int-thinkos_thread_create_inf-intvoid-task_ptr-void-task_arg-const-struct-thinkos_thread_inf-inf}

{}

{\#include
\textless{}}{\protect\hyperlink{h.pkwqa1}{thinkos.h}}{\protect\hyperlink{h.pkwqa1}{\textgreater{}}}

{create a new thread with extend information }

{}

\subparagraph{\texorpdfstring{{Parameters:}}{Parameters:}}\label{parameters-1}

\protect\hypertarget{t.4f5f8686134758035dc5df693ca2b3c02acdada7}{}{}\protect\hypertarget{t.1}{}{}

\begin{longtable}[]{@{}ll@{}}
\toprule
\begin{minipage}[t]{0.47\columnwidth}\raggedright\strut
{task\_ptr}{~}\strut
\end{minipage} & \begin{minipage}[t]{0.47\columnwidth}\raggedright\strut
{}\strut
\end{minipage}\tabularnewline
\begin{minipage}[t]{0.47\columnwidth}\raggedright\strut
{task\_arg}{~}\strut
\end{minipage} & \begin{minipage}[t]{0.47\columnwidth}\raggedright\strut
{}\strut
\end{minipage}\tabularnewline
\begin{minipage}[t]{0.47\columnwidth}\raggedright\strut
{inf}{~}\strut
\end{minipage} & \begin{minipage}[t]{0.47\columnwidth}\raggedright\strut
{}\strut
\end{minipage}\tabularnewline
\bottomrule
\end{longtable}

\subparagraph{\texorpdfstring{{Returns:}}{Returns:}}\label{returns-1}

{\protect\hyperlink{h.3s49zyc}{THINKOS\_ENOSYS}}{\protect\hyperlink{h.3s49zyc}{~}}{if
call is not implemented,
}{\protect\hyperlink{h.2fk6b3p}{THINKOS\_OK}}{\protect\hyperlink{h.2fk6b3p}{~}}{otherwise.
}

\paragraph{\texorpdfstring{{int const thinkos\_thread\_self (void
)}}{int const thinkos\_thread\_self (void )}}\label{int-const-thinkos_thread_self-void}

{}

{\#include
\textless{}}{\protect\hyperlink{h.pkwqa1}{C:/devel/yard-ice/src/sdk/include/thinkos.h}}{\protect\hyperlink{h.pkwqa1}{\textgreater{}}}

{obtaind a handler for the calling thread }

{}

\subparagraph{\texorpdfstring{{Returns:}}{Returns:}}\label{returns-2}

{This function always succeeds, returning the calling thread's ID. }

\paragraph{}

{\lstinline{int thinkos\_cancel (unsigned int ~thread\_id, int ~code)}\\

{\lstinline{\#include thinkos.h}}\\

\subparagraph{request a thread to terminate }

{}

\subparagraph{\texorpdfstring{{Parameters:}}{Parameters:}}\label{parameters-2}

\protect\hypertarget{t.92b92c82f7d931014629a9c6b6ef3d6c02d24480}{}{}\protect\hypertarget{t.2}{}{}

\begin{longtable}[]{@{}ll@{}}
\toprule
\begin{minipage}[t]{0.47\columnwidth}\raggedright\strut
{thread\_id}{~}\strut
\end{minipage} & \begin{minipage}[t]{0.47\columnwidth}\raggedright\strut
{thread handler }\strut
\end{minipage}\tabularnewline
\begin{minipage}[t]{0.47\columnwidth}\raggedright\strut
{code}{~}\strut
\end{minipage} & \begin{minipage}[t]{0.47\columnwidth}\raggedright\strut
{return code }\strut
\end{minipage}\tabularnewline
\bottomrule
\end{longtable}

\subparagraph{\texorpdfstring{{Returns:}}{Returns:}}\label{returns-3}

{\protect\hyperlink{h.3s49zyc}{THINKOS\_ENOSYS}}{\protect\hyperlink{h.3s49zyc}{~}}{if
call is not implemented,
}{\protect\hyperlink{h.2fk6b3p}{THINKOS\_OK}}{\protect\hyperlink{h.2fk6b3p}{~}}{otherwise.
}

\paragraph{\texorpdfstring{{int thinkos\_exit (int
~code)}}{int thinkos\_exit (int ~code)}}\label{int-thinkos_exit-int-code}

{}

{\#include
\textless{}}{\protect\hyperlink{h.pkwqa1}{C:/devel/yard-ice/src/sdk/include/thinkos.h}}{\protect\hyperlink{h.pkwqa1}{\textgreater{}}}

{cause the thread to terminate }

{}

\subparagraph{\texorpdfstring{{Parameters:}}{Parameters:}}\label{parameters-3}

\protect\hypertarget{t.bc406d26b6ef312038598e48649c4000f5d76cf9}{}{}\protect\hypertarget{t.3}{}{}

\begin{longtable}[]{@{}ll@{}}
\toprule
\begin{minipage}[t]{0.47\columnwidth}\raggedright\strut
{code}{~}\strut
\end{minipage} & \begin{minipage}[t]{0.47\columnwidth}\raggedright\strut
{return code }\strut
\end{minipage}\tabularnewline
\bottomrule
\end{longtable}

\subparagraph{\texorpdfstring{{Returns:}}{Returns:}}\label{returns-4}

{\protect\hyperlink{h.3s49zyc}{THINKOS\_ENOSYS}}{\protect\hyperlink{h.3s49zyc}{~}}{if
call is not implemented,
}{\protect\hyperlink{h.2fk6b3p}{THINKOS\_OK}}{\protect\hyperlink{h.2fk6b3p}{~}}{otherwise.
}

\paragraph{\texorpdfstring{{int thinkos\_join (unsigned int
~thread\_id)}}{int thinkos\_join (unsigned int ~thread\_id)}}\label{int-thinkos_join-unsigned-int-thread_id}

{}

{\#include
\textless{}}{\protect\hyperlink{h.pkwqa1}{C:/devel/yard-ice/src/sdk/include/thinkos.h}}{\protect\hyperlink{h.pkwqa1}{\textgreater{}}}

{join with a terminated thread }

{The
}{\protect\hyperlink{h.1t3h5sf}{thinkos\_join()}}{\protect\hyperlink{h.1t3h5sf}{~}}{function
waits for the thread specified by }{thread\_id}{~ to terminate. If that
thread has already terminated, then
}{\protect\hyperlink{h.1t3h5sf}{thinkos\_join()}}{\protect\hyperlink{h.1t3h5sf}{~}}{returns
immediately.}

\subparagraph{\texorpdfstring{{Parameters:}}{Parameters:}}\label{parameters-4}

\protect\hypertarget{t.d782c80b316c585f37cfeff504ed527240d62369}{}{}\protect\hypertarget{t.4}{}{}

\begin{longtable}[]{@{}ll@{}}
\toprule
\begin{minipage}[t]{0.47\columnwidth}\raggedright\strut
{thread\_id}{~}\strut
\end{minipage} & \begin{minipage}[t]{0.47\columnwidth}\raggedright\strut
{thread handler }\strut
\end{minipage}\tabularnewline
\bottomrule
\end{longtable}

\subparagraph{\texorpdfstring{{Returns:}}{Returns:}}\label{returns-5}

{\protect\hyperlink{h.3s49zyc}{THINKOS\_ENOSYS}}{\protect\hyperlink{h.3s49zyc}{~}}{if
call is not implemented,
}{\protect\hyperlink{h.2fk6b3p}{THINKOS\_OK}}{\protect\hyperlink{h.2fk6b3p}{~}}{otherwise.
}

\paragraph{\texorpdfstring{{int thinkos\_pause (unsigned int
~thread\_id)}}{int thinkos\_pause (unsigned int ~thread\_id)}}\label{int-thinkos_pause-unsigned-int-thread_id}

{}

{\#include
\textless{}}{\protect\hyperlink{h.pkwqa1}{C:/devel/yard-ice/src/sdk/include/thinkos.h}}{\protect\hyperlink{h.pkwqa1}{\textgreater{}}}

{pause the thread execution }

{}

\subparagraph{\texorpdfstring{{Parameters:}}{Parameters:}}\label{parameters-5}

\protect\hypertarget{t.d782c80b316c585f37cfeff504ed527240d62369}{}{}\protect\hypertarget{t.5}{}{}

\begin{longtable}[]{@{}ll@{}}
\toprule
\begin{minipage}[t]{0.47\columnwidth}\raggedright\strut
{thread\_id}{~}\strut
\end{minipage} & \begin{minipage}[t]{0.47\columnwidth}\raggedright\strut
{thread handler }\strut
\end{minipage}\tabularnewline
\bottomrule
\end{longtable}

\subparagraph{\texorpdfstring{{Returns:}}{Returns:}}\label{returns-6}

{\protect\hyperlink{h.3s49zyc}{THINKOS\_ENOSYS}}{\protect\hyperlink{h.3s49zyc}{~}}{if
call is not implemented,
}{\protect\hyperlink{h.2fk6b3p}{THINKOS\_OK}}{\protect\hyperlink{h.2fk6b3p}{~}}{otherwise.
}

\paragraph{\texorpdfstring{{int thinkos\_resume (unsigned int
~thread\_id)}}{int thinkos\_resume (unsigned int ~thread\_id)}}\label{int-thinkos_resume-unsigned-int-thread_id}

{}

{\#include
\textless{}}{\protect\hyperlink{h.pkwqa1}{C:/devel/yard-ice/src/sdk/include/thinkos.h}}{\protect\hyperlink{h.pkwqa1}{\textgreater{}}}

{resume a thread operation }

{}

\subparagraph{\texorpdfstring{{Parameters:}}{Parameters:}}\label{parameters-6}

\protect\hypertarget{t.d782c80b316c585f37cfeff504ed527240d62369}{}{}\protect\hypertarget{t.6}{}{}

\begin{longtable}[]{@{}ll@{}}
\toprule
\begin{minipage}[t]{0.47\columnwidth}\raggedright\strut
{thread\_id}{~}\strut
\end{minipage} & \begin{minipage}[t]{0.47\columnwidth}\raggedright\strut
{thread handler }\strut
\end{minipage}\tabularnewline
\bottomrule
\end{longtable}

\subparagraph{\texorpdfstring{{Returns:}}{Returns:}}\label{returns-7}

{\protect\hyperlink{h.3s49zyc}{THINKOS\_ENOSYS}}{\protect\hyperlink{h.3s49zyc}{~}}{if
call is not implemented,
}{\protect\hyperlink{h.2fk6b3p}{THINKOS\_OK}}{\protect\hyperlink{h.2fk6b3p}{~}}{otherwise.
}

\paragraph{\texorpdfstring{{void thinkos\_yield (void
)}}{void thinkos\_yield (void )}}\label{void-thinkos_yield-void}

{}

{\#include
\textless{}}{\protect\hyperlink{h.pkwqa1}{C:/devel/yard-ice/src/sdk/include/thinkos.h}}{\protect\hyperlink{h.pkwqa1}{\textgreater{}}}

{causes the calling thread to relinquish the CPU. The thread is moved to
the end of the queue for its static priority and a new thread gets to
run. }

\paragraph{\texorpdfstring{{}}{}}\label{section}

\subsection{\texorpdfstring{{Time related
calls}}{Time related calls}}\label{time-related-calls}

\subsubsection{\texorpdfstring{{Functions}}{Functions}}\label{h.3rdcrjn}

{int
}{\protect\hyperlink{h.26in1rg}{thinkos\_sleep}}{\protect\hyperlink{h.26in1rg}{~}}{(unsigned
int ms)}

{pause for a specified amount of time }

{uint32\_t
}{\protect\hyperlink{h.lnxbz9}{thinkos\_clock}}{\protect\hyperlink{h.lnxbz9}{~}}{(void)}

{retrieve the ThikOS clock time }

{int
}{\protect\hyperlink{h.35nkun2}{thinkos\_alarm}}{\protect\hyperlink{h.35nkun2}{~}}{(uint32\_t
clock)}

{pause until a specified alarm time }

{}

\subsubsection{\texorpdfstring{{Detailed
Description}}{Detailed Description}}\label{detailed-description-1}

{}

{}

\subsubsection{\texorpdfstring{{Function
Documentation}}{Function Documentation}}\label{function-documentation-1}

\paragraph{\texorpdfstring{{int thinkos\_sleep (unsigned int
~ms)}}{int thinkos\_sleep (unsigned int ~ms)}}\label{int-thinkos_sleep-unsigned-int-ms}

{}

{\#include
\textless{}}{\protect\hyperlink{h.pkwqa1}{C:/devel/yard-ice/src/sdk/include/thinkos.h}}{\protect\hyperlink{h.pkwqa1}{\textgreater{}}}

{pause for a specified amount of time }

{The
}{\protect\hyperlink{h.26in1rg}{thinkos\_sleep()}}{\protect\hyperlink{h.26in1rg}{~}}{function
suspends execution of the calling thread for (at least) ms milliseconds.
The sleep may be lengthened slightly by any system activity or by the
time spent processing the call or by the granularity of system timer.}

\subparagraph{\texorpdfstring{{Parameters:}}{Parameters:}}\label{parameters-7}

\protect\hypertarget{t.e4727bf18c12225c1cdcee7b27c6f0d762fa7610}{}{}\protect\hypertarget{t.7}{}{}

\begin{longtable}[]{@{}ll@{}}
\toprule
\begin{minipage}[t]{0.47\columnwidth}\raggedright\strut
{ms}{~}\strut
\end{minipage} & \begin{minipage}[t]{0.47\columnwidth}\raggedright\strut
{waiting time in milliseconds }\strut
\end{minipage}\tabularnewline
\bottomrule
\end{longtable}

\subparagraph{\texorpdfstring{{Returns:}}{Returns:}}\label{returns-8}

{\protect\hyperlink{h.3s49zyc}{THINKOS\_ENOSYS}}{\protect\hyperlink{h.3s49zyc}{~}}{if
call is not implemented,
}{\protect\hyperlink{h.2fk6b3p}{THINKOS\_OK}}{\protect\hyperlink{h.2fk6b3p}{~}}{otherwise.
}

\paragraph{\texorpdfstring{{uint32\_t thinkos\_clock (void
)}}{uint32\_t thinkos\_clock (void )}}\label{uint32_t-thinkos_clock-void}

{}

{\#include
\textless{}}{\protect\hyperlink{h.pkwqa1}{C:/devel/yard-ice/src/sdk/include/thinkos.h}}{\protect\hyperlink{h.pkwqa1}{\textgreater{}}}

{retrieve the ThikOS clock time }

{The ThikOS clock is a tick counter with a resolution of 1 milisecond,
which cannot be set and is allways running. It represents monotonic time
since the system power up.}

\subparagraph{\texorpdfstring{{Returns:}}{Returns:}}\label{returns-9}

{the ThinkOS monotonic clock. }

\paragraph{\texorpdfstring{{int thinkos\_alarm (uint32\_t
~clock)}}{int thinkos\_alarm (uint32\_t ~clock)}}\label{int-thinkos_alarm-uint32_t-clock}

{}

{\#include
\textless{}}{\protect\hyperlink{h.pkwqa1}{C:/devel/yard-ice/src/sdk/include/thinkos.h}}{\protect\hyperlink{h.pkwqa1}{\textgreater{}}}

{pause until a specified alarm time }

{This function should be used in conjucntion with
}{\protect\hyperlink{h.lnxbz9}{thinkos\_clock()}}{\protect\hyperlink{h.lnxbz9}{.}}{~\ldots{}}

\subparagraph{\texorpdfstring{{Parameters:}}{Parameters:}}\label{parameters-8}

\protect\hypertarget{t.281362b1a9b4665596996d7811a599a33fd3578c}{}{}\protect\hypertarget{t.8}{}{}

\begin{longtable}[]{@{}ll@{}}
\toprule
\begin{minipage}[t]{0.47\columnwidth}\raggedright\strut
{clock}{~}\strut
\end{minipage} & \begin{minipage}[t]{0.47\columnwidth}\raggedright\strut
{time for the thread to be waked up. }\strut
\end{minipage}\tabularnewline
\bottomrule
\end{longtable}

\subparagraph{\texorpdfstring{{Returns:}}{Returns:}}\label{returns-10}

{\protect\hyperlink{h.3s49zyc}{THINKOS\_ENOSYS}}{\protect\hyperlink{h.3s49zyc}{~}}{if
call is not implemented,
}{\protect\hyperlink{h.2fk6b3p}{THINKOS\_OK}}{\protect\hyperlink{h.2fk6b3p}{~}}{otherwise.
}

\paragraph{\texorpdfstring{{}}{}}\label{section-1}

\subsection{\texorpdfstring{{Mutexes}}{Mutexes}}\label{mutexes}

\subsubsection{\texorpdfstring{{Functions}}{Functions}}\label{h.1ksv4uv}

{int
}{\protect\hyperlink{h.44sinio}{thinkos\_mutex\_alloc}}{\protect\hyperlink{h.44sinio}{~}}{(void)}

{alloc a mutex }

{int
}{\protect\hyperlink{h.2jxsxqh}{thinkos\_mutex\_free}}{\protect\hyperlink{h.2jxsxqh}{~}}{(int
mutex)}

{release a mutex }

{int
}{\protect\hyperlink{h.z337ya}{thinkos\_mutex\_lock}}{\protect\hyperlink{h.z337ya}{~}}{(int
mutex)}

{lock a mutex, blocking if already locked. }

{int
}{\protect\hyperlink{h.3j2qqm3}{thinkos\_mutex\_trylock}}{\protect\hyperlink{h.3j2qqm3}{~}}{(int
mutex)}

{try to lock a mutex, non blocking. }

{int
}{\protect\hyperlink{h.1y810tw}{thinkos\_mutex\_timedlock}}{\protect\hyperlink{h.1y810tw}{~}}{(int
mutex, unsigned int ms)}

{lock a mutex, blocking for limited time. }

{int
}{\protect\hyperlink{h.4i7ojhp}{thinkos\_mutex\_unlock}}{\protect\hyperlink{h.4i7ojhp}{~}}{(int
mutex)}

{unlock the mutex }

{}

\subsubsection{\texorpdfstring{{Detailed
Description}}{Detailed Description}}\label{detailed-description-2}

{}

{}

\subsubsection{\texorpdfstring{{Function
Documentation}}{Function Documentation}}\label{function-documentation-2}

\paragraph{\texorpdfstring{{int thinkos\_mutex\_alloc (void
)}}{int thinkos\_mutex\_alloc (void )}}\label{int-thinkos_mutex_alloc-void}

{}

{\#include
\textless{}}{\protect\hyperlink{h.pkwqa1}{C:/devel/yard-ice/src/sdk/include/thinkos.h}}{\protect\hyperlink{h.pkwqa1}{\textgreater{}}}

{alloc a mutex }

{}

\subparagraph{\texorpdfstring{{Returns:}}{Returns:}}\label{returns-11}

{\protect\hyperlink{h.3s49zyc}{THINKOS\_ENOSYS}}{\protect\hyperlink{h.3s49zyc}{~}}{if
call is not implemented,
}{\protect\hyperlink{h.2fk6b3p}{THINKOS\_OK}}{\protect\hyperlink{h.2fk6b3p}{~}}{otherwise.
}

\paragraph{\texorpdfstring{{int thinkos\_mutex\_free (int
~mutex)}}{int thinkos\_mutex\_free (int ~mutex)}}\label{int-thinkos_mutex_free-int-mutex}

{}

{\#include
\textless{}}{\protect\hyperlink{h.pkwqa1}{C:/devel/yard-ice/src/sdk/include/thinkos.h}}{\protect\hyperlink{h.pkwqa1}{\textgreater{}}}

{release a mutex }

{}

\subparagraph{\texorpdfstring{{Parameters:}}{Parameters:}}\label{parameters-9}

\protect\hypertarget{t.b0b883f74d2462ca881bd7f4d0508cf144bbde9b}{}{}\protect\hypertarget{t.9}{}{}

\begin{longtable}[]{@{}ll@{}}
\toprule
\begin{minipage}[t]{0.47\columnwidth}\raggedright\strut
{mutex}{~}\strut
\end{minipage} & \begin{minipage}[t]{0.47\columnwidth}\raggedright\strut
{mutual exclusion handler }\strut
\end{minipage}\tabularnewline
\bottomrule
\end{longtable}

\subparagraph{\texorpdfstring{{Returns:}}{Returns:}}\label{returns-12}

{\protect\hyperlink{h.3s49zyc}{THINKOS\_ENOSYS}}{\protect\hyperlink{h.3s49zyc}{~}}{if
call is not implemented,
}{\protect\hyperlink{h.2fk6b3p}{THINKOS\_OK}}{\protect\hyperlink{h.2fk6b3p}{~}}{otherwise.
}

\paragraph{\texorpdfstring{{int thinkos\_mutex\_lock (int
~mutex)}}{int thinkos\_mutex\_lock (int ~mutex)}}\label{int-thinkos_mutex_lock-int-mutex}

{}

{\#include
\textless{}}{\protect\hyperlink{h.pkwqa1}{C:/devel/yard-ice/src/sdk/include/thinkos.h}}{\protect\hyperlink{h.pkwqa1}{\textgreater{}}}

{lock a mutex, blocking if already locked. }

{}

\subparagraph{\texorpdfstring{{Parameters:}}{Parameters:}}\label{parameters-10}

\protect\hypertarget{t.b0b883f74d2462ca881bd7f4d0508cf144bbde9b}{}{}\protect\hypertarget{t.10}{}{}

\begin{longtable}[]{@{}ll@{}}
\toprule
\begin{minipage}[t]{0.47\columnwidth}\raggedright\strut
{mutex}{~}\strut
\end{minipage} & \begin{minipage}[t]{0.47\columnwidth}\raggedright\strut
{mutual exclusion handler }\strut
\end{minipage}\tabularnewline
\bottomrule
\end{longtable}

\subparagraph{\texorpdfstring{{Returns:}}{Returns:}}\label{returns-13}

{\protect\hyperlink{h.3s49zyc}{THINKOS\_ENOSYS}}{\protect\hyperlink{h.3s49zyc}{~}}{if
call is not implemented,
}{\protect\hyperlink{h.2fk6b3p}{THINKOS\_OK}}{\protect\hyperlink{h.2fk6b3p}{~}}{otherwise.
}

\paragraph{\texorpdfstring{{int thinkos\_mutex\_trylock (int
~mutex)}}{int thinkos\_mutex\_trylock (int ~mutex)}}\label{int-thinkos_mutex_trylock-int-mutex}

{}

{\#include
\textless{}}{\protect\hyperlink{h.pkwqa1}{C:/devel/yard-ice/src/sdk/include/thinkos.h}}{\protect\hyperlink{h.pkwqa1}{\textgreater{}}}

{try to lock a mutex, non blocking. }

{}

\subparagraph{\texorpdfstring{{Parameters:}}{Parameters:}}\label{parameters-11}

\protect\hypertarget{t.b0b883f74d2462ca881bd7f4d0508cf144bbde9b}{}{}\protect\hypertarget{t.11}{}{}

\begin{longtable}[]{@{}ll@{}}
\toprule
\begin{minipage}[t]{0.47\columnwidth}\raggedright\strut
{mutex}{~}\strut
\end{minipage} & \begin{minipage}[t]{0.47\columnwidth}\raggedright\strut
{mutual exclusion handler }\strut
\end{minipage}\tabularnewline
\bottomrule
\end{longtable}

\subparagraph{\texorpdfstring{{Returns:}}{Returns:}}\label{returns-14}

{\protect\hyperlink{h.3s49zyc}{THINKOS\_ENOSYS}}{\protect\hyperlink{h.3s49zyc}{~}}{if
call is not implemented,
}{\protect\hyperlink{h.2fk6b3p}{THINKOS\_OK}}{\protect\hyperlink{h.2fk6b3p}{~}}{otherwise.
}

\paragraph{\texorpdfstring{{int thinkos\_mutex\_timedlock (int ~mutex,
unsigned int
~ms)}}{int thinkos\_mutex\_timedlock (int ~mutex, unsigned int ~ms)}}\label{int-thinkos_mutex_timedlock-int-mutex-unsigned-int-ms}

{}

{\#include
\textless{}}{\protect\hyperlink{h.pkwqa1}{C:/devel/yard-ice/src/sdk/include/thinkos.h}}{\protect\hyperlink{h.pkwqa1}{\textgreater{}}}

{lock a mutex, blocking for limited time. }

{}

\subparagraph{\texorpdfstring{{Parameters:}}{Parameters:}}\label{parameters-12}

\protect\hypertarget{t.de941ac38142c0150a822438d9a2a2aec068d18f}{}{}\protect\hypertarget{t.12}{}{}

\begin{longtable}[]{@{}ll@{}}
\toprule
\begin{minipage}[t]{0.47\columnwidth}\raggedright\strut
{mutex}{~}\strut
\end{minipage} & \begin{minipage}[t]{0.47\columnwidth}\raggedright\strut
{mutual exclusion handler }\strut
\end{minipage}\tabularnewline
\begin{minipage}[t]{0.47\columnwidth}\raggedright\strut
{ms}{~}\strut
\end{minipage} & \begin{minipage}[t]{0.47\columnwidth}\raggedright\strut
{waiting time in milliseconds }\strut
\end{minipage}\tabularnewline
\bottomrule
\end{longtable}

\subparagraph{\texorpdfstring{{Returns:}}{Returns:}}\label{returns-15}

{\protect\hyperlink{h.3s49zyc}{THINKOS\_ENOSYS}}{\protect\hyperlink{h.3s49zyc}{~}}{if
call is not implemented,
}{\protect\hyperlink{h.2fk6b3p}{THINKOS\_OK}}{\protect\hyperlink{h.2fk6b3p}{~}}{otherwise.
}

\paragraph{\texorpdfstring{{int thinkos\_mutex\_unlock (int
~mutex)}}{int thinkos\_mutex\_unlock (int ~mutex)}}\label{int-thinkos_mutex_unlock-int-mutex}

{}

{\#include
\textless{}}{\protect\hyperlink{h.pkwqa1}{C:/devel/yard-ice/src/sdk/include/thinkos.h}}{\protect\hyperlink{h.pkwqa1}{\textgreater{}}}

{unlock the mutex }

{}

\subparagraph{\texorpdfstring{{Parameters:}}{Parameters:}}\label{parameters-13}

\protect\hypertarget{t.b0b883f74d2462ca881bd7f4d0508cf144bbde9b}{}{}\protect\hypertarget{t.13}{}{}

\begin{longtable}[]{@{}ll@{}}
\toprule
\begin{minipage}[t]{0.47\columnwidth}\raggedright\strut
{mutex}{~}\strut
\end{minipage} & \begin{minipage}[t]{0.47\columnwidth}\raggedright\strut
{mutual exclusion handler }\strut
\end{minipage}\tabularnewline
\bottomrule
\end{longtable}

\subparagraph{\texorpdfstring{{Returns:}}{Returns:}}\label{returns-16}

{\protect\hyperlink{h.3s49zyc}{THINKOS\_ENOSYS}}{\protect\hyperlink{h.3s49zyc}{~}}{if
call is not implemented,
}{\protect\hyperlink{h.2fk6b3p}{THINKOS\_OK}}{\protect\hyperlink{h.2fk6b3p}{~}}{otherwise.
}

\paragraph{\texorpdfstring{{}}{}}\label{section-2}

\subsection{\texorpdfstring{{Conditional
Variables}}{Conditional Variables}}\label{conditional-variables}

\subsubsection{\texorpdfstring{{Functions}}{Functions}}\label{h.2xcytpi}

{int
}{\protect\hyperlink{h.1ci93xb}{thinkos\_cond\_alloc}}{\protect\hyperlink{h.1ci93xb}{~}}{(void)}

{alloc a conditional variable }

{int
}{\protect\hyperlink{h.3whwml4}{thinkos\_cond\_free}}{\protect\hyperlink{h.3whwml4}{~}}{(int
cond)}

{release a conditional variable }

{int
}{\protect\hyperlink{h.2bn6wsx}{thinkos\_cond\_wait}}{\protect\hyperlink{h.2bn6wsx}{~}}{(int
cond, int mutex)}

{wait for a conditional variable }

{int
}{\protect\hyperlink{h.qsh70q}{thinkos\_cond\_timedwait}}{\protect\hyperlink{h.qsh70q}{~}}{(int
cond, int mutex, unsigned int ms)}

{wait for a conditional variable or timeout }

{int
}{\protect\hyperlink{h.3as4poj}{thinkos\_cond\_signal}}{\protect\hyperlink{h.3as4poj}{~}}{(int
cond)}

{signal a conditional variable wake a single thread }

{int
}{\protect\hyperlink{h.1pxezwc}{thinkos\_cond\_broadcast}}{\protect\hyperlink{h.1pxezwc}{~}}{(int
cond)}

{signal a conditional variable wake all waiting threads }

{}

\subsubsection{\texorpdfstring{{Detailed
Description}}{Detailed Description}}\label{detailed-description-3}

{}

{}

\subsubsection{\texorpdfstring{{Function
Documentation}}{Function Documentation}}\label{function-documentation-3}

\paragraph{\texorpdfstring{{int thinkos\_cond\_alloc (void
)}}{int thinkos\_cond\_alloc (void )}}\label{int-thinkos_cond_alloc-void}

{}

{\#include
\textless{}}{\protect\hyperlink{h.pkwqa1}{C:/devel/yard-ice/src/sdk/include/thinkos.h}}{\protect\hyperlink{h.pkwqa1}{\textgreater{}}}

{alloc a conditional variable }

{}

\subparagraph{\texorpdfstring{{Returns:}}{Returns:}}\label{returns-17}

{\protect\hyperlink{h.3s49zyc}{THINKOS\_ENOSYS}}{\protect\hyperlink{h.3s49zyc}{~}}{if
call is not implemented,
}{\protect\hyperlink{h.2fk6b3p}{THINKOS\_OK}}{\protect\hyperlink{h.2fk6b3p}{~}}{otherwise.
}

\paragraph{\texorpdfstring{{int thinkos\_cond\_free (int
~cond)}}{int thinkos\_cond\_free (int ~cond)}}\label{int-thinkos_cond_free-int-cond}

{}

{\#include
\textless{}}{\protect\hyperlink{h.pkwqa1}{C:/devel/yard-ice/src/sdk/include/thinkos.h}}{\protect\hyperlink{h.pkwqa1}{\textgreater{}}}

{release a conditional variable }

{}

\subparagraph{\texorpdfstring{{Parameters:}}{Parameters:}}\label{parameters-14}

\protect\hypertarget{t.5f6a6301995a69995f1253a58a8df876856672ec}{}{}\protect\hypertarget{t.14}{}{}

\begin{longtable}[]{@{}ll@{}}
\toprule
\begin{minipage}[t]{0.47\columnwidth}\raggedright\strut
{cond}{~}\strut
\end{minipage} & \begin{minipage}[t]{0.47\columnwidth}\raggedright\strut
{conditional variable handler }\strut
\end{minipage}\tabularnewline
\bottomrule
\end{longtable}

\subparagraph{\texorpdfstring{{Returns:}}{Returns:}}\label{returns-18}

{\protect\hyperlink{h.3s49zyc}{THINKOS\_ENOSYS}}{\protect\hyperlink{h.3s49zyc}{~}}{if
call is not implemented,
}{\protect\hyperlink{h.2fk6b3p}{THINKOS\_OK}}{\protect\hyperlink{h.2fk6b3p}{~}}{otherwise.
}

\paragraph{\texorpdfstring{{int thinkos\_cond\_wait (int ~cond, int
~mutex)}}{int thinkos\_cond\_wait (int ~cond, int ~mutex)}}\label{int-thinkos_cond_wait-int-cond-int-mutex}

{}

{\#include
\textless{}}{\protect\hyperlink{h.pkwqa1}{C:/devel/yard-ice/src/sdk/include/thinkos.h}}{\protect\hyperlink{h.pkwqa1}{\textgreater{}}}

{wait for a conditional variable }

{}

\subparagraph{\texorpdfstring{{Parameters:}}{Parameters:}}\label{parameters-15}

\protect\hypertarget{t.abe173c86a250b55a35b20a1a938b3e0ed603c37}{}{}\protect\hypertarget{t.15}{}{}

\begin{longtable}[]{@{}ll@{}}
\toprule
\begin{minipage}[t]{0.47\columnwidth}\raggedright\strut
{cond}{~}\strut
\end{minipage} & \begin{minipage}[t]{0.47\columnwidth}\raggedright\strut
{conditional variable handler }\strut
\end{minipage}\tabularnewline
\begin{minipage}[t]{0.47\columnwidth}\raggedright\strut
{mutex}{~}\strut
\end{minipage} & \begin{minipage}[t]{0.47\columnwidth}\raggedright\strut
{mutex handler }\strut
\end{minipage}\tabularnewline
\bottomrule
\end{longtable}

\subparagraph{\texorpdfstring{{Returns:}}{Returns:}}\label{returns-19}

{\protect\hyperlink{h.3s49zyc}{THINKOS\_ENOSYS}}{\protect\hyperlink{h.3s49zyc}{~}}{if
call is not implemented,
}{\protect\hyperlink{h.2fk6b3p}{THINKOS\_OK}}{\protect\hyperlink{h.2fk6b3p}{~}}{otherwise.
}

\paragraph{\texorpdfstring{{int thinkos\_cond\_timedwait (int ~cond, int
~mutex, unsigned int
~ms)}}{int thinkos\_cond\_timedwait (int ~cond, int ~mutex, unsigned int ~ms)}}\label{int-thinkos_cond_timedwait-int-cond-int-mutex-unsigned-int-ms}

{}

{\#include
\textless{}}{\protect\hyperlink{h.pkwqa1}{C:/devel/yard-ice/src/sdk/include/thinkos.h}}{\protect\hyperlink{h.pkwqa1}{\textgreater{}}}

{wait for a conditional variable or timeout }

{}

\subparagraph{\texorpdfstring{{Parameters:}}{Parameters:}}\label{parameters-16}

\protect\hypertarget{t.858b18386d0696f6efe16a7a32d5dcded6652bc5}{}{}\protect\hypertarget{t.16}{}{}

\begin{longtable}[]{@{}ll@{}}
\toprule
\begin{minipage}[t]{0.47\columnwidth}\raggedright\strut
{cond}{~}\strut
\end{minipage} & \begin{minipage}[t]{0.47\columnwidth}\raggedright\strut
{conditional variable handler }\strut
\end{minipage}\tabularnewline
\begin{minipage}[t]{0.47\columnwidth}\raggedright\strut
{mutex}{~}\strut
\end{minipage} & \begin{minipage}[t]{0.47\columnwidth}\raggedright\strut
{mutex handler }\strut
\end{minipage}\tabularnewline
\begin{minipage}[t]{0.47\columnwidth}\raggedright\strut
{ms}{~}\strut
\end{minipage} & \begin{minipage}[t]{0.47\columnwidth}\raggedright\strut
{waiting time in milliseconds }\strut
\end{minipage}\tabularnewline
\bottomrule
\end{longtable}

\subparagraph{\texorpdfstring{{Returns:}}{Returns:}}\label{returns-20}

{\protect\hyperlink{h.3s49zyc}{THINKOS\_ENOSYS}}{\protect\hyperlink{h.3s49zyc}{~}}{if
call is not implemented,
}{\protect\hyperlink{h.2fk6b3p}{THINKOS\_OK}}{\protect\hyperlink{h.2fk6b3p}{~}}{otherwise.
}

\paragraph{\texorpdfstring{{int thinkos\_cond\_signal (int
~cond)}}{int thinkos\_cond\_signal (int ~cond)}}\label{int-thinkos_cond_signal-int-cond}

{}

{\#include
\textless{}}{\protect\hyperlink{h.pkwqa1}{C:/devel/yard-ice/src/sdk/include/thinkos.h}}{\protect\hyperlink{h.pkwqa1}{\textgreater{}}}

{signal a conditional variable wake a single thread }

{}

\subparagraph{\texorpdfstring{{Parameters:}}{Parameters:}}\label{parameters-17}

\protect\hypertarget{t.5f6a6301995a69995f1253a58a8df876856672ec}{}{}\protect\hypertarget{t.17}{}{}

\begin{longtable}[]{@{}ll@{}}
\toprule
\begin{minipage}[t]{0.47\columnwidth}\raggedright\strut
{cond}{~}\strut
\end{minipage} & \begin{minipage}[t]{0.47\columnwidth}\raggedright\strut
{conditional variable handler }\strut
\end{minipage}\tabularnewline
\bottomrule
\end{longtable}

\subparagraph{\texorpdfstring{{Returns:}}{Returns:}}\label{returns-21}

{\protect\hyperlink{h.3s49zyc}{THINKOS\_ENOSYS}}{\protect\hyperlink{h.3s49zyc}{~}}{if
call is not implemented,
}{\protect\hyperlink{h.2fk6b3p}{THINKOS\_OK}}{\protect\hyperlink{h.2fk6b3p}{~}}{otherwise.
}

\paragraph{\texorpdfstring{{int thinkos\_cond\_broadcast (int
~cond)}}{int thinkos\_cond\_broadcast (int ~cond)}}\label{int-thinkos_cond_broadcast-int-cond}

{}

{\#include
\textless{}}{\protect\hyperlink{h.pkwqa1}{C:/devel/yard-ice/src/sdk/include/thinkos.h}}{\protect\hyperlink{h.pkwqa1}{\textgreater{}}}

{signal a conditional variable wake all waiting threads }

{}

\subparagraph{\texorpdfstring{{Parameters:}}{Parameters:}}\label{parameters-18}

\protect\hypertarget{t.5f6a6301995a69995f1253a58a8df876856672ec}{}{}\protect\hypertarget{t.18}{}{}

\begin{longtable}[]{@{}ll@{}}
\toprule
\begin{minipage}[t]{0.47\columnwidth}\raggedright\strut
{cond}{~}\strut
\end{minipage} & \begin{minipage}[t]{0.47\columnwidth}\raggedright\strut
{conditional variable handler }\strut
\end{minipage}\tabularnewline
\bottomrule
\end{longtable}

\subparagraph{\texorpdfstring{{Returns:}}{Returns:}}\label{returns-22}

{\protect\hyperlink{h.3s49zyc}{THINKOS\_ENOSYS}}{\protect\hyperlink{h.3s49zyc}{~}}{if
call is not implemented,
}{\protect\hyperlink{h.2fk6b3p}{THINKOS\_OK}}{\protect\hyperlink{h.2fk6b3p}{~}}{otherwise.
}

\paragraph{\texorpdfstring{{}}{}}\label{section-3}

\subsection{\texorpdfstring{{Semaphores}}{Semaphores}}\label{semaphores}

\subsubsection{\texorpdfstring{{Functions}}{Functions}}\label{h.49x2ik5}

{int
}{\protect\hyperlink{h.2p2csry}{thinkos\_sem\_alloc}}{\protect\hyperlink{h.2p2csry}{~}}{(unsigned
int val)}

{alloc a semaphore }

{int
}{\protect\hyperlink{h.147n2zr}{thinkos\_sem\_free}}{\protect\hyperlink{h.147n2zr}{~}}{(int
sem)}

{release a semaphore }

{int
}{\protect\hyperlink{h.3o7alnk}{thinkos\_sem\_init}}{\protect\hyperlink{h.3o7alnk}{~}}{(int
sem, unsigned int val)}

{set the initial value of a semaphore }

{int
}{\protect\hyperlink{h.23ckvvd}{thinkos\_sem\_wait}}{\protect\hyperlink{h.23ckvvd}{~}}{(int
sem)}

{wait for a semaphore to be signaled }

{int
}{\protect\hyperlink{h.ihv636}{thinkos\_sem\_timedwait}}{\protect\hyperlink{h.ihv636}{~}}{(int
sem, unsigned int ms)}

{wait for a semaphore to be signaled or timeout }

{int
}{\protect\hyperlink{h.32hioqz}{thinkos\_sem\_post}}{\protect\hyperlink{h.32hioqz}{~}}{(int
sem)}

{signal a semaphore }

{void
}{\protect\hyperlink{h.1hmsyys}{thinkos\_sem\_post\_i}}{\protect\hyperlink{h.1hmsyys}{~}}{(int
sem)}

{signal a semaphore inside an interrupt handler }

{}

\subsubsection{\texorpdfstring{{Detailed
Description}}{Detailed Description}}\label{detailed-description-4}

{}

{}

\subsubsection{\texorpdfstring{{Function
Documentation}}{Function Documentation}}\label{function-documentation-4}

\paragraph{\texorpdfstring{{int thinkos\_sem\_alloc (unsigned int
~val)}}{int thinkos\_sem\_alloc (unsigned int ~val)}}\label{int-thinkos_sem_alloc-unsigned-int-val}

{}

{\#include
\textless{}}{\protect\hyperlink{h.pkwqa1}{C:/devel/yard-ice/src/sdk/include/thinkos.h}}{\protect\hyperlink{h.pkwqa1}{\textgreater{}}}

{alloc a semaphore }

{}

\subparagraph{\texorpdfstring{{Parameters:}}{Parameters:}}\label{parameters-19}

\protect\hypertarget{t.b02b07ec80924b1785f96fcda4f10f7c9466fa7b}{}{}\protect\hypertarget{t.19}{}{}

\begin{longtable}[]{@{}ll@{}}
\toprule
\begin{minipage}[t]{0.47\columnwidth}\raggedright\strut
{val}{~}\strut
\end{minipage} & \begin{minipage}[t]{0.47\columnwidth}\raggedright\strut
{initial semaphore value }\strut
\end{minipage}\tabularnewline
\bottomrule
\end{longtable}

\subparagraph{\texorpdfstring{{Returns:}}{Returns:}}\label{returns-23}

{\protect\hyperlink{h.3s49zyc}{THINKOS\_ENOSYS}}{\protect\hyperlink{h.3s49zyc}{~}}{if
call is not implemented,
}{\protect\hyperlink{h.2fk6b3p}{THINKOS\_OK}}{\protect\hyperlink{h.2fk6b3p}{~}}{otherwise.
}

\paragraph{\texorpdfstring{{int thinkos\_sem\_free (int
~sem)}}{int thinkos\_sem\_free (int ~sem)}}\label{int-thinkos_sem_free-int-sem}

{}

{\#include
\textless{}}{\protect\hyperlink{h.pkwqa1}{C:/devel/yard-ice/src/sdk/include/thinkos.h}}{\protect\hyperlink{h.pkwqa1}{\textgreater{}}}

{release a semaphore }

{}

\subparagraph{\texorpdfstring{{Parameters:}}{Parameters:}}\label{parameters-20}

\protect\hypertarget{t.42f0ad63e27678006eae2b4c683934157041a4d7}{}{}\protect\hypertarget{t.20}{}{}

\begin{longtable}[]{@{}ll@{}}
\toprule
\begin{minipage}[t]{0.47\columnwidth}\raggedright\strut
{sem}{~}\strut
\end{minipage} & \begin{minipage}[t]{0.47\columnwidth}\raggedright\strut
{semaphore handler }\strut
\end{minipage}\tabularnewline
\bottomrule
\end{longtable}

\subparagraph{\texorpdfstring{{Returns:}}{Returns:}}\label{returns-24}

{\protect\hyperlink{h.3s49zyc}{THINKOS\_ENOSYS}}{\protect\hyperlink{h.3s49zyc}{~}}{if
call is not implemented,
}{\protect\hyperlink{h.2fk6b3p}{THINKOS\_OK}}{\protect\hyperlink{h.2fk6b3p}{~}}{otherwise.
}

\paragraph{\texorpdfstring{{int thinkos\_sem\_init (int ~sem, unsigned
int
~val)}}{int thinkos\_sem\_init (int ~sem, unsigned int ~val)}}\label{int-thinkos_sem_init-int-sem-unsigned-int-val}

{}

{\#include
\textless{}}{\protect\hyperlink{h.pkwqa1}{C:/devel/yard-ice/src/sdk/include/thinkos.h}}{\protect\hyperlink{h.pkwqa1}{\textgreater{}}}

{set the initial value of a semaphore }

{}

\subparagraph{\texorpdfstring{{Parameters:}}{Parameters:}}\label{parameters-21}

\protect\hypertarget{t.4fd8384b48c431f3cad5f382b558ebd8eed50495}{}{}\protect\hypertarget{t.21}{}{}

\begin{longtable}[]{@{}ll@{}}
\toprule
\begin{minipage}[t]{0.47\columnwidth}\raggedright\strut
{sem}{~}\strut
\end{minipage} & \begin{minipage}[t]{0.47\columnwidth}\raggedright\strut
{semaphore handler }\strut
\end{minipage}\tabularnewline
\begin{minipage}[t]{0.47\columnwidth}\raggedright\strut
{val}{~}\strut
\end{minipage} & \begin{minipage}[t]{0.47\columnwidth}\raggedright\strut
{}\strut
\end{minipage}\tabularnewline
\bottomrule
\end{longtable}

\subparagraph{\texorpdfstring{{Returns:}}{Returns:}}\label{returns-25}

{\protect\hyperlink{h.3s49zyc}{THINKOS\_ENOSYS}}{\protect\hyperlink{h.3s49zyc}{~}}{if
call is not implemented,
}{\protect\hyperlink{h.2fk6b3p}{THINKOS\_OK}}{\protect\hyperlink{h.2fk6b3p}{~}}{otherwise.
}

\paragraph{\texorpdfstring{{int thinkos\_sem\_wait (int
~sem)}}{int thinkos\_sem\_wait (int ~sem)}}\label{int-thinkos_sem_wait-int-sem}

{}

{\#include
\textless{}}{\protect\hyperlink{h.pkwqa1}{C:/devel/yard-ice/src/sdk/include/thinkos.h}}{\protect\hyperlink{h.pkwqa1}{\textgreater{}}}

{wait for a semaphore to be signaled }

{}

\subparagraph{\texorpdfstring{{Parameters:}}{Parameters:}}\label{parameters-22}

\protect\hypertarget{t.42f0ad63e27678006eae2b4c683934157041a4d7}{}{}\protect\hypertarget{t.22}{}{}

\begin{longtable}[]{@{}ll@{}}
\toprule
\begin{minipage}[t]{0.47\columnwidth}\raggedright\strut
{sem}{~}\strut
\end{minipage} & \begin{minipage}[t]{0.47\columnwidth}\raggedright\strut
{semaphore handler }\strut
\end{minipage}\tabularnewline
\bottomrule
\end{longtable}

\subparagraph{\texorpdfstring{{Returns:}}{Returns:}}\label{returns-26}

{\protect\hyperlink{h.3s49zyc}{THINKOS\_ENOSYS}}{\protect\hyperlink{h.3s49zyc}{~}}{if
call is not implemented,
}{\protect\hyperlink{h.2fk6b3p}{THINKOS\_OK}}{\protect\hyperlink{h.2fk6b3p}{~}}{otherwise.
}

\paragraph{\texorpdfstring{{int thinkos\_sem\_timedwait (int ~sem,
unsigned int
~ms)}}{int thinkos\_sem\_timedwait (int ~sem, unsigned int ~ms)}}\label{int-thinkos_sem_timedwait-int-sem-unsigned-int-ms}

{}

{\#include
\textless{}}{\protect\hyperlink{h.pkwqa1}{C:/devel/yard-ice/src/sdk/include/thinkos.h}}{\protect\hyperlink{h.pkwqa1}{\textgreater{}}}

{wait for a semaphore to be signaled or timeout }

{}

\subparagraph{\texorpdfstring{{Parameters:}}{Parameters:}}\label{parameters-23}

\protect\hypertarget{t.33a353b3a299d454126b7888c7a887c8e8d7a8ee}{}{}\protect\hypertarget{t.23}{}{}

\begin{longtable}[]{@{}ll@{}}
\toprule
\begin{minipage}[t]{0.47\columnwidth}\raggedright\strut
{sem}{~}\strut
\end{minipage} & \begin{minipage}[t]{0.47\columnwidth}\raggedright\strut
{semaphore handler }\strut
\end{minipage}\tabularnewline
\begin{minipage}[t]{0.47\columnwidth}\raggedright\strut
{ms}{~}\strut
\end{minipage} & \begin{minipage}[t]{0.47\columnwidth}\raggedright\strut
{}\strut
\end{minipage}\tabularnewline
\bottomrule
\end{longtable}

\subparagraph{\texorpdfstring{{Returns:}}{Returns:}}\label{returns-27}

{\protect\hyperlink{h.3s49zyc}{THINKOS\_ENOSYS}}{\protect\hyperlink{h.3s49zyc}{~}}{if
call is not implemented,
}{\protect\hyperlink{h.2fk6b3p}{THINKOS\_OK}}{\protect\hyperlink{h.2fk6b3p}{~}}{otherwise.
}

\paragraph{\texorpdfstring{{int thinkos\_sem\_post (int
~sem)}}{int thinkos\_sem\_post (int ~sem)}}\label{int-thinkos_sem_post-int-sem}

{}

{\#include
\textless{}}{\protect\hyperlink{h.pkwqa1}{C:/devel/yard-ice/src/sdk/include/thinkos.h}}{\protect\hyperlink{h.pkwqa1}{\textgreater{}}}

{signal a semaphore }

{}

\subparagraph{\texorpdfstring{{Parameters:}}{Parameters:}}\label{parameters-24}

\protect\hypertarget{t.42f0ad63e27678006eae2b4c683934157041a4d7}{}{}\protect\hypertarget{t.24}{}{}

\begin{longtable}[]{@{}ll@{}}
\toprule
\begin{minipage}[t]{0.47\columnwidth}\raggedright\strut
{sem}{~}\strut
\end{minipage} & \begin{minipage}[t]{0.47\columnwidth}\raggedright\strut
{semaphore handler }\strut
\end{minipage}\tabularnewline
\bottomrule
\end{longtable}

\subparagraph{\texorpdfstring{{Returns:}}{Returns:}}\label{returns-28}

{\protect\hyperlink{h.3s49zyc}{THINKOS\_ENOSYS}}{\protect\hyperlink{h.3s49zyc}{~}}{if
call is not implemented,
}{\protect\hyperlink{h.2fk6b3p}{THINKOS\_OK}}{\protect\hyperlink{h.2fk6b3p}{~}}{otherwise.
}

\paragraph{\texorpdfstring{{void thinkos\_sem\_post\_i (int
~sem)}}{void thinkos\_sem\_post\_i (int ~sem)}}\label{void-thinkos_sem_post_i-int-sem}

{}

{\#include
\textless{}}{\protect\hyperlink{h.pkwqa1}{C:/devel/yard-ice/src/sdk/include/thinkos.h}}{\protect\hyperlink{h.pkwqa1}{\textgreater{}}}

{signal a semaphore inside an interrupt handler }

{}

\subparagraph{\texorpdfstring{{Parameters:}}{Parameters:}}\label{parameters-25}

\protect\hypertarget{t.42f0ad63e27678006eae2b4c683934157041a4d7}{}{}\protect\hypertarget{t.25}{}{}

\begin{longtable}[]{@{}ll@{}}
\toprule
\begin{minipage}[t]{0.47\columnwidth}\raggedright\strut
{sem}{~}\strut
\end{minipage} & \begin{minipage}[t]{0.47\columnwidth}\raggedright\strut
{semaphore handler }\strut
\end{minipage}\tabularnewline
\bottomrule
\end{longtable}

\paragraph{\texorpdfstring{{}}{}}\label{section-4}

\subsection{\texorpdfstring{{Event sets}}{Event sets}}\label{event-sets}

\subsubsection{\texorpdfstring{{Functions}}{Functions}}\label{h.41mghml}

{int
}{\protect\hyperlink{h.2grqrue}{thinkos\_ev\_alloc}}{\protect\hyperlink{h.2grqrue}{~}}{(void)}

{alloc an event set }

{int
}{\protect\hyperlink{h.vx1227}{thinkos\_ev\_free}}{\protect\hyperlink{h.vx1227}{~}}{(int
set)}

{release an event set }

{int
}{\protect\hyperlink{h.3fwokq0}{thinkos\_ev\_wait}}{\protect\hyperlink{h.3fwokq0}{~}}{(int
set)}

{wait for an event }

{int
}{\protect\hyperlink{h.1v1yuxt}{thinkos\_ev\_timedwait}}{\protect\hyperlink{h.1v1yuxt}{~}}{(int
set, unsigned int ms)}

{wait for an event or timeout }

{int
}{\protect\hyperlink{h.4f1mdlm}{thinkos\_ev\_raise}}{\protect\hyperlink{h.4f1mdlm}{~}}{(int
set, int ev)}

{signal an event }

{void
}{\protect\hyperlink{h.2u6wntf}{thinkos\_ev\_raise\_i}}{\protect\hyperlink{h.2u6wntf}{~}}{(int
set, int ev)}

{signal an event from inside an interrupt handler }

{int
}{\protect\hyperlink{h.19c6y18}{thinkos\_ev\_mask}}{\protect\hyperlink{h.19c6y18}{~}}{(int
set, int ev, int val)}

{mask/unmask an event }

{int
}{\protect\hyperlink{h.3tbugp1}{thinkos\_ev\_clear}}{\protect\hyperlink{h.3tbugp1}{~}}{(int
set, int ev)}

{clear an event }

{}

\subsubsection{\texorpdfstring{{Detailed
Description}}{Detailed Description}}\label{detailed-description-5}

{}

{}

\subsubsection{\texorpdfstring{{Function
Documentation}}{Function Documentation}}\label{function-documentation-5}

\paragraph{\texorpdfstring{{int thinkos\_ev\_alloc (void
)}}{int thinkos\_ev\_alloc (void )}}\label{int-thinkos_ev_alloc-void}

{}

{\#include
\textless{}}{\protect\hyperlink{h.pkwqa1}{C:/devel/yard-ice/src/sdk/include/thinkos.h}}{\protect\hyperlink{h.pkwqa1}{\textgreater{}}}

{alloc an event set }

{}

\subparagraph{\texorpdfstring{{Returns:}}{Returns:}}\label{returns-29}

{\protect\hyperlink{h.3s49zyc}{THINKOS\_ENOSYS}}{\protect\hyperlink{h.3s49zyc}{~}}{if
call is not implemented,
}{\protect\hyperlink{h.2fk6b3p}{THINKOS\_OK}}{\protect\hyperlink{h.2fk6b3p}{~}}{otherwise.
}

\paragraph{\texorpdfstring{{int thinkos\_ev\_free (int
~set)}}{int thinkos\_ev\_free (int ~set)}}\label{int-thinkos_ev_free-int-set}

{}

{\#include
\textless{}}{\protect\hyperlink{h.pkwqa1}{C:/devel/yard-ice/src/sdk/include/thinkos.h}}{\protect\hyperlink{h.pkwqa1}{\textgreater{}}}

{release an event set }

{}

\subparagraph{\texorpdfstring{{Parameters:}}{Parameters:}}\label{parameters-26}

\protect\hypertarget{t.fa7a662ba3715d766286afd3239f3de6307c79a3}{}{}\protect\hypertarget{t.26}{}{}

\begin{longtable}[]{@{}ll@{}}
\toprule
\begin{minipage}[t]{0.47\columnwidth}\raggedright\strut
{set}{~}\strut
\end{minipage} & \begin{minipage}[t]{0.47\columnwidth}\raggedright\strut
{event set handler }\strut
\end{minipage}\tabularnewline
\bottomrule
\end{longtable}

\subparagraph{\texorpdfstring{{Returns:}}{Returns:}}\label{returns-30}

{\protect\hyperlink{h.3s49zyc}{THINKOS\_ENOSYS}}{\protect\hyperlink{h.3s49zyc}{~}}{if
call is not implemented,
}{\protect\hyperlink{h.2fk6b3p}{THINKOS\_OK}}{\protect\hyperlink{h.2fk6b3p}{~}}{otherwise.
}

\paragraph{\texorpdfstring{{int thinkos\_ev\_wait (int
~set)}}{int thinkos\_ev\_wait (int ~set)}}\label{int-thinkos_ev_wait-int-set}

{}

{\#include
\textless{}}{\protect\hyperlink{h.pkwqa1}{C:/devel/yard-ice/src/sdk/include/thinkos.h}}{\protect\hyperlink{h.pkwqa1}{\textgreater{}}}

{wait for an event }

{}

\subparagraph{\texorpdfstring{{Parameters:}}{Parameters:}}\label{parameters-27}

\protect\hypertarget{t.fa7a662ba3715d766286afd3239f3de6307c79a3}{}{}\protect\hypertarget{t.27}{}{}

\begin{longtable}[]{@{}ll@{}}
\toprule
\begin{minipage}[t]{0.47\columnwidth}\raggedright\strut
{set}{~}\strut
\end{minipage} & \begin{minipage}[t]{0.47\columnwidth}\raggedright\strut
{event set handler }\strut
\end{minipage}\tabularnewline
\bottomrule
\end{longtable}

\subparagraph{\texorpdfstring{{Returns:}}{Returns:}}\label{returns-31}

{\protect\hyperlink{h.3s49zyc}{THINKOS\_ENOSYS}}{\protect\hyperlink{h.3s49zyc}{~}}{if
call is not implemented, an event otherwise. }

\paragraph{\texorpdfstring{{int thinkos\_ev\_timedwait (int ~set,
unsigned int
~ms)}}{int thinkos\_ev\_timedwait (int ~set, unsigned int ~ms)}}\label{int-thinkos_ev_timedwait-int-set-unsigned-int-ms}

{}

{\#include
\textless{}}{\protect\hyperlink{h.pkwqa1}{C:/devel/yard-ice/src/sdk/include/thinkos.h}}{\protect\hyperlink{h.pkwqa1}{\textgreater{}}}

{wait for an event or timeout }

{}

\subparagraph{\texorpdfstring{{Parameters:}}{Parameters:}}\label{parameters-28}

\protect\hypertarget{t.798629b4d033e81359313f6211f814d0e4c1651b}{}{}\protect\hypertarget{t.28}{}{}

\begin{longtable}[]{@{}ll@{}}
\toprule
\begin{minipage}[t]{0.47\columnwidth}\raggedright\strut
{set}{~}\strut
\end{minipage} & \begin{minipage}[t]{0.47\columnwidth}\raggedright\strut
{event set handler }\strut
\end{minipage}\tabularnewline
\begin{minipage}[t]{0.47\columnwidth}\raggedright\strut
{ms}{~}\strut
\end{minipage} & \begin{minipage}[t]{0.47\columnwidth}\raggedright\strut
{time to wait in milliseconds }\strut
\end{minipage}\tabularnewline
\bottomrule
\end{longtable}

\subparagraph{\texorpdfstring{{Returns:}}{Returns:}}\label{returns-32}

{\protect\hyperlink{h.3s49zyc}{THINKOS\_ENOSYS}}{\protect\hyperlink{h.3s49zyc}{~}}{if
call is not implemented, an event otherwise. }

\paragraph{\texorpdfstring{{int thinkos\_ev\_raise (int ~set, int
~ev)}}{int thinkos\_ev\_raise (int ~set, int ~ev)}}\label{int-thinkos_ev_raise-int-set-int-ev}

{}

{\#include
\textless{}}{\protect\hyperlink{h.pkwqa1}{C:/devel/yard-ice/src/sdk/include/thinkos.h}}{\protect\hyperlink{h.pkwqa1}{\textgreater{}}}

{signal an event }

{}

\subparagraph{\texorpdfstring{{Parameters:}}{Parameters:}}\label{parameters-29}

\protect\hypertarget{t.edc44c26eda9f56c94bb61cc6c572ae53f165af7}{}{}\protect\hypertarget{t.29}{}{}

\begin{longtable}[]{@{}ll@{}}
\toprule
\begin{minipage}[t]{0.47\columnwidth}\raggedright\strut
{set}{~}\strut
\end{minipage} & \begin{minipage}[t]{0.47\columnwidth}\raggedright\strut
{event set handler }\strut
\end{minipage}\tabularnewline
\begin{minipage}[t]{0.47\columnwidth}\raggedright\strut
{ev}{~}\strut
\end{minipage} & \begin{minipage}[t]{0.47\columnwidth}\raggedright\strut
{event identifier }\strut
\end{minipage}\tabularnewline
\bottomrule
\end{longtable}

\subparagraph{\texorpdfstring{{Returns:}}{Returns:}}\label{returns-33}

{\protect\hyperlink{h.3s49zyc}{THINKOS\_ENOSYS}}{\protect\hyperlink{h.3s49zyc}{~}}{if
call is not implemented,
}{\protect\hyperlink{h.2fk6b3p}{THINKOS\_OK}}{\protect\hyperlink{h.2fk6b3p}{~}}{otherwise.
}

\paragraph{\texorpdfstring{{void thinkos\_ev\_raise\_i (int ~set, int
~ev)}}{void thinkos\_ev\_raise\_i (int ~set, int ~ev)}}\label{void-thinkos_ev_raise_i-int-set-int-ev}

{}

{\#include
\textless{}}{\protect\hyperlink{h.pkwqa1}{C:/devel/yard-ice/src/sdk/include/thinkos.h}}{\protect\hyperlink{h.pkwqa1}{\textgreater{}}}

{signal an event from inside an interrupt handler }

{}

\subparagraph{\texorpdfstring{{Parameters:}}{Parameters:}}\label{parameters-30}

\protect\hypertarget{t.edc44c26eda9f56c94bb61cc6c572ae53f165af7}{}{}\protect\hypertarget{t.30}{}{}

\begin{longtable}[]{@{}ll@{}}
\toprule
\begin{minipage}[t]{0.47\columnwidth}\raggedright\strut
{set}{~}\strut
\end{minipage} & \begin{minipage}[t]{0.47\columnwidth}\raggedright\strut
{event set handler }\strut
\end{minipage}\tabularnewline
\begin{minipage}[t]{0.47\columnwidth}\raggedright\strut
{ev}{~}\strut
\end{minipage} & \begin{minipage}[t]{0.47\columnwidth}\raggedright\strut
{event identifier }\strut
\end{minipage}\tabularnewline
\bottomrule
\end{longtable}

\paragraph{\texorpdfstring{{int thinkos\_ev\_mask (int ~set, int ~ev,
int
~val)}}{int thinkos\_ev\_mask (int ~set, int ~ev, int ~val)}}\label{int-thinkos_ev_mask-int-set-int-ev-int-val}

{}

{\#include
\textless{}}{\protect\hyperlink{h.pkwqa1}{C:/devel/yard-ice/src/sdk/include/thinkos.h}}{\protect\hyperlink{h.pkwqa1}{\textgreater{}}}

{mask/unmask an event }

{}

\subparagraph{\texorpdfstring{{Parameters:}}{Parameters:}}\label{parameters-31}

\protect\hypertarget{t.067f774366820f43dbfd2abca19aca2a5f35fd2e}{}{}\protect\hypertarget{t.31}{}{}

\begin{longtable}[]{@{}ll@{}}
\toprule
\begin{minipage}[t]{0.47\columnwidth}\raggedright\strut
{set}{~}\strut
\end{minipage} & \begin{minipage}[t]{0.47\columnwidth}\raggedright\strut
{event set handler }\strut
\end{minipage}\tabularnewline
\begin{minipage}[t]{0.47\columnwidth}\raggedright\strut
{ev}{~}\strut
\end{minipage} & \begin{minipage}[t]{0.47\columnwidth}\raggedright\strut
{event identifier }\strut
\end{minipage}\tabularnewline
\begin{minipage}[t]{0.47\columnwidth}\raggedright\strut
{val}{~}\strut
\end{minipage} & \begin{minipage}[t]{0.47\columnwidth}\raggedright\strut
{mask value. 1 enable the event, 0 disable the event. }\strut
\end{minipage}\tabularnewline
\bottomrule
\end{longtable}

\subparagraph{\texorpdfstring{{Returns:}}{Returns:}}\label{returns-34}

{\protect\hyperlink{h.3s49zyc}{THINKOS\_ENOSYS}}{\protect\hyperlink{h.3s49zyc}{~}}{if
call is not implemented,
}{\protect\hyperlink{h.2fk6b3p}{THINKOS\_OK}}{\protect\hyperlink{h.2fk6b3p}{~}}{otherwise.
}

\paragraph{\texorpdfstring{{int thinkos\_ev\_clear (int ~set, int
~ev)}}{int thinkos\_ev\_clear (int ~set, int ~ev)}}\label{int-thinkos_ev_clear-int-set-int-ev}

{}

{\#include
\textless{}}{\protect\hyperlink{h.pkwqa1}{C:/devel/yard-ice/src/sdk/include/thinkos.h}}{\protect\hyperlink{h.pkwqa1}{\textgreater{}}}

{clear an event }

{}

\subparagraph{\texorpdfstring{{Parameters:}}{Parameters:}}\label{parameters-32}

\protect\hypertarget{t.edc44c26eda9f56c94bb61cc6c572ae53f165af7}{}{}\protect\hypertarget{t.32}{}{}

\begin{longtable}[]{@{}ll@{}}
\toprule
\begin{minipage}[t]{0.47\columnwidth}\raggedright\strut
{set}{~}\strut
\end{minipage} & \begin{minipage}[t]{0.47\columnwidth}\raggedright\strut
{event set handler }\strut
\end{minipage}\tabularnewline
\begin{minipage}[t]{0.47\columnwidth}\raggedright\strut
{ev}{~}\strut
\end{minipage} & \begin{minipage}[t]{0.47\columnwidth}\raggedright\strut
{event identifier }\strut
\end{minipage}\tabularnewline
\bottomrule
\end{longtable}

\subparagraph{\texorpdfstring{{Returns:}}{Returns:}}\label{returns-35}

{\protect\hyperlink{h.3s49zyc}{THINKOS\_ENOSYS}}{\protect\hyperlink{h.3s49zyc}{~}}{if
call is not implemented,
}{\protect\hyperlink{h.2fk6b3p}{THINKOS\_OK}}{\protect\hyperlink{h.2fk6b3p}{~}}{otherwise.
}

\paragraph{\texorpdfstring{{}}{}}\label{section-5}

\subsection{\texorpdfstring{{Flags}}{Flags}}\label{flags}

\subsubsection{\texorpdfstring{{Functions}}{Functions}}\label{h.28h4qwu}

{int
}{\protect\hyperlink{h.nmf14n}{thinkos\_flag\_alloc}}{\protect\hyperlink{h.nmf14n}{~}}{(void)}

{alloc a flag }

{int
}{\protect\hyperlink{h.37m2jsg}{thinkos\_flag\_free}}{\protect\hyperlink{h.37m2jsg}{~}}{(int
flag)}

{release a flag }

{int
}{\protect\hyperlink{h.1mrcu09}{thinkos\_flag\_val}}{\protect\hyperlink{h.1mrcu09}{~}}{(int
flag)}

{return the flag value }

{int
}{\protect\hyperlink{h.46r0co2}{thinkos\_flag\_set}}{\protect\hyperlink{h.46r0co2}{~}}{(int
flag)}

{set a flag }

{int
}{\protect\hyperlink{h.2lwamvv}{thinkos\_flag\_clr}}{\protect\hyperlink{h.2lwamvv}{~}}{(int
flag)}

{clear a flag }

{int
}{\protect\hyperlink{h.111kx3o}{thinkos\_flag\_give}}{\protect\hyperlink{h.111kx3o}{~}}{(int
flag)}

{signal a flag }

{void
}{\protect\hyperlink{h.3l18frh}{thinkos\_flag\_give\_i}}{\protect\hyperlink{h.3l18frh}{~}}{(int
flag)}

{signal a flag form inside an interrupt handler }

{int
}{\protect\hyperlink{h.206ipza}{thinkos\_flag\_take}}{\protect\hyperlink{h.206ipza}{~}}{(int
flag)}

{wait for flag to be signaled }

{int
}{\protect\hyperlink{h.4k668n3}{thinkos\_flag\_timedtake}}{\protect\hyperlink{h.4k668n3}{~}}{(int
flag, unsigned int ms)}

{wait for flag or timeout }

{}

\subsubsection{\texorpdfstring{{Detailed
Description}}{Detailed Description}}\label{detailed-description-6}

{}

{}

\subsubsection{\texorpdfstring{{Function
Documentation}}{Function Documentation}}\label{function-documentation-6}

\paragraph{\texorpdfstring{{int thinkos\_flag\_alloc (void
)}}{int thinkos\_flag\_alloc (void )}}\label{int-thinkos_flag_alloc-void}

{}

{\#include
\textless{}}{\protect\hyperlink{h.pkwqa1}{C:/devel/yard-ice/src/sdk/include/thinkos.h}}{\protect\hyperlink{h.pkwqa1}{\textgreater{}}}

{alloc a flag }

{}

\subparagraph{\texorpdfstring{{Returns:}}{Returns:}}\label{returns-36}

{\protect\hyperlink{h.3s49zyc}{THINKOS\_ENOSYS}}{\protect\hyperlink{h.3s49zyc}{~}}{if
call is not implemented,
}{\protect\hyperlink{h.2fk6b3p}{THINKOS\_OK}}{\protect\hyperlink{h.2fk6b3p}{~}}{otherwise.
}

\paragraph{\texorpdfstring{{int thinkos\_flag\_free (int
~flag)}}{int thinkos\_flag\_free (int ~flag)}}\label{int-thinkos_flag_free-int-flag}

{}

{\#include
\textless{}}{\protect\hyperlink{h.pkwqa1}{C:/devel/yard-ice/src/sdk/include/thinkos.h}}{\protect\hyperlink{h.pkwqa1}{\textgreater{}}}

{release a flag }

{}

\subparagraph{\texorpdfstring{{Parameters:}}{Parameters:}}\label{parameters-33}

\protect\hypertarget{t.383f7e723be615a1ab251cbb171dc44f8c48c9f4}{}{}\protect\hypertarget{t.33}{}{}

\begin{longtable}[]{@{}ll@{}}
\toprule
\begin{minipage}[t]{0.47\columnwidth}\raggedright\strut
{flag}{~}\strut
\end{minipage} & \begin{minipage}[t]{0.47\columnwidth}\raggedright\strut
{}\strut
\end{minipage}\tabularnewline
\bottomrule
\end{longtable}

\subparagraph{\texorpdfstring{{Returns:}}{Returns:}}\label{returns-37}

{\protect\hyperlink{h.3s49zyc}{THINKOS\_ENOSYS}}{\protect\hyperlink{h.3s49zyc}{~}}{if
call is not implemented,
}{\protect\hyperlink{h.2fk6b3p}{THINKOS\_OK}}{\protect\hyperlink{h.2fk6b3p}{~}}{otherwise.
}

\paragraph{\texorpdfstring{{int thinkos\_flag\_val (int
~flag)}}{int thinkos\_flag\_val (int ~flag)}}\label{int-thinkos_flag_val-int-flag}

{}

{\#include
\textless{}}{\protect\hyperlink{h.pkwqa1}{C:/devel/yard-ice/src/sdk/include/thinkos.h}}{\protect\hyperlink{h.pkwqa1}{\textgreater{}}}

{return the flag value }

{}

\subparagraph{\texorpdfstring{{Parameters:}}{Parameters:}}\label{parameters-34}

\protect\hypertarget{t.383f7e723be615a1ab251cbb171dc44f8c48c9f4}{}{}\protect\hypertarget{t.34}{}{}

\begin{longtable}[]{@{}ll@{}}
\toprule
\begin{minipage}[t]{0.47\columnwidth}\raggedright\strut
{flag}{~}\strut
\end{minipage} & \begin{minipage}[t]{0.47\columnwidth}\raggedright\strut
{}\strut
\end{minipage}\tabularnewline
\bottomrule
\end{longtable}

\subparagraph{\texorpdfstring{{Returns:}}{Returns:}}\label{returns-38}

{\protect\hyperlink{h.3s49zyc}{THINKOS\_ENOSYS}}{\protect\hyperlink{h.3s49zyc}{~}}{if
call is not implemented,
}{\protect\hyperlink{h.2fk6b3p}{THINKOS\_OK}}{\protect\hyperlink{h.2fk6b3p}{~}}{otherwise.
}

\paragraph{\texorpdfstring{{int thinkos\_flag\_set (int
~flag)}}{int thinkos\_flag\_set (int ~flag)}}\label{int-thinkos_flag_set-int-flag}

{}

{\#include
\textless{}}{\protect\hyperlink{h.pkwqa1}{C:/devel/yard-ice/src/sdk/include/thinkos.h}}{\protect\hyperlink{h.pkwqa1}{\textgreater{}}}

{set a flag }

{}

\subparagraph{\texorpdfstring{{Parameters:}}{Parameters:}}\label{parameters-35}

\protect\hypertarget{t.383f7e723be615a1ab251cbb171dc44f8c48c9f4}{}{}\protect\hypertarget{t.35}{}{}

\begin{longtable}[]{@{}ll@{}}
\toprule
\begin{minipage}[t]{0.47\columnwidth}\raggedright\strut
{flag}{~}\strut
\end{minipage} & \begin{minipage}[t]{0.47\columnwidth}\raggedright\strut
{}\strut
\end{minipage}\tabularnewline
\bottomrule
\end{longtable}

\subparagraph{\texorpdfstring{{Returns:}}{Returns:}}\label{returns-39}

{\protect\hyperlink{h.3s49zyc}{THINKOS\_ENOSYS}}{\protect\hyperlink{h.3s49zyc}{~}}{if
call is not implemented,
}{\protect\hyperlink{h.2fk6b3p}{THINKOS\_OK}}{\protect\hyperlink{h.2fk6b3p}{~}}{otherwise.
}

\paragraph{\texorpdfstring{{int thinkos\_flag\_clr (int
~flag)}}{int thinkos\_flag\_clr (int ~flag)}}\label{int-thinkos_flag_clr-int-flag}

{}

{\#include
\textless{}}{\protect\hyperlink{h.pkwqa1}{C:/devel/yard-ice/src/sdk/include/thinkos.h}}{\protect\hyperlink{h.pkwqa1}{\textgreater{}}}

{clear a flag }

{}

\subparagraph{\texorpdfstring{{Parameters:}}{Parameters:}}\label{parameters-36}

\protect\hypertarget{t.383f7e723be615a1ab251cbb171dc44f8c48c9f4}{}{}\protect\hypertarget{t.36}{}{}

\begin{longtable}[]{@{}ll@{}}
\toprule
\begin{minipage}[t]{0.47\columnwidth}\raggedright\strut
{flag}{~}\strut
\end{minipage} & \begin{minipage}[t]{0.47\columnwidth}\raggedright\strut
{}\strut
\end{minipage}\tabularnewline
\bottomrule
\end{longtable}

\subparagraph{\texorpdfstring{{Returns:}}{Returns:}}\label{returns-40}

{\protect\hyperlink{h.3s49zyc}{THINKOS\_ENOSYS}}{\protect\hyperlink{h.3s49zyc}{~}}{if
call is not implemented,
}{\protect\hyperlink{h.2fk6b3p}{THINKOS\_OK}}{\protect\hyperlink{h.2fk6b3p}{~}}{otherwise.
}

\paragraph{\texorpdfstring{{int thinkos\_flag\_give (int
~flag)}}{int thinkos\_flag\_give (int ~flag)}}\label{int-thinkos_flag_give-int-flag}

{}

{\#include
\textless{}}{\protect\hyperlink{h.pkwqa1}{C:/devel/yard-ice/src/sdk/include/thinkos.h}}{\protect\hyperlink{h.pkwqa1}{\textgreater{}}}

{signal a flag }

{}

\subparagraph{\texorpdfstring{{Parameters:}}{Parameters:}}\label{parameters-37}

\protect\hypertarget{t.383f7e723be615a1ab251cbb171dc44f8c48c9f4}{}{}\protect\hypertarget{t.37}{}{}

\begin{longtable}[]{@{}ll@{}}
\toprule
\begin{minipage}[t]{0.47\columnwidth}\raggedright\strut
{flag}{~}\strut
\end{minipage} & \begin{minipage}[t]{0.47\columnwidth}\raggedright\strut
{}\strut
\end{minipage}\tabularnewline
\bottomrule
\end{longtable}

\subparagraph{\texorpdfstring{{Returns:}}{Returns:}}\label{returns-41}

{\protect\hyperlink{h.3s49zyc}{THINKOS\_ENOSYS}}{\protect\hyperlink{h.3s49zyc}{~}}{if
call is not implemented,
}{\protect\hyperlink{h.2fk6b3p}{THINKOS\_OK}}{\protect\hyperlink{h.2fk6b3p}{~}}{otherwise.
}

\paragraph{\texorpdfstring{{void thinkos\_flag\_give\_i (int
~flag)}}{void thinkos\_flag\_give\_i (int ~flag)}}\label{void-thinkos_flag_give_i-int-flag}

{}

{\#include
\textless{}}{\protect\hyperlink{h.pkwqa1}{C:/devel/yard-ice/src/sdk/include/thinkos.h}}{\protect\hyperlink{h.pkwqa1}{\textgreater{}}}

{signal a flag form inside an interrupt handler }

{}

\subparagraph{\texorpdfstring{{Parameters:}}{Parameters:}}\label{parameters-38}

\protect\hypertarget{t.383f7e723be615a1ab251cbb171dc44f8c48c9f4}{}{}\protect\hypertarget{t.38}{}{}

\begin{longtable}[]{@{}ll@{}}
\toprule
\begin{minipage}[t]{0.47\columnwidth}\raggedright\strut
{flag}{~}\strut
\end{minipage} & \begin{minipage}[t]{0.47\columnwidth}\raggedright\strut
{}\strut
\end{minipage}\tabularnewline
\bottomrule
\end{longtable}

\paragraph{\texorpdfstring{{int thinkos\_flag\_take (int
~flag)}}{int thinkos\_flag\_take (int ~flag)}}\label{int-thinkos_flag_take-int-flag}

{}

{\#include
\textless{}}{\protect\hyperlink{h.pkwqa1}{C:/devel/yard-ice/src/sdk/include/thinkos.h}}{\protect\hyperlink{h.pkwqa1}{\textgreater{}}}

{wait for flag to be signaled }

{}

\subparagraph{\texorpdfstring{{Parameters:}}{Parameters:}}\label{parameters-39}

\protect\hypertarget{t.383f7e723be615a1ab251cbb171dc44f8c48c9f4}{}{}\protect\hypertarget{t.39}{}{}

\begin{longtable}[]{@{}ll@{}}
\toprule
\begin{minipage}[t]{0.47\columnwidth}\raggedright\strut
{flag}{~}\strut
\end{minipage} & \begin{minipage}[t]{0.47\columnwidth}\raggedright\strut
{}\strut
\end{minipage}\tabularnewline
\bottomrule
\end{longtable}

\paragraph{\texorpdfstring{{int thinkos\_flag\_timedtake (int ~flag,
unsigned int
~ms)}}{int thinkos\_flag\_timedtake (int ~flag, unsigned int ~ms)}}\label{int-thinkos_flag_timedtake-int-flag-unsigned-int-ms}

{}

{\#include
\textless{}}{\protect\hyperlink{h.pkwqa1}{C:/devel/yard-ice/src/sdk/include/thinkos.h}}{\protect\hyperlink{h.pkwqa1}{\textgreater{}}}

{wait for flag or timeout }

{}

\subparagraph{\texorpdfstring{{Parameters:}}{Parameters:}}\label{parameters-40}

\protect\hypertarget{t.aed66b1bf3dc3447b75e535bc6b7a8c236003a69}{}{}\protect\hypertarget{t.40}{}{}

\begin{longtable}[]{@{}ll@{}}
\toprule
\begin{minipage}[t]{0.47\columnwidth}\raggedright\strut
{flag}{~}\strut
\end{minipage} & \begin{minipage}[t]{0.47\columnwidth}\raggedright\strut
{}\strut
\end{minipage}\tabularnewline
\begin{minipage}[t]{0.47\columnwidth}\raggedright\strut
{ms}{~}\strut
\end{minipage} & \begin{minipage}[t]{0.47\columnwidth}\raggedright\strut
{}\strut
\end{minipage}\tabularnewline
\bottomrule
\end{longtable}

\subparagraph{\texorpdfstring{{Returns:}}{Returns:}}\label{returns-42}

{\protect\hyperlink{h.3s49zyc}{THINKOS\_ENOSYS}}{\protect\hyperlink{h.3s49zyc}{~}}{if
call is not implemented,
}{\protect\hyperlink{h.2fk6b3p}{THINKOS\_OK}}{\protect\hyperlink{h.2fk6b3p}{~}}{otherwise.
}

\paragraph{\texorpdfstring{{}}{}}\label{section-6}

\subsection{\texorpdfstring{{Gates}}{Gates}}\label{gates}

\subsubsection{\texorpdfstring{{Functions}}{Functions}}\label{h.2zbgiuw}

{int
}{\protect\hyperlink{h.1egqt2p}{thinkos\_gate\_alloc}}{\protect\hyperlink{h.1egqt2p}{~}}{(void)}

{Alloc a gate synchronization object. }

{int
}{\protect\hyperlink{h.3ygebqi}{thinkos\_gate\_free}}{\protect\hyperlink{h.3ygebqi}{~}}{(int
gate)}

{Frees the gate synchronization object. }

{int
}{\protect\hyperlink{h.2dlolyb}{thinkos\_gate\_wait}}{\protect\hyperlink{h.2dlolyb}{~}}{(int
gate)}

{Wait for a gate to open. }

{int
}{\protect\hyperlink{h.sqyw64}{thinkos\_gate\_timedwait}}{\protect\hyperlink{h.sqyw64}{~}}{(int
gate, unsigned int ms)}

{Wait for a gate to open or a timeout. }

{int
}{\protect\hyperlink{h.3cqmetx}{thinkos\_gate\_open}}{\protect\hyperlink{h.3cqmetx}{~}}{(int
gate)}

{Open or signal the gate. }

{void
}{\protect\hyperlink{h.1rvwp1q}{thinkos\_gate\_open\_i}}{\protect\hyperlink{h.1rvwp1q}{~}}{(int
gate)}

{Open or signal the gate from inside an interrupt handler. }

{int
}{\protect\hyperlink{h.4bvk7pj}{thinkos\_gate\_close}}{\protect\hyperlink{h.4bvk7pj}{~}}{(int
gate)}

{Close the gate if the gate is }{OPEN}{~ or remove pending signaling it
the gate is @ LOCKED. }

{int
}{\protect\hyperlink{h.2r0uhxc}{thinkos\_gate\_exit}}{\protect\hyperlink{h.2r0uhxc}{~}}{(int
gate, unsigned int open)}

{Exit the gate, leaving the gate, optionally leaving it open or closed.
}

{}

\subsubsection{\texorpdfstring{{Detailed
Description}}{Detailed Description}}\label{detailed-description-7}

{Gates are syncronization objects which provide a convenient way of
creating mutual exclusion acess to code blocks signaled by interrupt
handlers\ldots{}}

{A gate have a lock flag and a signal flag. A gate can be in one of the
following states:}

{CLOSED:}{~ no threads crossed the gate yet.}

{LOCKED:}{~ a thread entered the gate, closed and locked it.}

{OPENED:}{~ no threads are waiting in the gate, the first thread to call
}{\protect\hyperlink{h.2dlolyb}{thinkos\_gate\_wait()}}{\protect\hyperlink{h.2dlolyb}{~}}{~will
cross the gate.}

{SIGNALED:}{~ a thread crossed the gate and locked it, but the gate
received a signal to open. When the thread exits the gate the gate will
stay open. }

{}

\subsubsection{\texorpdfstring{{Function
Documentation}}{Function Documentation}}\label{function-documentation-7}

\paragraph{\texorpdfstring{{int thinkos\_gate\_alloc (void
)}}{int thinkos\_gate\_alloc (void )}}\label{int-thinkos_gate_alloc-void}

{}

{\#include
\textless{}}{\protect\hyperlink{h.pkwqa1}{C:/devel/yard-ice/src/sdk/include/thinkos.h}}{\protect\hyperlink{h.pkwqa1}{\textgreater{}}}

{Alloc a gate synchronization object. }

{}

\subparagraph{\texorpdfstring{{Returns:}}{Returns:}}\label{returns-43}

{return a handler for a new gate object, or a negative value if an error
ocurred.}

{Errors:}

{\protect\hyperlink{h.3s49zyc}{THINKOS\_ENOSYS}}{\protect\hyperlink{h.3s49zyc}{~}}{if
the system call is not enabled.}

{\protect\hyperlink{h.meukdy}{THINKOS\_ENOMEM}}{\protect\hyperlink{h.meukdy}{~}}{no
gates left in the gate pool. }

\paragraph{\texorpdfstring{{int thinkos\_gate\_free (int
~gate)}}{int thinkos\_gate\_free (int ~gate)}}\label{int-thinkos_gate_free-int-gate}

{}

{\#include
\textless{}}{\protect\hyperlink{h.pkwqa1}{C:/devel/yard-ice/src/sdk/include/thinkos.h}}{\protect\hyperlink{h.pkwqa1}{\textgreater{}}}

{Frees the gate synchronization object. }

{}

\subparagraph{\texorpdfstring{{Parameters:}}{Parameters:}}\label{parameters-41}

\protect\hypertarget{t.732390f47e0bc9de7bbae841837625969c4e7db0}{}{}\protect\hypertarget{t.41}{}{}

\begin{longtable}[]{@{}ll@{}}
\toprule
\begin{minipage}[t]{0.47\columnwidth}\raggedright\strut
{gate}{~}\strut
\end{minipage} & \begin{minipage}[t]{0.47\columnwidth}\raggedright\strut
{handler for a gate object which must have been returned by a previous
call to
}{\protect\hyperlink{h.1egqt2p}{thinkos\_gate\_alloc()}}{\protect\hyperlink{h.1egqt2p}{~}}{.
}\strut
\end{minipage}\tabularnewline
\bottomrule
\end{longtable}

\subparagraph{\texorpdfstring{{Returns:}}{Returns:}}\label{returns-44}

{returns
}{\protect\hyperlink{h.2fk6b3p}{THINKOS\_OK}}{\protect\hyperlink{h.2fk6b3p}{~}}{on
sucess. On error a negative code value is returned. an error ocurred.}

{Errors:}

{\protect\hyperlink{h.1tuee74}{THINKOS\_EINVAL}}{\protect\hyperlink{h.1tuee74}{~}}{gate}{~
is not a valid gate handler.}

{\protect\hyperlink{h.3s49zyc}{THINKOS\_ENOSYS}}{\protect\hyperlink{h.3s49zyc}{~}}{not
implemented. }

\paragraph{\texorpdfstring{{int thinkos\_gate\_wait (int
~gate)}}{int thinkos\_gate\_wait (int ~gate)}}\label{int-thinkos_gate_wait-int-gate}

{}

{\#include
\textless{}}{\protect\hyperlink{h.pkwqa1}{C:/devel/yard-ice/src/sdk/include/thinkos.h}}{\protect\hyperlink{h.pkwqa1}{\textgreater{}}}

{Wait for a gate to open. }

{If the gate is open this function return imediatelly, otherwise it will
block the calling thread.}

\subparagraph{\texorpdfstring{{Parameters:}}{Parameters:}}\label{parameters-42}

\protect\hypertarget{t.63aac5addc95ec2809d3272ba50ab5f93508edb0}{}{}\protect\hypertarget{t.42}{}{}

\begin{longtable}[]{@{}ll@{}}
\toprule
\begin{minipage}[t]{0.47\columnwidth}\raggedright\strut
{gate}{~}\strut
\end{minipage} & \begin{minipage}[t]{0.47\columnwidth}\raggedright\strut
{The gate descriptor. }\strut
\end{minipage}\tabularnewline
\bottomrule
\end{longtable}

\subparagraph{\texorpdfstring{{Returns:}}{Returns:}}\label{returns-45}

{\protect\hyperlink{h.1tuee74}{THINKOS\_EINVAL}}{\protect\hyperlink{h.1tuee74}{~}}{if
}{gate}{~ is invalid,
}{\protect\hyperlink{h.2fk6b3p}{THINKOS\_OK}}{\protect\hyperlink{h.2fk6b3p}{~}}{otherwise.
}

\paragraph{\texorpdfstring{{int thinkos\_gate\_timedwait (int ~gate,
unsigned int
~ms)}}{int thinkos\_gate\_timedwait (int ~gate, unsigned int ~ms)}}\label{int-thinkos_gate_timedwait-int-gate-unsigned-int-ms}

{}

{\#include
\textless{}}{\protect\hyperlink{h.pkwqa1}{C:/devel/yard-ice/src/sdk/include/thinkos.h}}{\protect\hyperlink{h.pkwqa1}{\textgreater{}}}

{Wait for a gate to open or a timeout. }

{If the gate is open this function return imediatelly, otherwise it will
block the calling thread.}

\subparagraph{\texorpdfstring{{Parameters:}}{Parameters:}}\label{parameters-43}

\protect\hypertarget{t.acc072d2f3bb4c2c47837e0a954ea1c3c6ac8b37}{}{}\protect\hypertarget{t.43}{}{}

\begin{longtable}[]{@{}ll@{}}
\toprule
\begin{minipage}[t]{0.47\columnwidth}\raggedright\strut
{gate}{~}\strut
\end{minipage} & \begin{minipage}[t]{0.47\columnwidth}\raggedright\strut
{The gate descriptor. }\strut
\end{minipage}\tabularnewline
\begin{minipage}[t]{0.47\columnwidth}\raggedright\strut
{ms}{~}\strut
\end{minipage} & \begin{minipage}[t]{0.47\columnwidth}\raggedright\strut
{Timeout ins milliseconds. }\strut
\end{minipage}\tabularnewline
\bottomrule
\end{longtable}

\subparagraph{\texorpdfstring{{Returns:}}{Returns:}}\label{returns-46}

{}

{\protect\hyperlink{h.2fk6b3p}{THINKOS\_OK}}{\protect\hyperlink{h.2fk6b3p}{~}}{is
returned on sucess. On error a negative code value is returned.}

{\protect\hyperlink{h.1tuee74}{THINKOS\_EINVAL}}{\protect\hyperlink{h.1tuee74}{:}}{~}{gate}{~
is not a valid handler.}

{\protect\hyperlink{h.upglbi}{THINKOS\_ETIMEDOUT}}{\protect\hyperlink{h.upglbi}{:}}{~timer
expired before the }{gate}{~ opens.}

{\protect\hyperlink{h.3s49zyc}{THINKOS\_ENOSYS}}{\protect\hyperlink{h.3s49zyc}{:}}{~syscall
not implemented. }

\paragraph{\texorpdfstring{{int thinkos\_gate\_open (int
~gate)}}{int thinkos\_gate\_open (int ~gate)}}\label{int-thinkos_gate_open-int-gate}

{}

{\#include
\textless{}}{\protect\hyperlink{h.pkwqa1}{C:/devel/yard-ice/src/sdk/include/thinkos.h}}{\protect\hyperlink{h.pkwqa1}{\textgreater{}}}

{Open or signal the gate. }

{The resulting gate's state will depend on the current gate state and
whether there are threads waiting at the gate. There are four possible
scenarios \ldots{} :}

{1.~~~~~~~~the gate is open already, then this function does nothing.}

{2.~~~~~~~~the gate is closed and no threads are waiting it will open
the gate, allowing the next thread to call }{gate\_open()}{~ to enter
the gate.}

{3.~~~~~~~~the gate is closed and at least one thread is waiting it will
allow the thread to cross the gate, in this case the gate will be
locked.}

{4.~~~~~~~~a thread crossed the gate (gate state is }{LOCKED}{~), then
the gate will be signaled to open when the gate is unlocked.}

\subparagraph{\texorpdfstring{{Parameters:}}{Parameters:}}\label{parameters-44}

\protect\hypertarget{t.63aac5addc95ec2809d3272ba50ab5f93508edb0}{}{}\protect\hypertarget{t.44}{}{}

\begin{longtable}[]{@{}ll@{}}
\toprule
\begin{minipage}[t]{0.47\columnwidth}\raggedright\strut
{gate}{~}\strut
\end{minipage} & \begin{minipage}[t]{0.47\columnwidth}\raggedright\strut
{The gate descriptor. }\strut
\end{minipage}\tabularnewline
\bottomrule
\end{longtable}

\subparagraph{\texorpdfstring{{Returns:}}{Returns:}}\label{returns-47}

{\protect\hyperlink{h.1tuee74}{THINKOS\_EINVAL}}{\protect\hyperlink{h.1tuee74}{~}}{if
}{gate}{~ is invalid,
}{\protect\hyperlink{h.2fk6b3p}{THINKOS\_OK}}{\protect\hyperlink{h.2fk6b3p}{~}}{otherwise.
}

\paragraph{\texorpdfstring{{void thinkos\_gate\_open\_i (int
~gate)}}{void thinkos\_gate\_open\_i (int ~gate)}}\label{void-thinkos_gate_open_i-int-gate}

{}

{\#include
\textless{}}{\protect\hyperlink{h.pkwqa1}{C:/devel/yard-ice/src/sdk/include/thinkos.h}}{\protect\hyperlink{h.pkwqa1}{\textgreater{}}}

{Open or signal the gate from inside an interrupt handler. }

{This call is similar to the
}{\protect\hyperlink{h.3cqmetx}{thinkos\_gate\_open()}}{\protect\hyperlink{h.3cqmetx}{~}}{~except
that it is safe to ba called from inside an interrupt handler.}

\subparagraph{\texorpdfstring{{Parameters:}}{Parameters:}}\label{parameters-45}

\protect\hypertarget{t.5f99e327007a50c0bedc6d95e651021a3e43c836}{}{}\protect\hypertarget{t.45}{}{}

\begin{longtable}[]{@{}ll@{}}
\toprule
\begin{minipage}[t]{0.47\columnwidth}\raggedright\strut
{gate}{~}\strut
\end{minipage} & \begin{minipage}[t]{0.47\columnwidth}\raggedright\strut
{The gate descriptor.}\strut
\end{minipage}\tabularnewline
\bottomrule
\end{longtable}

{Warning:}{~ no argument validation is performed. }

\paragraph{\texorpdfstring{{int thinkos\_gate\_close (int
~gate)}}{int thinkos\_gate\_close (int ~gate)}}\label{int-thinkos_gate_close-int-gate}

{}

{\#include
\textless{}}{\protect\hyperlink{h.pkwqa1}{C:/devel/yard-ice/src/sdk/include/thinkos.h}}{\protect\hyperlink{h.pkwqa1}{\textgreater{}}}

{Close the gate if the gate is }{OPEN}{~ or remove pending signaling it
the gate is @ LOCKED. }

{}

\subparagraph{\texorpdfstring{{Parameters:}}{Parameters:}}\label{parameters-46}

\protect\hypertarget{t.63aac5addc95ec2809d3272ba50ab5f93508edb0}{}{}\protect\hypertarget{t.46}{}{}

\begin{longtable}[]{@{}ll@{}}
\toprule
\begin{minipage}[t]{0.47\columnwidth}\raggedright\strut
{gate}{~}\strut
\end{minipage} & \begin{minipage}[t]{0.47\columnwidth}\raggedright\strut
{The gate descriptor. }\strut
\end{minipage}\tabularnewline
\bottomrule
\end{longtable}

\subparagraph{\texorpdfstring{{Returns:}}{Returns:}}\label{returns-48}

{\protect\hyperlink{h.1tuee74}{THINKOS\_EINVAL}}{\protect\hyperlink{h.1tuee74}{~}}{if
}{gate}{~ is invalid,
}{\protect\hyperlink{h.2fk6b3p}{THINKOS\_OK}}{\protect\hyperlink{h.2fk6b3p}{~}}{otherwise.
}

\paragraph{\texorpdfstring{{int thinkos\_gate\_exit (int ~gate, unsigned
int
~open)}}{int thinkos\_gate\_exit (int ~gate, unsigned int ~open)}}\label{int-thinkos_gate_exit-int-gate-unsigned-int-open}

{}

{\#include
\textless{}}{\protect\hyperlink{h.pkwqa1}{C:/devel/yard-ice/src/sdk/include/thinkos.h}}{\protect\hyperlink{h.pkwqa1}{\textgreater{}}}

{Exit the gate, leaving the gate, optionally leaving it open or closed.
}

{}

\subparagraph{\texorpdfstring{{Parameters:}}{Parameters:}}\label{parameters-47}

\protect\hypertarget{t.ed499f740a0883130710bbb983d50c5a5eccdc42}{}{}\protect\hypertarget{t.47}{}{}

\begin{longtable}[]{@{}ll@{}}
\toprule
\begin{minipage}[t]{0.47\columnwidth}\raggedright\strut
{gate}{~}\strut
\end{minipage} & \begin{minipage}[t]{0.47\columnwidth}\raggedright\strut
{The gate descriptor. }\strut
\end{minipage}\tabularnewline
\begin{minipage}[t]{0.47\columnwidth}\raggedright\strut
{open}{~}\strut
\end{minipage} & \begin{minipage}[t]{0.47\columnwidth}\raggedright\strut
{Indicate the state of the gate on exit.}\strut
\end{minipage}\tabularnewline
\bottomrule
\end{longtable}

{open}{~ \textgreater{} 0, the gate will be left open, allowing for
another thread to enter the gate.}

{open}{~ == 0, the gate will stay closed if not signaled, in wich case
it will open accordingly. }

\protect\hypertarget{t.1aebe24dda1b56316fca09844ec7bf63e4290b0d}{}{}\protect\hypertarget{t.48}{}{}

\begin{longtable}[]{@{}l@{}}
\toprule
\begin{minipage}[t]{0.97\columnwidth}\raggedright\strut
{}\strut
\end{minipage}\tabularnewline
\bottomrule
\end{longtable}

\subparagraph{\texorpdfstring{{Returns:}}{Returns:}}\label{returns-49}

{\protect\hyperlink{h.1tuee74}{THINKOS\_EINVAL}}{\protect\hyperlink{h.1tuee74}{~}}{if
}{gate}{~ is invalid,
}{\protect\hyperlink{h.2fk6b3p}{THINKOS\_OK}}{\protect\hyperlink{h.2fk6b3p}{~}}{otherwise.
}

\paragraph{\texorpdfstring{{}}{}}\label{section-7}

\subsection{\texorpdfstring{{Interrupt
Requests}}{Interrupt Requests}}\label{interrupt-requests}

\subsubsection{\texorpdfstring{{Functions}}{Functions}}\label{h.1664s55}

{int
}{\protect\hyperlink{h.3q5sasy}{thinkos\_irq\_wait}}{\protect\hyperlink{h.3q5sasy}{~}}{(int
irq)}

{wait for interrupt }

{int
}{\protect\hyperlink{h.25b2l0r}{thinkos\_irq\_register}}{\protect\hyperlink{h.25b2l0r}{~}}{(int
irq, int pri, void(*isr)(void))}

{reqister an interrupt handler }

{}

\subsubsection{\texorpdfstring{{Detailed
Description}}{Detailed Description}}\label{detailed-description-8}

{}

{}

\subsubsection{\texorpdfstring{{Function
Documentation}}{Function Documentation}}\label{function-documentation-8}

\paragraph{\texorpdfstring{{int thinkos\_irq\_wait (int
~irq)}}{int thinkos\_irq\_wait (int ~irq)}}\label{int-thinkos_irq_wait-int-irq}

{}

{\#include
\textless{}}{\protect\hyperlink{h.pkwqa1}{C:/devel/yard-ice/src/sdk/include/thinkos.h}}{\protect\hyperlink{h.pkwqa1}{\textgreater{}}}

{wait for interrupt }

{}

\subparagraph{\texorpdfstring{{Parameters:}}{Parameters:}}\label{parameters-48}

\protect\hypertarget{t.c2957ce4f831c006182672b15100a2a6a33851a6}{}{}\protect\hypertarget{t.49}{}{}

\begin{longtable}[]{@{}ll@{}}
\toprule
\begin{minipage}[t]{0.47\columnwidth}\raggedright\strut
{irq}{~}\strut
\end{minipage} & \begin{minipage}[t]{0.47\columnwidth}\raggedright\strut
{}\strut
\end{minipage}\tabularnewline
\bottomrule
\end{longtable}

\subparagraph{\texorpdfstring{{Returns:}}{Returns:}}\label{returns-50}

{\protect\hyperlink{h.3s49zyc}{THINKOS\_ENOSYS}}{\protect\hyperlink{h.3s49zyc}{~}}{if
call is not implemented,
}{\protect\hyperlink{h.2fk6b3p}{THINKOS\_OK}}{\protect\hyperlink{h.2fk6b3p}{~}}{otherwise.
}

\paragraph{\texorpdfstring{{int thinkos\_irq\_register (int ~irq, int
~pri, void(*)(void)
~isr)}}{int thinkos\_irq\_register (int ~irq, int ~pri, void(*)(void) ~isr)}}\label{int-thinkos_irq_register-int-irq-int-pri-voidvoid-isr}

{}

{\#include
\textless{}}{\protect\hyperlink{h.pkwqa1}{C:/devel/yard-ice/src/sdk/include/thinkos.h}}{\protect\hyperlink{h.pkwqa1}{\textgreater{}}}

{reqister an interrupt handler }

{}

\subparagraph{\texorpdfstring{{Parameters:}}{Parameters:}}\label{parameters-49}

\protect\hypertarget{t.bfce7faafba4041e0a1d16655da9fc068af23510}{}{}\protect\hypertarget{t.50}{}{}

\begin{longtable}[]{@{}ll@{}}
\toprule
\begin{minipage}[t]{0.47\columnwidth}\raggedright\strut
{irq}{~}\strut
\end{minipage} & \begin{minipage}[t]{0.47\columnwidth}\raggedright\strut
{}\strut
\end{minipage}\tabularnewline
\begin{minipage}[t]{0.47\columnwidth}\raggedright\strut
{pri}{~}\strut
\end{minipage} & \begin{minipage}[t]{0.47\columnwidth}\raggedright\strut
{}\strut
\end{minipage}\tabularnewline
\begin{minipage}[t]{0.47\columnwidth}\raggedright\strut
{isr}{~}\strut
\end{minipage} & \begin{minipage}[t]{0.47\columnwidth}\raggedright\strut
{}\strut
\end{minipage}\tabularnewline
\bottomrule
\end{longtable}

\subparagraph{\texorpdfstring{{Returns:}}{Returns:}}\label{returns-51}

{\protect\hyperlink{h.3s49zyc}{THINKOS\_ENOSYS}}{\protect\hyperlink{h.3s49zyc}{~}}{if
call is not implemented,
}{\protect\hyperlink{h.2fk6b3p}{THINKOS\_OK}}{\protect\hyperlink{h.2fk6b3p}{~}}{otherwise.
}

\paragraph{\texorpdfstring{{}}{}}\label{section-8}

\subsection{\texorpdfstring{{OS Monitor and
Cotrol.}}{OS Monitor and Cotrol.}}\label{os-monitor-and-cotrol.}

\subsubsection{\texorpdfstring{{Functions}}{Functions}}\label{h.kgcv8k}

{int
}{\protect\hyperlink{h.34g0dwd}{thinkos\_clocks}}{\protect\hyperlink{h.34g0dwd}{~}}{(uint32\_t
*clk{[}{]})}

{get hardware clocks frequencies }

{int
}{\protect\hyperlink{h.1jlao46}{thinkos\_udelay\_factor}}{\protect\hyperlink{h.1jlao46}{~}}{(int32\_t
*factor)}

{get udelay calibration factor }

{void
}{\protect\hyperlink{h.43ky6rz}{\_\_attribute\_\_}}{\protect\hyperlink{h.43ky6rz}{~}}{((noreturn))
thinkos\_abort(void)}

{abort the operating system. }

{}

\subsubsection{\texorpdfstring{{Detailed
Description}}{Detailed Description}}\label{detailed-description-9}

{These calls return information about the ThinkOS configuratiion and
operational parameters. }

{}

\subsubsection{\texorpdfstring{{Function
Documentation}}{Function Documentation}}\label{function-documentation-9}

\paragraph{\texorpdfstring{{int thinkos\_clocks (uint32\_t *
~clk{[}{]})}}{int thinkos\_clocks (uint32\_t * ~clk{[}{]})}}\label{int-thinkos_clocks-uint32_t-clk}

{}

{\#include
\textless{}}{\protect\hyperlink{h.pkwqa1}{C:/devel/yard-ice/src/sdk/include/thinkos.h}}{\protect\hyperlink{h.pkwqa1}{\textgreater{}}}

{get hardware clocks frequencies }

{}

\subparagraph{\texorpdfstring{{Parameters:}}{Parameters:}}\label{parameters-50}

\protect\hypertarget{t.866b9b535e5516e27514680fa2775da750ad842a}{}{}\protect\hypertarget{t.51}{}{}

\begin{longtable}[]{@{}ll@{}}
\toprule
\begin{minipage}[t]{0.47\columnwidth}\raggedright\strut
{clk{[}{]}}{~}\strut
\end{minipage} & \begin{minipage}[t]{0.47\columnwidth}\raggedright\strut
{pointer to an array to receive the clocks frequency list. }\strut
\end{minipage}\tabularnewline
\bottomrule
\end{longtable}

\subparagraph{\texorpdfstring{{Returns:}}{Returns:}}\label{returns-52}

{\protect\hyperlink{h.3s49zyc}{THINKOS\_ENOSYS}}{\protect\hyperlink{h.3s49zyc}{~}}{if
call is not implemented,
}{\protect\hyperlink{h.2fk6b3p}{THINKOS\_OK}}{\protect\hyperlink{h.2fk6b3p}{~}}{otherwise.
}

\paragraph{\texorpdfstring{{int thinkos\_udelay\_factor (int32\_t *
~factor)}}{int thinkos\_udelay\_factor (int32\_t * ~factor)}}\label{int-thinkos_udelay_factor-int32_t-factor}

{}

{\#include
\textless{}}{\protect\hyperlink{h.pkwqa1}{C:/devel/yard-ice/src/sdk/include/thinkos.h}}{\protect\hyperlink{h.pkwqa1}{\textgreater{}}}

{get udelay calibration factor }

{}

\subparagraph{\texorpdfstring{{Parameters:}}{Parameters:}}\label{parameters-51}

\protect\hypertarget{t.ba04b86e0085018617d8457d329191a496a93af3}{}{}\protect\hypertarget{t.52}{}{}

\begin{longtable}[]{@{}ll@{}}
\toprule
\begin{minipage}[t]{0.47\columnwidth}\raggedright\strut
{factor}{~}\strut
\end{minipage} & \begin{minipage}[t]{0.47\columnwidth}\raggedright\strut
{pointer to an integer. }\strut
\end{minipage}\tabularnewline
\bottomrule
\end{longtable}

\subparagraph{\texorpdfstring{{Returns:}}{Returns:}}\label{returns-53}

{\protect\hyperlink{h.3s49zyc}{THINKOS\_ENOSYS}}{\protect\hyperlink{h.3s49zyc}{~}}{if
call is not implemented,
}{\protect\hyperlink{h.2fk6b3p}{THINKOS\_OK}}{\protect\hyperlink{h.2fk6b3p}{~}}{otherwise.
}

\paragraph{\texorpdfstring{{void \_\_attribute\_\_ ((noreturn)
)}}{void \_\_attribute\_\_ ((noreturn) )}}\label{void-__attribute__-noreturn}

{}

{\#include
\textless{}}{\protect\hyperlink{h.pkwqa1}{C:/devel/yard-ice/src/sdk/include/thinkos.h}}{\protect\hyperlink{h.pkwqa1}{\textgreater{}}}

{abort the operating system. }

{}

\subparagraph{\texorpdfstring{{Returns:}}{Returns:}}\label{returns-54}

{This function does not return. }

{}

\section{\texorpdfstring{{Data Structure
Documentation}}{Data Structure Documentation}}\label{data-structure-documentation}

\subsection{\texorpdfstring{{thinkos\_thread\_inf Struct
Reference}}{thinkos\_thread\_inf Struct Reference}}\label{thinkos_thread_inf-struct-reference}

{}

{\#include \textless{}thinkos.h\textgreater{}}

\subsubsection{\texorpdfstring{{Data
Fields}}{Data Fields}}\label{data-fields}

{void * }{\protect\hyperlink{h.xvir7l}{stack\_ptr}}

{union \{}

{~ uint32\_t }{\protect\hyperlink{h.3hv69ve}{opt}}

{~ struct \{}

{~ ~ uint16\_t }{\protect\hyperlink{h.1x0gk37}{stack\_size}}

{~ ~ uint8\_t }{\protect\hyperlink{h.4h042r0}{priority}}

{~ ~ uint8\_t
}{\protect\hyperlink{h.2w5ecyt}{thread\_id}}{\protect\hyperlink{h.2w5ecyt}{:}}{~7}

{~ ~ uint8\_t
}{\protect\hyperlink{h.1baon6m}{paused}}{\protect\hyperlink{h.1baon6m}{:}}{~1}

{~ \} }

{\}; }

{char
}{\protect\hyperlink{h.2afmg28}{tag}}{\protect\hyperlink{h.2afmg28}{~}}{{[}8{]}}

{}

\subsubsection{\texorpdfstring{{Field
Documentation}}{Field Documentation}}\label{field-documentation}

\paragraph{\texorpdfstring{{void*
thinkos\_thread\_inf::stack\_ptr}}{void* thinkos\_thread\_inf::stack\_ptr}}\label{void-thinkos_thread_infstack_ptr}

\hypertarget{h.xvir7l}{\paragraph{\texorpdfstring{{uint32\_t
thinkos\_thread\_inf::opt}}{uint32\_t thinkos\_thread\_inf::opt}}\label{h.xvir7l}}

\hypertarget{h.3hv69ve}{\paragraph{\texorpdfstring{{uint16\_t
thinkos\_thread\_inf::stack\_size}}{uint16\_t thinkos\_thread\_inf::stack\_size}}\label{h.3hv69ve}}

\hypertarget{h.1x0gk37}{\paragraph{\texorpdfstring{{uint8\_t
thinkos\_thread\_inf::priority}}{uint8\_t thinkos\_thread\_inf::priority}}\label{h.1x0gk37}}

\hypertarget{h.4h042r0}{\paragraph{\texorpdfstring{{uint8\_t
thinkos\_thread\_inf::thread\_id}}{uint8\_t thinkos\_thread\_inf::thread\_id}}\label{h.4h042r0}}

\hypertarget{h.2w5ecyt}{\paragraph{\texorpdfstring{{uint8\_t
thinkos\_thread\_inf::paused}}{uint8\_t thinkos\_thread\_inf::paused}}\label{h.2w5ecyt}}

\hypertarget{h.1baon6m}{\paragraph{\texorpdfstring{{union \{ \ldots{} \}
}}{union \{ \ldots{} \} }}\label{h.1baon6m}}

\paragraph{\texorpdfstring{{char
thinkos\_thread\_inf::tag{[}8{]}}}{char thinkos\_thread\_inf::tag{[}8{]}}}\label{h.3vac5uf}

{}

\section{\texorpdfstring{{File
Documentation}}{File Documentation}}\label{file-documentation}

\subsection{\texorpdfstring{{C:/devel/yard-ice/src/sdk/include/thinkos.h}}{C:/devel/yard-ice/src/sdk/include/thinkos.h}}\label{cdevelyard-icesrcsdkincludethinkos.h}

{}

{ThinkOS API. }

{\#include \textless{}stdint.h\textgreater{}}

{\#include \textless{}thinkos\_svc.h\textgreater{}}

\subsubsection{\texorpdfstring{{Data
Structures}}{Data Structures}}\label{data-structures-1}

{struct }{\protect\hyperlink{h.2iq8gzs}{thinkos\_thread\_inf}}

\subsubsection{\texorpdfstring{{Macros}}{Macros}}\label{macros}

{\#define
}{\protect\hyperlink{h.39kk8xu}{IRQ\_PRIORITY\_HIGHEST}}{\protect\hyperlink{h.39kk8xu}{~}}{~(1
\textless{}\textless{} 5)}

{\#define
}{\protect\hyperlink{h.1opuj5n}{IRQ\_PRIORITY\_VERY\_HIGH}}{\protect\hyperlink{h.1opuj5n}{~}}{~(2
\textless{}\textless{} 5)}

{\#define
}{\protect\hyperlink{h.48pi1tg}{IRQ\_PRIORITY\_HIGH}}{\protect\hyperlink{h.48pi1tg}{~}}{~(3
\textless{}\textless{} 5)}

{\#define
}{\protect\hyperlink{h.2nusc19}{IRQ\_PRIORITY\_REGULAR}}{\protect\hyperlink{h.2nusc19}{~}}{~(4
\textless{}\textless{} 5)}

{\#define
}{\protect\hyperlink{h.1302m92}{IRQ\_PRIORITY\_LOW}}{\protect\hyperlink{h.1302m92}{~}}{~(5
\textless{}\textless{} 5)}

{\#define
}{\protect\hyperlink{h.3mzq4wv}{IRQ\_PRIORITY\_VERY\_LOW}}{\protect\hyperlink{h.3mzq4wv}{~}}{~(6
\textless{}\textless{} 5)}

{\#define
}{\protect\hyperlink{h.2250f4o}{THINKOS\_OPT\_PRIORITY}}{\protect\hyperlink{h.2250f4o}{(}}{VAL)~
(((VAL) \& 0xff) \textless{}\textless{} 16)}

{\#define
}{\protect\hyperlink{h.haapch}{THINKOS\_OPT\_ID}}{\protect\hyperlink{h.haapch}{(}}{VAL)~
(((VAL) \& 0x07f) \textless{}\textless{} 24)}

{\#define
}{\protect\hyperlink{h.319y80a}{THINKOS\_OPT\_PAUSED}}{\protect\hyperlink{h.319y80a}{~}}{~(1
\textless{}\textless{} 31) /* don't run at startup */}

{\#define
}{\protect\hyperlink{h.1gf8i83}{THINKOS\_OPT\_STACK\_SIZE}}{\protect\hyperlink{h.1gf8i83}{(}}{VAL)~
((VAL) \& 0xffff)}

\subsubsection{\texorpdfstring{{Enumerations}}{Enumerations}}\label{enumerations}

{enum }{\protect\hyperlink{h.40ew0vw}{thinkos\_err}}{\{ }

{\protect\hyperlink{h.2fk6b3p}{THINKOS\_OK}}{= 0, }

{\protect\hyperlink{h.upglbi}{THINKOS\_ETIMEDOUT}}{= -1, }

{\protect\hyperlink{h.3ep43zb}{THINKOS\_EINTR}}{= -2, }

{\protect\hyperlink{h.1tuee74}{THINKOS\_EINVAL}}{= -3, }

{\protect\hyperlink{h.4du1wux}{THINKOS\_EAGAIN}}{= -4, }

{\protect\hyperlink{h.2szc72q}{THINKOS\_EDEADLK}}{= -5, }

{\protect\hyperlink{h.184mhaj}{THINKOS\_EPERM}}{= -6, }

{\protect\hyperlink{h.3s49zyc}{THINKOS\_ENOSYS}}{= -7, }

{\protect\hyperlink{h.279ka65}{THINKOS\_EFAULT}}{= -8, }

{\protect\hyperlink{h.meukdy}{THINKOS\_ENOMEM}}{= -9 }

{\}}

{enum
}{\protect\hyperlink{h.36ei31r}{thinkos\_obj\_kind}}{\protect\hyperlink{h.36ei31r}{~}}{\{
}

{\protect\hyperlink{h.1ljsd9k}{THINKOS\_OBJ\_READY}}{\protect\hyperlink{h.1ljsd9k}{~}}{=
0, }

{\protect\hyperlink{h.45jfvxd}{THINKOS\_OBJ\_TMSHARE}}{\protect\hyperlink{h.45jfvxd}{,}}{~}

{\protect\hyperlink{h.2koq656}{THINKOS\_OBJ\_CLOCK}}{\protect\hyperlink{h.2koq656}{,}}{~}

{\protect\hyperlink{h.zu0gcz}{THINKOS\_OBJ\_MUTEX}}{\protect\hyperlink{h.zu0gcz}{,}}{~}

{\protect\hyperlink{h.3jtnz0s}{THINKOS\_OBJ\_COND}}{\protect\hyperlink{h.3jtnz0s}{,}}{~}

{\protect\hyperlink{h.1yyy98l}{THINKOS\_OBJ\_SEMAPHORE}}{\protect\hyperlink{h.1yyy98l}{,}}{~}

{\protect\hyperlink{h.4iylrwe}{THINKOS\_OBJ\_EVENT}}{\protect\hyperlink{h.4iylrwe}{,}}{~}

{\protect\hyperlink{h.2y3w247}{THINKOS\_OBJ\_FLAG}}{\protect\hyperlink{h.2y3w247}{,}}{~}

{\protect\hyperlink{h.1d96cc0}{THINKOS\_OBJ\_GATE}}{\protect\hyperlink{h.1d96cc0}{,}}{~}

{\protect\hyperlink{h.3x8tuzt}{THINKOS\_OBJ\_JOIN}}{\protect\hyperlink{h.3x8tuzt}{,}}{~}

{\protect\hyperlink{h.2ce457m}{THINKOS\_OBJ\_CONREAD}}{\protect\hyperlink{h.2ce457m}{,}}{~}

{\protect\hyperlink{h.rjefff}{THINKOS\_OBJ\_CONWRITE}}{\protect\hyperlink{h.rjefff}{,}}{~}

{\protect\hyperlink{h.3bj1y38}{THINKOS\_OBJ\_PAUSED}}{\protect\hyperlink{h.3bj1y38}{,}}{~}

{\protect\hyperlink{h.1qoc8b1}{THINKOS\_OBJ\_CANCELED}}{\protect\hyperlink{h.1qoc8b1}{,}}{~}

{\protect\hyperlink{h.4anzqyu}{THINKOS\_OBJ\_FAULT}}{\protect\hyperlink{h.4anzqyu}{,}}{~}

{\protect\hyperlink{h.2pta16n}{THINKOS\_OBJ\_INVALID}}{\protect\hyperlink{h.2pta16n}{~}}

{\}}

\subsubsection{\texorpdfstring{{Functions}}{Functions}}\label{functions}

{int
}{\protect\hyperlink{h.14ykbeg}{thinkos\_init}}{\protect\hyperlink{h.14ykbeg}{~}}{(unsigned
int opt)}

{Initializes the }{ThinkOS}{~ library. }

{int
}{\protect\hyperlink{h.1fob9te}{thinkos\_thread\_create}}{(int(*task\_ptr)(void
*), void *task\_arg, void *stack\_ptr, unsigned int opt)}

{create a new thread }

{int
}{\protect\hyperlink{h.3znysh7}{thinkos\_thread\_create\_inf}}{(int(*task\_ptr)(void
*), void *task\_arg, const struct
}{\protect\hyperlink{h.2iq8gzs}{thinkos\_thread\_inf}}{\protect\hyperlink{h.2iq8gzs}{~}}{*inf)}

{create a new thread with extend information }

{int const
}{\protect\hyperlink{h.2et92p0}{thinkos\_thread\_self}}{(void)}

{obtaind a handler for the calling thread }

{int
}{\protect\hyperlink{h.tyjcwt}{thinkos\_cancel}}{\protect\hyperlink{h.tyjcwt}{~}}{(unsigned
int thread\_id, int code)}

{request a thread to terminate }

{int
}{\protect\hyperlink{h.3dy6vkm}{thinkos\_exit}}{\protect\hyperlink{h.3dy6vkm}{~}}{(int
code)}

{cause the thread to terminate }

{int
}{\protect\hyperlink{h.1t3h5sf}{thinkos\_join}}{\protect\hyperlink{h.1t3h5sf}{~}}{(unsigned
int thread\_id)}

{join with a terminated thread }

{int
}{\protect\hyperlink{h.4d34og8}{thinkos\_pause}}{\protect\hyperlink{h.4d34og8}{~}}{(unsigned
int thread\_id)}

{pause the thread execution }

{int
}{\protect\hyperlink{h.2s8eyo1}{thinkos\_resume}}{\protect\hyperlink{h.2s8eyo1}{~}}{(unsigned
int thread\_id)}

{resume a thread operation }

{void
}{\protect\hyperlink{h.17dp8vu}{thinkos\_yield}}{\protect\hyperlink{h.17dp8vu}{~}}{(void)}

{causes the calling thread to relinquish the CPU. The thread is moved to
the end of the queue for its static priority and a new thread gets to
run. }

{int
}{\protect\hyperlink{h.26in1rg}{thinkos\_sleep}}{\protect\hyperlink{h.26in1rg}{~}}{(unsigned
int ms)}

{pause for a specified amount of time }

{uint32\_t
}{\protect\hyperlink{h.lnxbz9}{thinkos\_clock}}{\protect\hyperlink{h.lnxbz9}{~}}{(void)}

{retrieve the ThikOS clock time }

{int
}{\protect\hyperlink{h.35nkun2}{thinkos\_alarm}}{\protect\hyperlink{h.35nkun2}{~}}{(uint32\_t
clock)}

{pause until a specified alarm time }

{int
}{\protect\hyperlink{h.44sinio}{thinkos\_mutex\_alloc}}{\protect\hyperlink{h.44sinio}{~}}{(void)}

{alloc a mutex }

{int
}{\protect\hyperlink{h.2jxsxqh}{thinkos\_mutex\_free}}{\protect\hyperlink{h.2jxsxqh}{~}}{(int
mutex)}

{release a mutex }

{int
}{\protect\hyperlink{h.z337ya}{thinkos\_mutex\_lock}}{\protect\hyperlink{h.z337ya}{~}}{(int
mutex)}

{lock a mutex, blocking if already locked. }

{int
}{\protect\hyperlink{h.3j2qqm3}{thinkos\_mutex\_trylock}}{\protect\hyperlink{h.3j2qqm3}{~}}{(int
mutex)}

{try to lock a mutex, non blocking. }

{int
}{\protect\hyperlink{h.1y810tw}{thinkos\_mutex\_timedlock}}{\protect\hyperlink{h.1y810tw}{~}}{(int
mutex, unsigned int ms)}

{lock a mutex, blocking for limited time. }

{int
}{\protect\hyperlink{h.4i7ojhp}{thinkos\_mutex\_unlock}}{\protect\hyperlink{h.4i7ojhp}{~}}{(int
mutex)}

{unlock the mutex }

{int
}{\protect\hyperlink{h.1ci93xb}{thinkos\_cond\_alloc}}{\protect\hyperlink{h.1ci93xb}{~}}{(void)}

{alloc a conditional variable }

{int
}{\protect\hyperlink{h.3whwml4}{thinkos\_cond\_free}}{\protect\hyperlink{h.3whwml4}{~}}{(int
cond)}

{release a conditional variable }

{int
}{\protect\hyperlink{h.2bn6wsx}{thinkos\_cond\_wait}}{\protect\hyperlink{h.2bn6wsx}{~}}{(int
cond, int mutex)}

{wait for a conditional variable }

{int
}{\protect\hyperlink{h.qsh70q}{thinkos\_cond\_timedwait}}{\protect\hyperlink{h.qsh70q}{~}}{(int
cond, int mutex, unsigned int ms)}

{wait for a conditional variable or timeout }

{int
}{\protect\hyperlink{h.3as4poj}{thinkos\_cond\_signal}}{\protect\hyperlink{h.3as4poj}{~}}{(int
cond)}

{signal a conditional variable wake a single thread }

{int
}{\protect\hyperlink{h.1pxezwc}{thinkos\_cond\_broadcast}}{\protect\hyperlink{h.1pxezwc}{~}}{(int
cond)}

{signal a conditional variable wake all waiting threads }

{int
}{\protect\hyperlink{h.2p2csry}{thinkos\_sem\_alloc}}{\protect\hyperlink{h.2p2csry}{~}}{(unsigned
int val)}

{alloc a semaphore }

{int
}{\protect\hyperlink{h.147n2zr}{thinkos\_sem\_free}}{\protect\hyperlink{h.147n2zr}{~}}{(int
sem)}

{release a semaphore }

{int
}{\protect\hyperlink{h.3o7alnk}{thinkos\_sem\_init}}{\protect\hyperlink{h.3o7alnk}{~}}{(int
sem, unsigned int val)}

{set the initial value of a semaphore }

{int
}{\protect\hyperlink{h.23ckvvd}{thinkos\_sem\_wait}}{\protect\hyperlink{h.23ckvvd}{~}}{(int
sem)}

{wait for a semaphore to be signaled }

{int
}{\protect\hyperlink{h.ihv636}{thinkos\_sem\_timedwait}}{\protect\hyperlink{h.ihv636}{~}}{(int
sem, unsigned int ms)}

{wait for a semaphore to be signaled or timeout }

{int
}{\protect\hyperlink{h.32hioqz}{thinkos\_sem\_post}}{\protect\hyperlink{h.32hioqz}{~}}{(int
sem)}

{signal a semaphore }

{void
}{\protect\hyperlink{h.1hmsyys}{thinkos\_sem\_post\_i}}{\protect\hyperlink{h.1hmsyys}{~}}{(int
sem)}

{signal a semaphore inside an interrupt handler }

{int
}{\protect\hyperlink{h.2grqrue}{thinkos\_ev\_alloc}}{\protect\hyperlink{h.2grqrue}{~}}{(void)}

{alloc an event set }

{int
}{\protect\hyperlink{h.vx1227}{thinkos\_ev\_free}}{\protect\hyperlink{h.vx1227}{~}}{(int
set)}

{release an event set }

{int
}{\protect\hyperlink{h.3fwokq0}{thinkos\_ev\_wait}}{\protect\hyperlink{h.3fwokq0}{~}}{(int
set)}

{wait for an event }

{int
}{\protect\hyperlink{h.1v1yuxt}{thinkos\_ev\_timedwait}}{\protect\hyperlink{h.1v1yuxt}{~}}{(int
set, unsigned int ms)}

{wait for an event or timeout }

{int
}{\protect\hyperlink{h.4f1mdlm}{thinkos\_ev\_raise}}{\protect\hyperlink{h.4f1mdlm}{~}}{(int
set, int ev)}

{signal an event }

{void
}{\protect\hyperlink{h.2u6wntf}{thinkos\_ev\_raise\_i}}{\protect\hyperlink{h.2u6wntf}{~}}{(int
set, int ev)}

{signal an event from inside an interrupt handler }

{int
}{\protect\hyperlink{h.19c6y18}{thinkos\_ev\_mask}}{\protect\hyperlink{h.19c6y18}{~}}{(int
set, int ev, int val)}

{mask/unmask an event }

{int
}{\protect\hyperlink{h.3tbugp1}{thinkos\_ev\_clear}}{\protect\hyperlink{h.3tbugp1}{~}}{(int
set, int ev)}

{clear an event }

{int
}{\protect\hyperlink{h.nmf14n}{thinkos\_flag\_alloc}}{\protect\hyperlink{h.nmf14n}{~}}{(void)}

{alloc a flag }

{int
}{\protect\hyperlink{h.37m2jsg}{thinkos\_flag\_free}}{\protect\hyperlink{h.37m2jsg}{~}}{(int
flag)}

{release a flag }

{int
}{\protect\hyperlink{h.1mrcu09}{thinkos\_flag\_val}}{\protect\hyperlink{h.1mrcu09}{~}}{(int
flag)}

{return the flag value }

{int
}{\protect\hyperlink{h.46r0co2}{thinkos\_flag\_set}}{\protect\hyperlink{h.46r0co2}{~}}{(int
flag)}

{set a flag }

{int
}{\protect\hyperlink{h.2lwamvv}{thinkos\_flag\_clr}}{\protect\hyperlink{h.2lwamvv}{~}}{(int
flag)}

{clear a flag }

{int
}{\protect\hyperlink{h.111kx3o}{thinkos\_flag\_give}}{\protect\hyperlink{h.111kx3o}{~}}{(int
flag)}

{signal a flag }

{void
}{\protect\hyperlink{h.3l18frh}{thinkos\_flag\_give\_i}}{\protect\hyperlink{h.3l18frh}{~}}{(int
flag)}

{signal a flag form inside an interrupt handler }

{int
}{\protect\hyperlink{h.206ipza}{thinkos\_flag\_take}}{\protect\hyperlink{h.206ipza}{~}}{(int
flag)}

{wait for flag to be signaled }

{int
}{\protect\hyperlink{h.4k668n3}{thinkos\_flag\_timedtake}}{\protect\hyperlink{h.4k668n3}{~}}{(int
flag, unsigned int ms)}

{wait for flag or timeout }

{int
}{\protect\hyperlink{h.1egqt2p}{thinkos\_gate\_alloc}}{\protect\hyperlink{h.1egqt2p}{~}}{(void)}

{Alloc a gate synchronization object. }

{int
}{\protect\hyperlink{h.3ygebqi}{thinkos\_gate\_free}}{\protect\hyperlink{h.3ygebqi}{~}}{(int
gate)}

{Frees the gate synchronization object. }

{int
}{\protect\hyperlink{h.2dlolyb}{thinkos\_gate\_wait}}{\protect\hyperlink{h.2dlolyb}{~}}{(int
gate)}

{Wait for a gate to open. }

{int
}{\protect\hyperlink{h.sqyw64}{thinkos\_gate\_timedwait}}{\protect\hyperlink{h.sqyw64}{~}}{(int
gate, unsigned int ms)}

{Wait for a gate to open or a timeout. }

{int
}{\protect\hyperlink{h.3cqmetx}{thinkos\_gate\_open}}{\protect\hyperlink{h.3cqmetx}{~}}{(int
gate)}

{Open or signal the gate. }

{void
}{\protect\hyperlink{h.1rvwp1q}{thinkos\_gate\_open\_i}}{\protect\hyperlink{h.1rvwp1q}{~}}{(int
gate)}

{Open or signal the gate from inside an interrupt handler. }

{int
}{\protect\hyperlink{h.4bvk7pj}{thinkos\_gate\_close}}{\protect\hyperlink{h.4bvk7pj}{~}}{(int
gate)}

{Close the gate if the gate is }{OPEN}{~ or remove pending signaling it
the gate is @ LOCKED. }

{int
}{\protect\hyperlink{h.2r0uhxc}{thinkos\_gate\_exit}}{\protect\hyperlink{h.2r0uhxc}{~}}{(int
gate, unsigned int open)}

{Exit the gate, leaving the gate, optionally leaving it open or closed.
}

{int
}{\protect\hyperlink{h.3q5sasy}{thinkos\_irq\_wait}}{\protect\hyperlink{h.3q5sasy}{~}}{(int
irq)}

{wait for interrupt }

{int
}{\protect\hyperlink{h.25b2l0r}{thinkos\_irq\_register}}{\protect\hyperlink{h.25b2l0r}{~}}{(int
irq, int pri, void(*isr)(void))}

{reqister an interrupt handler }

{int
}{\protect\hyperlink{h.34g0dwd}{thinkos\_clocks}}{\protect\hyperlink{h.34g0dwd}{~}}{(uint32\_t
*clk{[}{]})}

{get hardware clocks frequencies }

{int
}{\protect\hyperlink{h.1jlao46}{thinkos\_udelay\_factor}}{\protect\hyperlink{h.1jlao46}{~}}{(int32\_t
*factor)}

{get udelay calibration factor }

{void
}{\protect\hyperlink{h.43ky6rz}{\_\_attribute\_\_}}{\protect\hyperlink{h.43ky6rz}{~}}{((noreturn))
thinkos\_abort(void)}

{abort the operating system. }

{}

\subparagraph{\texorpdfstring{{Author:}}{Author:}}\label{author}

{Robinson Mittmann
}{\href{mailto:bobmittmann@gmail.com}{\nolinkurl{bobmittmann@gmail.com}}}{~}

{}

\subsubsection{\texorpdfstring{{Macro Definition
Documentation}}{Macro Definition Documentation}}\label{macro-definition-documentation}

\paragraph{\texorpdfstring{{\#define IRQ\_PRIORITY\_HIGHEST~ (1
\textless{}\textless{}
5)}}{\#define IRQ\_PRIORITY\_HIGHEST~ (1 \textless{}\textless{} 5)}}\label{define-irq_priority_highest-1-5}

\hypertarget{h.39kk8xu}{\paragraph{\texorpdfstring{{\#define
IRQ\_PRIORITY\_VERY\_HIGH~ (2 \textless{}\textless{}
5)}}{\#define IRQ\_PRIORITY\_VERY\_HIGH~ (2 \textless{}\textless{} 5)}}\label{h.39kk8xu}}

\hypertarget{h.1opuj5n}{\paragraph{\texorpdfstring{{\#define
IRQ\_PRIORITY\_HIGH~ (3 \textless{}\textless{}
5)}}{\#define IRQ\_PRIORITY\_HIGH~ (3 \textless{}\textless{} 5)}}\label{h.1opuj5n}}

\hypertarget{h.48pi1tg}{\paragraph{\texorpdfstring{{\#define
IRQ\_PRIORITY\_REGULAR~ (4 \textless{}\textless{}
5)}}{\#define IRQ\_PRIORITY\_REGULAR~ (4 \textless{}\textless{} 5)}}\label{h.48pi1tg}}

\hypertarget{h.2nusc19}{\paragraph{\texorpdfstring{{\#define
IRQ\_PRIORITY\_LOW~ (5 \textless{}\textless{}
5)}}{\#define IRQ\_PRIORITY\_LOW~ (5 \textless{}\textless{} 5)}}\label{h.2nusc19}}

\hypertarget{h.1302m92}{\paragraph{\texorpdfstring{{\#define
IRQ\_PRIORITY\_VERY\_LOW~ (6 \textless{}\textless{}
5)}}{\#define IRQ\_PRIORITY\_VERY\_LOW~ (6 \textless{}\textless{} 5)}}\label{h.1302m92}}

\hypertarget{h.3mzq4wv}{\paragraph{\texorpdfstring{{\#define
THINKOS\_OPT\_PRIORITY( VAL)~ (((VAL) \& 0xff) \textless{}\textless{}
16)}}{\#define THINKOS\_OPT\_PRIORITY( VAL)~ (((VAL) \& 0xff) \textless{}\textless{} 16)}}\label{h.3mzq4wv}}

\hypertarget{h.2250f4o}{\paragraph{\texorpdfstring{{\#define
THINKOS\_OPT\_ID( VAL)~ (((VAL) \& 0x07f) \textless{}\textless{}
24)}}{\#define THINKOS\_OPT\_ID( VAL)~ (((VAL) \& 0x07f) \textless{}\textless{} 24)}}\label{h.2250f4o}}

\hypertarget{h.haapch}{\paragraph{\texorpdfstring{{\#define
THINKOS\_OPT\_PAUSED~ (1 \textless{}\textless{} 31) /* don't run at
startup
*/}}{\#define THINKOS\_OPT\_PAUSED~ (1 \textless{}\textless{} 31) /* don't run at startup */}}\label{h.haapch}}

\hypertarget{h.319y80a}{\paragraph{\texorpdfstring{{\#define
THINKOS\_OPT\_STACK\_SIZE( VAL)~ ((VAL) \&
0xffff)}}{\#define THINKOS\_OPT\_STACK\_SIZE( VAL)~ ((VAL) \& 0xffff)}}\label{h.319y80a}}

{}

\subsubsection{\texorpdfstring{{Enumeration Type
Documentation}}{Enumeration Type Documentation}}\label{enumeration-type-documentation}

\paragraph{\texorpdfstring{{enum
}{\protect\hyperlink{h.40ew0vw}{thinkos\_err}}}{enum thinkos\_err}}\label{enum-thinkos_err}

{\protect\hyperlink{AAAAAAAAAM}{}}

{Enumerator}

{THINKOS\_OK}{~ No error }

{THINKOS\_ETIMEDOUT}{~ System call timed out }

{THINKOS\_EINTR}{~ System call interrupted out }

{THINKOS\_EINVAL}{~ Invalid argument }

{THINKOS\_EAGAIN}{~ Non blocking call failed }

{THINKOS\_EDEADLK}{~ Deadlock condition detected }

{THINKOS\_EPERM}{~ }

{THINKOS\_ENOSYS}{~ Invalid system call }

{THINKOS\_EFAULT}{~ }

{THINKOS\_ENOMEM}{~ Resource pool exausted }

\paragraph{\texorpdfstring{{enum
}{\protect\hyperlink{h.36ei31r}{thinkos\_obj\_kind}}}{enum thinkos\_obj\_kind}}\label{enum-thinkos_obj_kind}

{\protect\hyperlink{AAAAAAAAAX}{}}

{Enumerator}

{THINKOS\_OBJ\_READY}{~ }

{THINKOS\_OBJ\_TMSHARE}{~ }

{THINKOS\_OBJ\_CLOCK}{~ }

{THINKOS\_OBJ\_MUTEX}{~ }

{THINKOS\_OBJ\_COND}{~ }

{THINKOS\_OBJ\_SEMAPHORE}{~ }

{THINKOS\_OBJ\_EVENT}{~ }

{THINKOS\_OBJ\_FLAG}{~ }

{THINKOS\_OBJ\_GATE}{~ }

{THINKOS\_OBJ\_JOIN}{~ }

{THINKOS\_OBJ\_CONREAD}{~ }

{THINKOS\_OBJ\_CONWRITE}{~ }

{THINKOS\_OBJ\_PAUSED}{~ }

{THINKOS\_OBJ\_CANCELED}{~ }

{THINKOS\_OBJ\_FAULT}{~ }

{THINKOS\_OBJ\_INVALID}{~ }

{}

\subsubsection{\texorpdfstring{{Function
Documentation}}{Function Documentation}}\label{function-documentation-10}

\paragraph{\texorpdfstring{{int thinkos\_init (unsigned int
~opt)}}{int thinkos\_init (unsigned int ~opt)}}\label{int-thinkos_init-unsigned-int-opt}

{}

{Initializes the }{ThinkOS}{~ library. }

{On return the current program execution thread turns into the first
thread of the system. }

\subparagraph{\texorpdfstring{{Returns:}}{Returns:}}\label{returns-55}

{THINKOS\_OK µ}

{}



\section{\texorpdfstring{{Module
Index}}{Module Index}}\label{module-index}

\subsection{\texorpdfstring{{Modules}}{Modules}}\label{modules}

{Here is a list of all modules}

{Threads~~~~~~~~}

{Time related calls~~~~~~~~}

{Mutexes~~~~~~~~}

{Conditional Variables~~~~~~~~}

{Semaphores~~~~~~~~}

{Event sets~~~~~~~~}

{Flags~~~~~~~~}

{Gates~~~~~~~~}

{Interrupt Requests~~~~~~~~}

{OS Monitor and Control.~~~~~~~~}

{}

\section{\texorpdfstring{{Data Structure
Index}}{Data Structure Index}}\label{data-structure-index}

\subsection{\texorpdfstring{{Data
Structures}}{Data Structures}}\label{data-structures}

{Here are the data structures with brief descriptions:}

{\protect\hyperlink{h.2iq8gzs}{thinkos\_thread\_inf}}{\protect\hyperlink{h.2iq8gzs}{~}}{~~~~~~~~~}

\section{\texorpdfstring{{Modules
Summary}}{Modules Summary}}\label{modules-summary}



\section {Tools}
% = SECTION ===================================================================
\section{Tools}

% == SUB-SECTION =========================================
\subsection{MSYS-2}

MSYS2 is software distribution and a building platform for Windows. It provides a Unix-like environment, a command-line interface and a software repository making it easier to install, use, build and port software on Windows.\\

The official website is:

\href{http://www.msys2.org/}{MSYS2 Website}\\

Useful information can be found at the project's wiki pages at: 

\href{https://github.com/msys2/msys2/wiki/MSYS2-introduction}{MSYS2 Introduction}\\

% == SUB-SECTION =========================================
\subsection{GNU Toolchain for ARM Processors}

ThinnkOS on Windows Host - this document is a guide on how to set-up the tools and environment to develop ThinkOS applications using a Windows computer as host............

The GNU toolchain is a broad collection of programming tools developed by the GNU Project. These tools form a toolchain (a suite of tools used in a serial manner) used for developing software applications and operating systems.

If you are a windows user, you can download and install the GNU Toolchain for ARM Processors from the site: https://launchpad.net/gcc-arm-embedded. At the time of writing this document the latest version known to work with YARD-ICE was 4.9-2015-q3, which can be downloaded directly from:

\begin{flushleft}
GNU Arm Embedded Toolchain at Launchpad:
[\href{https://launchpad.net/gcc-arm-embedded}{https://launchpad.net/gcc-arm-embedded}]\\
\end{flushleft}

\begin{flushleft}
[\href{https://developer.arm.com/open-source/gnu-toolchain/gnu-rm}{GNU Arm Embedded Toolchain}]\\

Downloads: [\href{https://developer.arm.com/open-source/gnu-toolchain/gnu-rm/downloads}{https://developer.arm.com/open-source/gnu-toolchain/gnu-rm/downloads}\\
\end{flushleft}

Run the downloaded file \textit{gcc-arm-none-eabi-6-2017-q2-update-win32-sha2.exe}, or newer, to start the installation.
In the next steps use the default values suggested by the wizard installer.
In the last screen, after the files were copied, select the "Add path to environment variable" option.

Alternatively a compressed (.zip) file can be downloaded (\textit{gcc-arm-none-eabi-6-2017-q2-update-win32.zip}) and extracted in a directory of choice...

\subsubsection{GNU Make}

\subsubsection{GNU Core Utils for Windows}

% == SUB-SECTION =========================================
\subsection{Python}

% == SUB-SECTION =========================================
\subsection{Julia}


% == SUB-SECTION =========================================
\subsection{Eclipse}

% == SUB-SECTION =========================================
\subsection{Tera-Term}


% == SUB-SECTION =========================================
\subsubsection{DCCLog}

\section{Apendix}

\section{Efficient Developer}

\subsection{Master your tools}

\subsection{Automated Code Generation}

Embedded systems ...
In some circumstances there are some pieces of code can be ...
better performance comprises of a repetition of code that is tedious to type and error prone. Also in some contexts a generic solution that can be computed at runtime can be too expensive in terms of CPU usage to be computed in real-time. Or may require too much memory. A common solution for these problems is to use lookup tables and to interpolate intermidiate values, store coefficients in FLASH.
Even if you could afford to calculate a complex dataset some systems require a very fast boot-up response, so storing pre computed values is a sure way of impproving things.
We present here two typical scenarios where an automated generated code is benneficial for embedded systems.

\subsubsection{Fixed-point sine function}

Example sine function using Python.


\subsubsection{IIR filter}

In this example we use a Julia script to calculate and automatically generate a very efficient C code for a low pass filter. By automating the code generation its easy to change the filter parameters, like cur off frequency, filter order without having to worry about synchronizing the design specs with your embedded code.

\subsubsection{Embedded HTML pages}

Here HTML pages, or any pages for that matter, are automatically compressed, converted into C arrays. Makefile rules take care of the dependency of your code and the web pages, so the code is regenerated and compiled whenever a page is modified.


\subsubsection{Autogenerated Code Comments}

Always add a comment in the code indicating that it is automatically generated and a summary of instructions on how to regenerate it.

Example:

\subsection{Makefiles}

\subsection{Text Editor}

\subsubsection{Multiple editor windows}

\subsection{Git - Revision Control}

\section {References}

\subsection {Gnu Tools for Window}

More information on the GNU Make and CoreUtils for Windows package can be found at:

\begin{itemize}
\tightlist
\item
  {\href{http://gnuwin32.sourceforge.net/packages/make.html}{http://gnuwin32.sourceforge.net/packages/make.htm}}
\item
  {\href{http://gnuwin32.sourceforge.net/packages/coreutils.htm}{http://gnuwin32.sourceforge.net/packages/coreutils.htm}}
\end{itemize}

\subsection{GNU ARM Eclipse Plug-Ins}

{Quick info on GNU ARM Eclipse Plug-Ins:}

\begin{itemize}
\tightlist
\item
  {Name: }{`\emph{GNU ARM Eclipse Plug-ins}'}
\item
  {Location:
  }{\href{http://gnuarmeclipse.sourceforge.net/updates}{http://gnuarmeclipse.sourceforge.net/updates}}
\item
  {Website:
  }{\href{http://gnuarmeclipse.livius.net/blog}{http://gnuarmeclipse.livius.net/blog}}
\end{itemize}


\subsection{GNU MCU Eclipse plug-ins}

These plug-ins provide Eclipse CDT (C/C++ Development Tooling) extensions for GNU embedded toolchains like the GNU Tools for ARM Embedded Processors used as reference in this document.

\begin{itemize}
\tightlist
\item
	{Eclipse update site:
	}{\href{http://gnu-mcu-eclipse.netlify.com/v4-neon-updates}{http://gnu-mcu-eclipse.netlify.com/v4-neon-updates}}
\item
	{Website:
	}{\href{https://gnu-mcu-eclipse.github.io}{https://gnu-mcu-eclipse.github.io}}
\end{itemize}

\subsection{Git Extensions}

From the project's website: \quote{Git Extensions is a toolkit aimed at making working with Git under Windows more intuitive (note that Git Extensions is also available on Linux and Macintosh OS X using Mono). The shell extension will integrate in Windows Explorer and presents a context menu on files and directories.}\\


\paragraph{Website: }\href{http://git-extensions-documentation.readthedocs.io/en/latest/git_extensions.html}{Git Extensions}


% == SUB-SECTION =========================================
\subsection{YARD-ICE}

The YARD-ICE project is a community based, open source hardware and software platform for a standalone remote debugger. The embedded platform uses \ThinkOS as bootloader and real-time operating system.

\begin{itemize}
\tightlist
\item
  {YARD-ICE Project Location:
  }{\href{https://github.com/bobmittmann/yard-ice}{https://github.com/bobmittmann/yard-ice}}
\item
  {Latest Release:
  }{\href{https://github.com/bobmittmann/yard-ice/archive/0.24.zip}{https://github.com/bobmittmann/yard-ice/archive/0.24.zip}}
\item
  {Git Repository (HTTPS):}
  {\href{https://github.com/bobmittmann/yard-ice.git}{https://github.com/bobmittmann/yard-ice.git}}
\item
  {Git Repository (SSH):}{\href{git@github.com:bobmittmann/yard-ice.git}{git@github.com:bobmittmann/yard-ice.git}}
\end{itemize}

\subsection{Useful Web References}

\begin{flushleft}
\href{https://en.wikipedia.org/wiki/Real-time_operating_system}{Real-time operating system} - Wikipedia - March, 2017
\end{flushleft}

\begin{flushleft}
\href{https://www.gnu.org/gnu/thegnuproject.en.html}{The GNU Project}\\
\end{flushleft}

\begin{flushleft}
\href{https://gcc.gnu.org/}{GCC, the GNU Compiler Collection}\\
\end{flushleft}

\begin{flushleft}
\href{https://github.com/msys2/msys2/wiki}{MSYS2}\\
\end{flushleft}

\begin{flushleft}
\href{https://git-scm.com/downloads}{Git}\\
\end{flushleft}

\begin{flushleft}
\href{https://www.kernel.org/doc/html/v4.10/process/coding-style.html}{Linux kernel coding style}\\
\end{flushleft}


\end{document}

